\documentclass[11pt,reqno]{amsart}

\usepackage{amsmath,amssymb,amsthm,mathtools}
\usepackage[T1]{fontenc}
\usepackage{lmodern}
\usepackage{microtype}
\usepackage[margin=1in]{geometry}
\usepackage{booktabs}
\usepackage{array}
\usepackage{graphicx}
\usepackage{xcolor}
\usepackage[colorlinks=true,linkcolor=blue!60!black,citecolor=blue!60!black,urlcolor=blue!60!black,bookmarks=false]{hyperref}

\newtheorem{theorem}{Theorem}[section]
\newtheorem{proposition}[theorem]{Proposition}
\newtheorem{lemma}[theorem]{Lemma}
\newtheorem{corollary}[theorem]{Corollary}
\theoremstyle{remark}
\newtheorem{remark}[theorem]{Remark}

\numberwithin{equation}{section}
\emergencystretch=1.5em

\newcommand{\N}{\mathbb{N}}
\newcommand{\Z}{\mathbb{Z}}
\newcommand{\Q}{\mathbb{Q}}
\newcommand{\R}{\mathbb{R}}
\newcommand{\C}{\mathbb{C}}
\newcommand{\F}{\mathbb{F}}
\newcommand{\Pset}{\mathcal{P}}
\newcommand{\Podd}{\mathcal{P}_{\mathrm{odd}}}
\newcommand{\Peven}{\mathcal{P}_{\mathrm{even}}}
\newcommand{\Sodd}{S_{\mathrm{odd}}}
\newcommand{\Seven}{S_{\mathrm{even}}}
\newcommand{\Dparity}{\Delta_{\mathrm{parity}}}
\newcommand{\Sparitys}{\Sigma_{\mathrm{parity}}^*}

\DeclareMathOperator{\core}{core}
\DeclareMathOperator{\rad}{rad}
\DeclareMathOperator{\radodd}{rad_{\mathrm{odd}}}
\DeclareMathOperator{\ord}{ord}
\DeclareMathOperator{\Gal}{Gal}
\DeclareMathOperator{\Cl}{Cl}
\DeclareMathOperator{\Pic}{Pic}
\DeclareMathOperator{\GL}{GL}
\DeclareMathOperator{\SL}{SL}
\DeclareMathOperator{\Imag}{Im}
\DeclareMathOperator{\rank}{rank}
\DeclareMathOperator{\li}{li}

\title[Arithmetic Structure and Counting Bounds for the Brocard--Ramanujan Equation]{Arithmetic Structure and Counting Bounds for the\\ Brocard--Ramanujan Equation}

\author{Prince Anand}
\address{}
\email{princeanand123pc@gmail.com}
\date{November 2026}
\subjclass[2020]{Primary 11D61, 11D85; Secondary 11R11, 11R23, 11J13, 11L07, 11G05.}
\keywords{Brocard--Ramanujan Diophantine equation, factorial valuation parity, Vaaler trigonometric polynomials, Pell equations, Chebyshev--Lucas polynomials, R\'edei reciprocity, $p$-rationality, Iwasawa invariants, Greenberg's conjecture, Pad\'e approximations, support-reduced Wronskian.}

\begin{document}

\begin{abstract}
Let $n! + 1 = m^2$ be the Brocard--Ramanujan Diophantine equation, and write $n! = q_n s_n^2$ with $q_n$ square-free and $m + s_n\sqrt{q_n} = \varepsilon_n^k$ in $\Z[\sqrt{q_n}]$. We prove that $k \in \{1, 2\}$ for all sufficiently large solutions (and $n \in \{4, 5, 7\}$), that $K_n = \Q(\sqrt{q_n})$ is non-$p$-rational ($\lambda_p^*(K_n) \ge 1$) for every prime $p \in [5, n/3]$ with $p \nmid k$, and that $X(K_{n,\infty}) \ne 0$ at every split prime $p \in \Peven(n) \cap [5, n/4]$ with $p \nmid k$; consequently, either Gras's $p$-rationality at a single prime $p \in (n/5, n/3]$, or Greenberg's conjecture ($\lambda_p(K_n) = 0$) with $n_0^{(r)}(K_n) > r + 1$ ($\forall\, r \ge 0$) at a single such split prime, implies that $(4, 5)$, $(5, 11)$, and $(7, 71)$ are the only solutions to $n! + 1 = m^2$. Unconditionally, we prove that the number of solutions $n \le N$ with $k = 2$ is $\ll_\varepsilon N^\varepsilon$, and that for any polynomial $P(x) \in \Z[x]$ of degree $r \ge 2$, the number of solutions $n \le N$ to $n! = P(x)$ is $\ll_{P,\varepsilon} N^{6/7+\varepsilon}$ (improving to $N^{2/3+\varepsilon}$ for larger degree gaps).
\end{abstract}

\maketitle

\section{Introduction}

\subsection{Background and State of the Problem}
In 1876 and again in 1885, Brocard \cite{Brocard1876,Brocard1885} asked for all positive integer solutions $(n, m)$ to the Diophantine equation
\begin{equation}\label{eq:1.1}
n! + 1 = m^2.
\end{equation}
In 1913, Ramanujan \cite{Ramanujan1913} independently posed the same problem and noted the three solutions $(n, m) \in \{(4, 5), (5, 11), (7, 71)\}$. While quadratic-residue sieves have verified that no further solutions exist for $n \le 10^9$ (Berndt and Galway \cite{BerndtGalway2000}), $n < 4 \times 10^{11}$ (Matson \cite{Matson2017}), and $n \le 10^{15}$ (Epstein and Glickman \cite{EpsteinGlickman2020}), determining whether \eqref{eq:1.1} has only finitely many solutions remains a longstanding open problem in Diophantine number theory.

For higher-degree and shifted variants, several finiteness theorems are known. Erd\H{o}s and Obl\'ath \cite{ErdosOblath1937} proved in 1937 that $n! = x^p \pm y^p$ ($\gcd(x, y) = 1$) and $n! \pm m! = x^p$ have only finitely many integer solutions for prime exponents $p \ge 3$, and Pollack and Shapiro \cite{PollackShapiro1973} showed that $n! + 1 = x^4$ has no solutions. Conditionally, Overholt \cite{Overholt1993} observed in 1993 that the weak $abc$ conjecture implies the finiteness of solutions to \eqref{eq:1.1}. D\k{a}browski \cite{Dabrowski1996} studied $n! + A = m^2$, proving via the Chebotarev density theorem that it has only finitely many solutions when $A$ is not a square and is conditionally finite under the $abc$ conjecture when $A = u^2$. Luca \cite{Luca2002} extended Overholt's conditional theorem to general polynomial-factorial equations $P(x) = n!$ of degree $r \ge 2$ under the $abc$ conjecture, and further multiplicative and additive variants were studied by Kihel and Luca \cite{KihelLuca2008}, D\k{a}browski and Ulas \cite{DabrowskiUlas2013}, and Krasniqi \cite{Krasniqi2026}. Unconditionally, Berend and Harmse \cite{BerendHarmse2006} proved that $P(x) = n!$ has only finitely many solutions whenever $P(x) \in \Z[x]$ is irreducible of degree $\ge 2$ or violates $p$-adic valuation-balance constraints. However, when $P(x)$ splits over $\Q$ into linear factors of equal multiplicity---as in $x^2 - 1 = (x - 1)(x + 1)$---local $p$-adic obstructions alone do not rule out infinitely many solutions.

A complementary line of inquiry seeks unconditional upper bounds on the counting function
\begin{equation}\label{eq:1.2}
\mathcal{N}_{P,s}(N) := \#\bigl\{1 \le n \le N : \exists\, x \in \Z \text{ such that } s \cdot n! = P(x)\bigr\}
\end{equation}
for a polynomial $P(x) \in \Z[x]$ of degree $r \ge 2$ and $s \in \Z \setminus \{0\}$, and in particular on $\mathcal{N}_{\mathrm{Brocard}}(N) := \mathcal{N}_{x^2-1,1}(N)$. Answering a question of Erd\H{o}s, Berend and Osgood \cite{BerendOsgood1992} proved in 1992 via Osgood's effective bounds on linear forms in Siegel $G$-functions that $\mathcal{N}_{P,s}(N) = o(N)$ as $N \to \infty$. Recently, Bui, Pratt, and Zaharescu \cite{BuiPrattZaharescu2023} introduced a three-variable Pad\'e approximation method for the algebraic functions $\omega_i(x) = \prod_{j=1}^{\beta_i-1} (1 - jx)^{-1/r}$, establishing the first unconditional power-saving bound
\begin{equation}\label{eq:1.3}
\mathcal{N}_{P,s}(N) \ll_{P,s} N^{33/34}.
\end{equation}

\subsection{Main Results}
To state our principal theorems, let $n \ge 4$ be an integer and write the canonical square-free / square factorization of $n!$ as
\begin{equation}\label{eq:1.4}
n! = q_n s_n^2, \qquad q_n := \core(n!) = \prod_{p \in \Podd(n)} p, \qquad s_n := \prod_{p \le n} p^{\lfloor v_p(n!)/2 \rfloor},
\end{equation}
where $\Podd(n) := \{p \le n : 2 \nmid v_p(n!)\}$ and $\Peven(n) := \{p \le n : 2 \mid v_p(n!)\}$ denote the sets of primes $p \le n$ dividing $n!$ to odd and even powers, respectively. For every integer $n \ge 4$ with $n \ne 6$ (note that $6! + 1 = 721 = 7 \times 103$ is not a square), $q_n \ge 6$ is square-free and divisible by a prime $p \equiv 3 \pmod 4$, so the negative Pell equation $x^2 - q_n y^2 = -1$ is insoluble, and the group of positive units of the quadratic order $\Z[\sqrt{q_n}] \subset K_n := \Q(\sqrt{q_n})$ is generated by a unique fundamental unit $\varepsilon_n = X_1(q_n) + Y_1(q_n)\sqrt{q_n} > 1$ of norm $X_1(q_n)^2 - q_n Y_1(q_n)^2 = +1$. Every positive integer solution $(n, m)$ to \eqref{eq:1.1} satisfies $m^2 - q_n s_n^2 = 1$ and therefore corresponds to a unique positive integer Pell index $k \ge 1$ such that
\begin{equation}\label{eq:1.5}
m + s_n\sqrt{q_n} = \varepsilon_n^k = \bigl(X_1(q_n) + Y_1(q_n)\sqrt{q_n}\bigr)^k.
\end{equation}
A priori, lower bounds on the fundamental regulator $\log \varepsilon_n > 0.2\,n$ imply $1 \le k < 2.5\log n$ for all $n \ge 8$. Our first main theorem proves that $k$ collapses unconditionally to $\{1, 2\}$ for all sufficiently large $n$ (and for $n \in \{4, 5, 7\}$).

\begin{theorem}[{Unconditional Collapse of the Pell Index}]\label{thm:1.1}
Let $(n, m)$ be a positive integer solution to $n! + 1 = m^2$ with $n \ge 4$, and let $k \ge 1$ be the Pell index defined by \eqref{eq:1.5}. For the known solutions we have $k(4) = 1$, $k(5) = 1$, and $k(7) = 2$, and for all sufficiently large solutions $(n, m)$,
\begin{equation}\label{eq:1.6}
k \in \{1, 2\}.
\end{equation}
Moreover, as $n \to \infty$, the fundamental regulator $\log \varepsilon_n$ satisfies the quantization law
\begin{equation}\label{eq:1.7}
\frac{\log \varepsilon_n}{n\log n} = \frac{1}{2k} + O\!\left(\frac{1}{\log n}\right),
\end{equation}
and no two solutions $4 \le n_1 < n_2$ share the same square-free factorial core ($q_{n_1} \ne q_{n_2}$).
\end{theorem}

This Pell embedding has a direct consequence for the $p$-adic arithmetic, abelian $p$-ramification, and Iwasawa theory of the totally real quadratic field $K_n = \Q(\sqrt{q_n})$. For every $n \ge 116$, every prime $p \in (n/5, n/3]$ (a non-empty set containing $(n/5, n/4]$ by Nagura's theorem \cite{Nagura1952}) satisfies $p > n/5 > 2.5\log n > k \ge 1$, so $p \nmid k$ unconditionally; for sufficiently large $n$, Theorem~\ref{thm:1.1} ($k \in \{1, 2\}$) ensures $p \nmid k$ for all primes $p \ge 5$. Taking the $p$-adic logarithm of $\varepsilon_n^{2k} = (1 + 2q_n s_n^2) + 2ms_n\sqrt{q_n}$ transfers the $p$-adic valuation $v_p(s_n) = \lfloor v_p(n!)/2 \rfloor$ of $s_n$ directly to the normalized $p$-adic regulator $\widetilde{R}_p(K_n)$, the Kubota--Leopoldt $p$-adic $L$-function $L_p(1, \chi_{D_n})$, and the abelian $p$-ramification torsion group $\mathcal{T}_{K_n, p} := \Gal(M_p(K_n)/K_{n,\infty})$ \cite{Coates1977,Gras2016,Taya1996}. Combining the unramified Euler factor ($\chi_{D_n}(p) \in \{\pm 1\}$ for $p \in \Peven(n)$) with the ramified Euler factor ($\chi_{D_n}(p) = 0$ for $p \in \Podd(n)$) yields a unified parity-independent valuation law across all primes $p \in [5, n/2]$ with $p \nmid k$, establishing non-$p$-rationality and positive $p$-ramified Iwasawa rank $\lambda_p^*(K_n) := \rank_{\Z_p} X_p(K_{n,\infty}) \ge 1$ for every prime $p \in [5, n/3]$ with $p \nmid k$, as well as non-triviality of the unramified Iwasawa module $X(K_{n,\infty}) = \Gal(L(K_{n,\infty})/K_{n,\infty}) = \varprojlim A_m$ at every split prime $p \in \Peven(n) \cap [5, n/4]$ with $p \nmid k$ \cite{FukudaKomatsu1986,Jaulent2024,Ozaki1997,Taya1996,Yanagisawa2026}.

\begin{theorem}[{Non-$p$-Rationality, Iwasawa Invariants, and Conditional Finiteness}]\label{thm:1.2}
Let $(n, m)$ be a positive integer solution to $n! + 1 = m^2$ with $n \ge 116$, and let $K_n = \Q(\sqrt{q_n})$ be the associated real quadratic field of fundamental discriminant $D_n \in \{q_n, 4q_n\}$, class number $h(q_n)$, primitive quadratic character $\chi_{D_n}$, and cyclotomic $\Z_p$-extension $K_{n,\infty}/K_n$ with Galois group $\Gamma = \Gal(K_{n,\infty}/K_n)$.
\begin{enumerate}
\item[(i)] \textbf{Parity-Independent Non-$p$-Rationality and $p$-Ramified Iwasawa Rank:} For every prime $p \in [5, n/2]$ with $p \nmid k$, the Kubota--Leopoldt $p$-adic $L$-function and the abelian $p$-ramification torsion group $\mathcal{T}_{K_n, p} = \Gal(M_p(K_n)/K_{n,\infty}) \simeq X_p(K_{n,\infty})_\Gamma$ satisfy
\begin{equation}\label{eq:1.8}
v_p\bigl(L_p(1, \chi_{D_n})\bigr) = v_p\bigl(\#\mathcal{T}_{K_n, p}\bigr) = v_p\bigl(h(q_n)\bigr) + \left\lceil \frac{v_p(n!)}{2} \right\rceil - 1.
\end{equation}
In particular, for every prime $p \in (n/5, n/3]$ (a non-empty set for all $n \ge 116$), and more generally for every prime $p \in [5, n/3]$ with $p \nmid k$, the field $K_n$ is non-$p$-rational ($\mathcal{T}_{K_n, p} \ne 0$) and the $p$-ramified Iwasawa module $X_p(K_{n,\infty}) = \Gal(M_p(K_{n,\infty})/K_{n,\infty})$ is a free $\Z_p$-module of rank
\[
\lambda_p^*(K_n) := \rank_{\Z_p} X_p(K_{n,\infty}) = \lambda_p(\chi_{D_n}) \ge 1.
\]
\item[(ii)] \textbf{Unramified Iwasawa Module at Split Primes:} For every prime $p \in \Peven(n) \cap [5, n/4]$ with $p \nmid k$ that splits in $K_n$ ($\chi_{D_n}(p) = +1$), the unramified Iwasawa module $X(K_{n,\infty}) = \varprojlim A_m$ satisfies
\[
\bigl|X(K_{n,\infty})_\Gamma\bigr| = \bigl|A_m^\Gamma\bigr| = \bigl|\mathcal{T}_{K_n, p}\bigr| = p^{v_p(h(q_n)) + \frac{1}{2}v_p(n!) - 1} \ge p^{\frac{1}{2}v_p(n!) - 1} \ge p \qquad \forall\, m \ge \frac{1}{2}v_p(n!) - 1,
\]
so $\lambda_p(K_n) + \nu_p(K_n) \ge \frac{1}{2}v_p(n!) - 1 \ge 1$ unconditionally, and $\lambda_p(K_n) \ge 1$ whenever Taya's $p$-unit invariant satisfies $n_0^{(r)}(K_n) > r + 1$ for all $r \ge 0$ (in particular, if the subgroup $D_r \subseteq A_r$ generated by the primes above $p$ is trivial for all $r \ge 0$).
\item[(iii)] \textbf{Conditional Finiteness:} Consequently, if for each $n \ge 116$ either \textup{(a)} $K_n$ satisfies Gras's $p$-rationality property ($\mathcal{T}_{K_n, p} = 0$, cf.\ \cite{Gras2016}) at a single prime $p \in (n/5, n/3]$, \textup{(b)} the unramified Iwasawa module $X(K_{n,\infty})$ is trivial ($\lambda_p(K_n) = \nu_p(K_n) = 0$ \cite{FukudaKomatsu1986,Jaulent2024,Yanagisawa2026}) at a single split prime $p \in \Peven(n) \cap [5, n/4]$ with $p \nmid k$, or \textup{(c)} Greenberg's conjecture ($\lambda_p(K_n) = 0$ \cite{Greenberg1976}) holds together with $n_0^{(r)}(K_n) > r + 1$ ($\forall\, r \ge 0$ \cite{Taya1996}) at a single split prime $p \in \Peven(n) \cap [5, n/4]$ with $p \nmid k$, then $n! + 1 = m^2$ has no solutions with $n \ge 116$, and $(4, 5)$, $(5, 11)$, and $(7, 71)$ are the only solutions to \eqref{eq:1.1}.
\end{enumerate}
\end{theorem}

Our third main theorem establishes two unconditional counting bounds: an $O_\varepsilon(N^\varepsilon)$ sparsity bound for solutions to \eqref{eq:1.1} in the critical regime $k = 2$, and an $O_{P,s,\varepsilon}(N^{6/7+\varepsilon})$ upper bound for the general polynomial-factorial counting function $\mathcal{N}_{P,s}(N)$ (improving \eqref{eq:1.3}, with still stronger exponents up to $O_{P,s,\varepsilon}(N^{2/3+\varepsilon})$ when $P(x)$ has degree gap $\nu \ge 3$).

\begin{theorem}[{Unconditional $N^\varepsilon$ and $N^{6/7+\varepsilon}$ Counting Bounds}]\label{thm:1.3}
For every $\varepsilon > 0$ and all $N \ge 1$, the following unconditional bounds hold:
\begin{enumerate}
\item[(i)] \textbf{Extreme Sparsity of the $k = 2$ Regime:} Any two solutions $4 \le n_1 < n_2$ to $n! + 1 = m^2$ in the $k = 2$ regime satisfy $n_2 - n_1 > n_1^{1-\varepsilon}$ for $n_1$ sufficiently large in terms of $\varepsilon$, and
\begin{equation}\label{eq:1.9}
\#\bigl\{1 \le n \le N : \exists\, m \in \N,\; n! + 1 = m^2 \text{ with } k = 2\bigr\} \ll_\varepsilon N^\varepsilon.
\end{equation}
\item[(ii)] \textbf{Global $N^{6/7+\varepsilon}$ and Monomial-Product Bounds for $n! = P(x)$:} For any polynomial $P(x) \in \Z[x]$ of degree $r \ge 2$ (and more generally for $s \cdot n! = P(x)$ with any $s \in \Z \setminus \{0\}$),
\begin{equation}\label{eq:1.10}
\mathcal{N}_{P,s}(N) \ll_{P,s,\varepsilon} N^{6/7+\varepsilon}, \qquad \text{and in particular} \qquad \mathcal{N}_{\mathrm{Brocard}}(N) \ll_\varepsilon N^{6/7+\varepsilon}.
\end{equation}
Furthermore, if $P(x) = a_r x^r + a_{r-\nu} x^{r-\nu} + \dots + a_0 \in \Z[x]$ ($a_r \ne 0$, $r \ge 3$) has degree gap $\nu := r - \deg(P - a_r x^r) \ge 3$ after depressing the $x^{r-1}$ term, then $\mathcal{N}_{P,s}(N) \ll_{P,s,\varepsilon} N^{5/6+\varepsilon}$ for $\nu = 3$, $N^{4/5+\varepsilon}$ for $\nu = 4$, $N^{3/4+\varepsilon}$ for $\nu \ge 6$, and $N^{2/3+\varepsilon}$ as $\nu \to \infty$.
\end{enumerate}
\end{theorem}

Beyond Theorems~\ref{thm:1.1}--\ref{thm:1.3}, we establish in the body of the paper several unconditional structural results governing the $k = 2$ and $k = 1$ regimes and the associated Frey--Hellegouarch elliptic curves:
\begin{itemize}
\item \textbf{Two-stage Euclidean quartic law in the $k = 2$ regime (Section~\ref{sec:4}, Theorem~\ref{thm:4.2}):} When $k = 2$, Euclidean division of the even-valuation product $A_n$ by the odd-valuation product $B_n$ obeys the exact quartic identity $A_n = B_n Z^4 + Z^2$ ($R_n^2 = Q_n$), a second-stage quotient law $\lfloor B_+ R_n / B_- \rfloor = Z_-^4$, and stretched-exponential midpoint core bounds $\delta_i(n) := \core(X_1(q_n) \mp 1) \ge \exp(n^{1/2-o(1)})$.
\item \textbf{Effective Landau--Siegel bounds and R\'edei--Reichardt 4-rank cohomology in the $k = 1$ regime (Section~\ref{sec:5}, Theorems~\ref{thm:5.2} and \ref{thm:5.3}):} The midpoint cores satisfy the effective unconditional lower bound $\delta_i^{\mathrm{mid}}(n) > \exp(c_0\sqrt{\log n})$, the norm equations $x^2 - q_n y^2 = \pm 2$ are insoluble, and the kernel of the ramified R\'edei symbol matrix $R_n^{\mathrm{ram}}$ over $\F_2$ satisfies $\dim_{\F_2} \ker(R_n^{\mathrm{ram}}) = 1 + \rank_4 \Cl^+(q_n)$.
\item \textbf{Level-2 Frey curve CM reduction and effective Tatuzawa class number explosion (Section~\ref{sec:8}, Theorems~\ref{thm:8.1} and \ref{thm:8.2}):} The 2-isogenous Frey curve $F_n^{(2)}$ satisfies $j(F_n^{(2)}) \equiv 1728 \pmod{X_1(q_n)^2}$, and the non-genus factor of the class group $\Cl(q_n)$ satisfies $h_{\mathrm{non\text{-}gen}}(q_n) \ge q_n^{1/2-\varepsilon-o(1)}$ with at most one exceptional index $n_*$.
\end{itemize}

\subsection{Principal Obstacles and Overview of the Method}
We now describe the main obstacles encountered in proving Theorems~\ref{thm:1.1}--\ref{thm:1.3} and the ideas used to overcome them.

\medskip
\noindent\textbf{1. Valuation-parity cancellation and the collapse of the Pell index.}
In \eqref{eq:1.5}, the Chebyshev polynomial factorization $s_n = Y_1(q_n)U_{k-1}(X_1(q_n))$ (where $U_{k-1}$ is the Chebyshev polynomial of the second kind) splits $n! = q_n s_n^2$ into the fundamental Pell factor $q_n Y_1(q_n)^2$ and the Chebyshev factor $U_{k-1}(X_1(q_n))^2$. By a $p$-adic Lifting-the-Exponent calculation (Lemma~\ref{lem:3.2}), every prime $p \in \Podd(n) \setminus \{3\}$ divides $q_n$ and satisfies $v_p(U_{k-1}(X_1(q_n))) = v_p(k)$, so $p$ is absorbed into $q_n Y_1(q_n)^2$ with valuation $v_p(n!) - 2v_p(k)$. Consequently,
\[
\log\bigl(q_n Y_1(q_n)^2\bigr) \ge \Sodd(n) - 2\log k, \qquad \text{where} \qquad \Sodd(n) := \sum_{\substack{3 \le p \le n \\ p \in \Podd(n)}} v_p(n!)\log p.
\]
On the other hand, $q_n Y_1(q_n)^2 < \varepsilon_n^2 = (m + s_n\sqrt{q_n})^{2/k} \asymp (n!)^{1/k}$, so $\log(q_n Y_1(q_n)^2) \le \frac{1}{k}n\log n + O(n)$. Thus, proving $k \in \{1, 2\}$ reduces to proving the asymptotic valuation-parity equidistribution $\Sodd(n) = (\frac{1}{2} + o(1))n\log n$, or equivalently $\Dparity(n) := \Seven(n) - \Sodd(n) = \sum_{3 \le p \le n} (-1)^{v_p(n!)} v_p(n!)\log p = o(n\log n)$.

The obstacle in estimating $\Dparity(n)$ is that Legendre's valuation formula $v_p(n!) = \sum_{r=1}^R \lfloor n/p^r \rfloor$ on a dyadic block $p \asymp P$ involves $R = \lfloor \frac{\log n}{\log P} \rfloor$ floor functions; expanding $(-1)^{v_p(n!)}$ via an $R$-fold product of Vaaler trigonometric polynomials of degree $H$ produces an $L^1$ coefficient norm $(\frac{2}{\pi}\log H)^R$ that blows up whenever $R \to \infty$, while on the transition region $p \approx n^{1/r}$ the highest active floor $\lfloor n/p^r \rfloor$ has derivative too small for van der Corput bounds. In Section~\ref{sec:2}, we overcome both difficulties simultaneously. First, on $3 \le p \le \sqrt{n}$, Legendre's base-$p$ digit formula $v_p(n!) = \frac{n - s_p(n)}{p - 1}$ replaces $v_p(n!)\log p$ by $n\frac{\log p}{p - 1}$ with an error of $O(\sqrt{n})$ without using exponential sums. Second, for any fixed $\delta \in (0, 1/3]$ with $R = \lfloor 1/\delta \rfloor$, on $(n^\delta, \sqrt{n}]$ we have $p^{R+1} > n$, so $v_p(n!) = \sum_{r=1}^R \lfloor n/p^r \rfloor$ involves at most $R$ floor functions. By excising the transition collars $\mathcal{C}_r(n) = (\frac{n^{1/r}}{(\log n)^{A_R/r}}, n^{1/r}]$ ($2 \le r \le R$) and partitioning each $\mathcal{C}_r(n)$ into $M_r(n) \le (\log n)^{A_R}$ subintervals $J_{r,m_r}$ on each of which $\lfloor n/p^r \rfloor = m_r$ is constant, the highest floor becomes constant on $J_{r,m_r}$ while all remaining active phases $n/p^s$ ($1 \le s \le r-1$) have derivative parameter $T_0 \ge n^{1/r} \ge n^\delta$, whereas on the complementary inter-collar intervals every active phase satisfies $T_0 \ge 2^{-R}(\log n)^{A_R}$. This establishes $\Sodd(n) = (\frac{1}{2} + o(1))n\log n$ and completes the proof of Theorem~\ref{thm:1.1} and Theorem~\ref{thm:1.2} in Section~\ref{sec:3}.

\medskip
\noindent\textbf{2. Quartic Euclidean rigidity, R\'edei cohomology, and the pairwise capacity diamond.}
With $k \in \{1, 2\}$ established, Sections~\ref{sec:4} and \ref{sec:5} analyze the fine arithmetic structure of the two regimes $k = 2$ and $k = 1$. When $k = 2$, $n! = 4(X_1(q_n) - 1)X_1(q_n)^2(X_1(q_n) + 1)$, and we prove that Euclidean division of the even-valuation product $A_n$ by the odd-valuation product $B_n$ obeys the two-stage quartic laws $A_n = B_n Z^4 + Z^2$ ($Q_n = Z^4$, $R_n = Z^2$) and $\lfloor B_+ R_n / B_- \rfloor = Z_-^4$, forcing the midpoint cores $\delta_i(n) = \core(X_1(q_n) \mp 1)$ to grow at least as fast as $\exp(n^{1/2-o(1)})$. When $k = 1$, writing $\frac{m-1}{2} = \delta_1(n)z_1^2$ and $\frac{m+1}{2} = \delta_2(n)z_2^2$ ($\delta_1(n)\delta_2(n) = q_n$), we prove effective Landau--Siegel sub-exponential lower bounds $\delta_i^{\mathrm{mid}}(n) > \exp(c_0\sqrt{\log n})$, an exact Galois cohomology sequence identifying the kernel of the ramified R\'edei matrix over $\F_2$ with $1 + \rank_4 \Cl^+(q_n)$, and a reciprocal cross-symmetry law equating the odd-valuation weight of $\frac{m+1}{2}$ with the even-valuation weight of $\frac{m-1}{2}$ up to $O(n)$.

In Section~\ref{sec:6}, we combine this cross-symmetry law with a universal $\SL_2(\Z)$ four-GCD decomposition $XYZW = n_1!/4$ for any pair of solutions $n_1 < n_2$. Because Vaaler's trigonometric polynomial for $(-1)^{\lfloor x \rfloor}$ is supported exclusively on odd harmonics $2m + 1 \ne 0$, the correlation sum $\sum_{p \le n_1} (-1)^{v_p(n_1!) + v_p(n_2!)} v_p(n_1!)\log p$ has non-vanishing diagonal phase $\frac{(2m_2+1)(n_2-n_1)}{2p} \ne 0$. This forces the odd-core polarization weights $(\theta(n_1), \theta(n_2))$ of any two solutions with $n_2 - n_1 = n_1^\alpha$ to lie inside the Capacity Diamond $\frac{1-\alpha}{2} - o(1) \le \theta(n_1) + \theta(n_2) \le \frac{3+\alpha}{2} + o(1)$. Since $\theta(n) = 0$ whenever $k = 2$, this immediately forces $\alpha \ge 1 - o(1)$, proving $n_2 - n_1 > n_1^{1-\varepsilon}$ and $\#\mathcal{B}_{k=2}(N) \ll_\varepsilon N^\varepsilon$ (Theorem~\ref{thm:1.3}(i)).

\medskip
\noindent\textbf{3. The support-reduced Wronskian and monomial-product Pad\'e approximations.}
To prove the global counting bound $\mathcal{N}_{P,s}(N) \ll_{P,s,\varepsilon} N^{6/7+\varepsilon}$ of Theorem~\ref{thm:1.3}(ii) in Section~\ref{sec:7}, we generalize the Pad\'e approximation method of Bui, Pratt, and Zaharescu \cite{BuiPrattZaharescu2023} from $m = 2$ shifts to $m + 1$ variables $\omega_0(x), \omega_1(x), \dots, \omega_m(x)$. In \cite{BuiPrattZaharescu2023}, the Pad\'e vanishing-slack parameter $\varepsilon_0$ in $O = \lfloor (m + 1 - \varepsilon_0)D \rfloor$ was restricted to $\varepsilon_0 \in (0, 1)$; for $\varepsilon_0 < 1$, the maximum saving over all $m \ge 2$ is capped at $\theta < 1/14$ ($O(N^{13/14})$). The reason for restricting $\varepsilon_0 < 1$ in \cite{BuiPrattZaharescu2023} was that when $\varepsilon_0 \ge 1$, Siegel's lemma may produce a Pad\'e tuple $(P_0, \dots, P_m)$ in which one or more polynomials $P_i(x)$ vanish identically, causing the full $(m + 1) \times (m + 1)$ Wronskian determinant to vanish identically. We resolve this by passing to the non-zero support $S := \{0 \le i \le m : P_i(x) \not\equiv 0\}$ ($s := |S|$) and forming the $s \times s$ \emph{support-reduced Wronskian} $\Delta_S(x) := \det(P_{i_a}^{[k]}(x))_{0 \le k < s,\, a \in S}$. Because the subfamily $\{\omega_i : i \in S\}$ is linearly independent over $\C(x)$ and every $P_i$ ($i \in S$) is non-zero, $\Delta_S(x) \not\equiv 0$ is guaranteed for all $\varepsilon_0 \in (0, m - 1)$ with $s \ge \lceil m + 1 - \varepsilon_0 \rceil \ge 2$, and because the solution vector $p_i = n_0^{\beta_i/r} x_i > 0$ has strictly positive coordinates at every index $0 \le i \le m$, its restriction to $S$ is never zero. Optimizing over $\varepsilon_0 \in (0, m - 1)$ at $m = 5$ and $\varepsilon_0 = 5 - \sqrt{5} \in (2, 3)$ yields $\theta_5(5 - \sqrt{5}) = \frac{7 - 3\sqrt{5}}{2} \approx 0.14590 > \frac{1}{7}$, reaching the $m = 5$ Pigeonhole cap $\theta_5^* = \frac{1}{7}$ and proving $\mathcal{N}_{P,s}(N) \ll_{P,s,\varepsilon} N^{6/7+\varepsilon}$. For polynomials with degree gap $\nu \ge 3$, we further replace the linear shift basis by the monomial-product basis $\Omega_{\mathbf{e}}(x) = \prod_{i=1}^d \omega_i(x)^{e_i}$ ($|\mathbf{e}| \le E \le \min(r - 1, \nu - 1)$), generating $\binom{d+E}{d}$ basis functions from only $d$ shifts and raising the Pigeonhole cap from $\frac{1}{m+2}$ to $\frac{1}{d+2}$.

Finally, Section~\ref{sec:8} establishes the properties of the associated Frey--Hellegouarch elliptic curves and the effective Tatuzawa class-group bounds.

\subsection{Notation}
Throughout the paper, $\N = \{1, 2, 3, \dots\}$ denotes the set of positive integers, $\Pset$ the set of primes, $\pi(x) = \sum_{p \le x} 1$, and $\vartheta(x) = \sum_{p \le x} \log p$. For a prime $p$ and a non-zero integer $a$, $v_p(a)$ denotes the exact $p$-adic valuation of $a$, and $p^e \parallel a$ means $v_p(a) = e$. For a positive integer $A = \prod_p p^{v_p(A)}$, we write $\core(A) := \prod_{2 \nmid v_p(A)} p$ for its square-free core, $\rad(A) := \prod_{p \mid A} p$ for its radical, $\omega(A) := \sum_{p \mid A} 1$ for its number of distinct prime factors, $P^+(A)$ for its largest prime factor, and $A^* := A / 2^{v_2(A)}$ for its odd part. We write $e(t) := e^{2\pi i t}$, and we use Vinogradov's notation $f \ll g$ interchangeably with Landau's notation $f = O(g)$, indicating dependence of implied constants on auxiliary parameters by subscripts.

\section{Valuation-Parity Equidistribution of Factorials}\label{sec:2}

For an integer $n \ge 4$, define the odd- and even-valuation prime weights of $\log(n!)$ over odd primes $3 \le p \le n$ together with their parity discrepancy:
\begin{equation}\label{eq:2.1}
\Sodd(n) := \sum_{\substack{3 \le p \le n \\ p \in \Podd(n)}} v_p(n!)\log p, \qquad \Seven(n) := \sum_{\substack{3 \le p \le n \\ p \in \Peven(n)}} v_p(n!)\log p,
\end{equation}
and
\begin{equation}\label{eq:2.2}
\Dparity(n) := \Seven(n) - \Sodd(n) = \sum_{3 \le p \le n} (-1)^{v_p(n!)} v_p(n!)\log p.
\end{equation}
Define the small-prime valuation-parity sum
\begin{equation}\label{eq:2.3}
\Sparitys(n) := \sum_{3 \le p \le \sqrt{n}} (-1)^{v_p(n!)-1} \frac{\log p}{p - 1}.
\end{equation}
In this section we establish the asymptotic equidistribution of $\Sodd(n)$ and $\Seven(n)$, the second-order formula for $\Dparity(n)$, and the multi-modulus power-band equidistribution law.

\subsection{Harmonic Prime Layers and the Square-Free Factorial Core}
We begin by partitioning the upper prime range $(\sqrt{n}, n]$ into harmonic layers on which $v_p(n!)$ is constant.

\begin{lemma}[{Harmonic Layer Stratification and Asymptotics of $\Podd(n)$ and $\log q_n$ \cite{Berend1997,Sander2001,Broughan2002}}]\label{lem:2.1}
Let $n \ge 4$ be an integer, and define the harmonic intervals $I_j := \bigl(\frac{n}{j+1}, \frac{n}{j}\bigr]$ for $j \ge 1$.
\begin{enumerate}
\item[(i)] For every integer $1 \le j \le \lfloor \sqrt{n} \rfloor - 1$ and every prime $p \in I_j = \bigl(\frac{n}{j+1}, \frac{n}{j}\bigr]$, we have $v_p(n!) = j$. In particular, on odd layers $I_{2r-1}$ ($1 \le 2r - 1 \le \lfloor \sqrt{n} \rfloor - 1$), every prime $p \in I_{2r-1}$ belongs to $\Podd(n)$ with $v_p(q_n) = 1$ and $v_p(s_n) = r - 1$; on even layers $I_{2r}$ ($2 \le 2r \le \lfloor \sqrt{n} \rfloor - 1$), every prime $p \in I_{2r}$ belongs to $\Peven(n)$ with $v_p(q_n) = 0$ and $v_p(s_n) = r$.
\item[(ii)] There exists an absolute effective constant $c_0 > 0$ such that as $n \to \infty$,
\begin{equation}\label{eq:2.4}
\log q_n = (\log 2)\,n + O\!\left(n\exp\bigl(-c_0\sqrt{\log n}\bigr)\right), \qquad |\Podd(n)| = (\log 2 + o(1))\,\pi(n).
\end{equation}
\end{enumerate}
\end{lemma}

\begin{proof}
(i) For $1 \le j \le \lfloor \sqrt{n} \rfloor - 1$, we have $j + 1 \le \lfloor \sqrt{n} \rfloor \le \sqrt{n}$, so every prime $p \in I_j = \bigl(\frac{n}{j+1}, \frac{n}{j}\bigr]$ strictly satisfies $p > \frac{n}{j+1} \ge \sqrt{n}$. Consequently, $p^2 > n$, so $\lfloor n/p^m \rfloor = 0$ for all $m \ge 2$. Since $\frac{n}{j+1} < p \le \frac{n}{j} \iff j \le \frac{n}{p} < j + 1$, Legendre's formula \cite{Legendre1798} gives
\[
v_p(n!) = \sum_{m=1}^\infty \left\lfloor \frac{n}{p^m} \right\rfloor = \left\lfloor \frac{n}{p} \right\rfloor = j.
\]
Writing $v_p(n!) = v_p(q_n) + 2v_p(s_n)$ with $v_p(q_n) \in \{0, 1\}$ yields the stated formulas on $I_{2r-1}$ and $I_{2r}$.

(ii) Split the prime range $p \le n$ at $\sqrt{n}$. The primes $p \le \sqrt{n}$ contribute at most $\sum_{p \le \sqrt{n}} \log p = \vartheta(\sqrt{n}) \ll \sqrt{n}$ to $\log q_n$ and at most $\pi(\sqrt{n}) \ll \frac{\sqrt{n}}{\log n}$ to $|\Podd(n)|$. On the upper range $(\sqrt{n}, n]$, set $J := \lfloor \exp(\frac{1}{2}c_1\sqrt{\log n}) \rfloor$, where $c_1 > 0$ is the constant in the de la Vall\'ee Poussin Prime Number Theorem \cite[Chapter~18]{Davenport2000}:
\begin{equation}\label{eq:2.5}
\vartheta(x) = x + O\!\left(x\exp\bigl(-c_1\sqrt{\log x}\bigr)\right), \qquad \pi(x) = \li(x) + O\!\left(x\exp\bigl(-c_1\sqrt{\log x}\bigr)\right).
\end{equation}
The primes in $(\sqrt{n}, n/J]$ contribute at most $\vartheta(n/J) \ll \frac{n}{J} = n\exp(-\frac{1}{2}c_1\sqrt{\log n})$ to $\log q_n$. On $(n/J, n]$, part (i) shows that a prime $p \in (n/J, n]$ belongs to $\Podd(n)$ if and only if $p \in I_{2r-1} = \bigl(\frac{n}{2r}, \frac{n}{2r-1}\bigr]$ for some $1 \le 2r - 1 \le J - 1$. Summing $\vartheta\bigl(\frac{n}{2r-1}\bigr) - \vartheta\bigl(\frac{n}{2r}\bigr)$ over $1 \le 2r - 1 \le J - 1$ via \eqref{eq:2.5} gives
\begin{align*}
\log q_n &= \sum_{1 \le 2r-1 \le J-1} \left(\vartheta\!\left(\frac{n}{2r - 1}\right) - \vartheta\!\left(\frac{n}{2r}\right)\right) + O\!\left(\frac{n}{J}\right) \\
&= n \sum_{1 \le 2r-1 \le J-1} \left(\frac{1}{2r - 1} - \frac{1}{2r}\right) + O\!\left(n(\log J)\exp\bigl(-c_1\sqrt{\log(n/J)}\bigr) + \frac{n}{J}\right).
\end{align*}
Since $\sum_{r=1}^\infty \bigl(\frac{1}{2r-1} - \frac{1}{2r}\bigr) = \sum_{j=1}^\infty \frac{(-1)^{j-1}}{j} = \log 2$ and $\sum_{2r > J} \bigl(\frac{1}{2r-1} - \frac{1}{2r}\bigr) = O(J^{-1})$, setting $c_0 := c_1/3$ yields the first relation in \eqref{eq:2.4}. Similarly, summing $\pi\bigl(\frac{n}{2r-1}\bigr) - \pi\bigl(\frac{n}{2r}\bigr) = \bigl(\frac{1}{2r-1} - \frac{1}{2r} + o(1)\bigr)\frac{n}{\log n}$ over $1 \le 2r - 1 \le \lfloor \log n \rfloor$ yields $|\Podd(n)| = (\log 2 + o(1))\pi(n)$.
\end{proof}

\subsection{Exact Base-$p$ Digit Reduction Below $\sqrt{n}$}
Whereas primes $p > \sqrt{n}$ satisfy $v_p(n!) = \lfloor n/p \rfloor$, primes $p \le \sqrt{n}$ have higher powers $p^2, p^3, \dots \le n$ contributing to $v_p(n!)$. The next lemma shows that Legendre's base-$p$ digit identity replaces the weighted sum $\sum_{3 \le p \le \sqrt{n}} (-1)^{v_p(n!)} v_p(n!)\log p$ by $-n\Sparitys(n)$ up to a deterministic error of $O(\sqrt{n})$, avoiding any Fourier expansion on $p \le \sqrt{n}$ when relating $\Dparity(n)$ to $\Sparitys(n)$.

\begin{lemma}[{Exact Legendre Base-$p$ Reduction on $3 \le p \le \sqrt{n}$}]\label{lem:2.2}
For every integer $n \ge 4$, write the base-$p$ expansion of $n$ as $n = \sum_{j=0}^{K_p} d_j^{(p)}(n)p^j$ ($0 \le d_j^{(p)}(n) \le p - 1$, $K_p = \lfloor \frac{\log n}{\log p} \rfloor$), and let $s_p(n) := \sum_{j=0}^{K_p} d_j^{(p)}(n)$. Across $3 \le p \le \sqrt{n}$, we have
\begin{equation}\label{eq:2.6}
\sum_{3 \le p \le \sqrt{n}} (-1)^{v_p(n!)} v_p(n!)\log p = -n\,\Sparitys(n) + E_{\mathrm{digit}}(n),
\end{equation}
where the digit-sum remainder $E_{\mathrm{digit}}(n) := \sum_{3 \le p \le \sqrt{n}} (-1)^{v_p(n!)-1} \frac{s_p(n)\log p}{p - 1}$ satisfies the explicit bound
\begin{equation}\label{eq:2.7}
\bigl|E_{\mathrm{digit}}(n)\bigr| \le 2\,\pi\bigl(\sqrt{n}\bigr)\log n = O\bigl(\sqrt{n}\bigr).
\end{equation}
\end{lemma}

\begin{proof}
Multiplying Legendre's identity $v_p(n!) = \frac{n - s_p(n)}{p - 1}$ by $(-1)^{v_p(n!)}\log p = -(-1)^{v_p(n!)-1}\log p$ and summing over all primes $3 \le p \le \sqrt{n}$ yields \eqref{eq:2.6}. For each prime $3 \le p \le \sqrt{n}$, since $0 \le d_j^{(p)}(n) \le p - 1$ for every $0 \le j \le K_p$, we have
\[
0 \le \frac{s_p(n)}{p - 1} \le K_p + 1 = \left\lfloor \frac{\log n}{\log p} \right\rfloor + 1 \le \frac{\log n}{\log p} + 1.
\]
Multiplying by $\log p \le \frac{1}{2}\log n$ gives $0 \le \frac{s_p(n)\log p}{p - 1} \le \log n + \log p \le \frac{3}{2}\log n < 2\log n$. Summing over the $\pi(\sqrt{n})$ primes in $[3, \sqrt{n}]$ and applying Chebyshev's bound $\pi(\sqrt{n}) \ll \frac{\sqrt{n}}{\log n}$ yields \eqref{eq:2.7}.
\end{proof}

\subsection{Collar-Decomposed Two-Variable Vaaler Sums and Equidistribution Theorems}
To evaluate the upper range $(\sqrt{n}, n]$ and cancel $\Sparitys(n)$ on $(n^\delta, \sqrt{n}]$, we approximate the square-wave parity function $(-1)^{\lfloor x \rfloor} = 2\psi(x) - 4\psi(x/2)$ ($x \notin \Z$, where $\psi(t) := t - \lfloor t \rfloor - \frac{1}{2}$) by Vaaler's extremal trigonometric polynomials \cite[Theorem~18]{Vaaler1985} with matched cutoff degrees $H$ and $2H + 1$.

\begin{lemma}[{Vaaler's Trigonometric Polynomial Approximation {\cite[Theorem~18]{Vaaler1985}}}]\label{lem:2.3}
For every integer $H \ge 1$, let
\[
\Psi_H(t) = \sum_{1 \le |h| \le H} c_h(H)\,e(ht), \qquad B_H(t) = \sum_{|h| \le H} b_h(H)\,e(ht),
\]
be Vaaler's trigonometric polynomials with $c_h(H) = -\frac{1}{2\pi i h} W\bigl(\frac{|h|}{H+1}\bigr)$ ($|c_h(H)| \le \frac{1}{2\pi |h|}$), $b_0(H) = \frac{1}{2H+2}$, and $|b_h(H)| \le \frac{1}{2H+2}$ satisfying $|\psi(t) - \Psi_H(t)| \le B_H(t)$ for all $t \in \R$. Define
\begin{equation}\label{eq:2.8}
\widetilde{\Psi}_H(x) := 2\Psi_H(x) - 4\Psi_{2H+1}\!\left(\frac{x}{2}\right), \qquad \widetilde{B}_H(x) := 2B_H(x) + 4B_{2H+1}\!\left(\frac{x}{2}\right).
\end{equation}
Then every even harmonic $h = 2m$ ($1 \le |m| \le H$) cancels identically in $\widetilde{\Psi}_H(x)$, so that
\[
\widetilde{\Psi}_H(x) = \sum_{\substack{1 \le |h| \le 2H+1 \\ 2 \nmid h}} \tilde{c}_h(H)\,e\!\left(\frac{hx}{2}\right), \qquad \widetilde{B}_H(x) = \frac{2}{H + 1} + \sum_{1 \le |h| \le 2H+1} \tilde{b}_h(H)\,e\!\left(\frac{hx}{2}\right),
\]
with $|\tilde{c}_h(H)| \le \frac{2}{\pi |h|}$ and $|\tilde{b}_h(H)| < \frac{3}{H+1}$, and for all $x \in \R$:
\begin{equation}\label{eq:2.9}
\bigl|(-1)^{\lfloor x \rfloor} - \widetilde{\Psi}_H(x)\bigr| \le \widetilde{B}_H(x), \qquad \sum_{\substack{1 \le |h| \le 2H+1 \\ 2 \nmid h}} \bigl|\tilde{c}_h(H)\bigr| \le \frac{4}{\pi}\bigl(\log(2H + 1) + 1\bigr).
\end{equation}
\end{lemma}

\begin{proof}
For every $x \notin \Z$, writing $\lfloor x \rfloor = 2\lfloor x/2 \rfloor + r$ with $r \in \{0, 1\}$ shows that $2\psi(x) - 4\psi(x/2) = 1 - 2r = (-1)^{\lfloor x \rfloor}$. By the triangle inequality, $|(-1)^{\lfloor x \rfloor} - \widetilde{\Psi}_H(x)| \le 2|\psi(x) - \Psi_H(x)| + 4|\psi(x/2) - \Psi_{2H+1}(x/2)| \le \widetilde{B}_H(x)$ for all $x \notin \Z$, and at integer points $x \in \Z$ (where $B_H(0) = \frac{1}{2}$ and $|\Psi_H(m)| = 0$) we have $\widetilde{B}_H(m) \ge 2(\frac{1}{2}) = 1 = |(-1)^m - \widetilde{\Psi}_H(m)|$. In $2\Psi_H(x) - 4\Psi_{2H+1}(x/2)$, the coefficient of $e(mx) = e((2m)x/2)$ for $1 \le |m| \le H$ is $2c_m(H) - 4c_{2m}(2H + 1) = -\frac{1}{\pi i m}\bigl(W(\frac{|m|}{H+1}) - W(\frac{2|m|}{2H+2})\bigr) = 0$ because $\frac{2|m|}{(2H+1)+1} = \frac{|m|}{H+1}$. Thus $\widetilde{\Psi}_H(x)$ is supported strictly on odd frequencies $h = 2m + 1$ ($1 \le |h| \le 2H + 1$) with $\tilde{c}_h(H) = -4c_h(2H + 1)$ ($|\tilde{c}_h(H)| \le \frac{2}{\pi |h|}$), while the zero-frequency term of $\widetilde{B}_H(x)$ is $2b_0(H) + 4b_0(2H + 1) = \frac{2}{2H+2} + \frac{4}{4H+4} = \frac{2}{H+1} < \frac{3}{H+1}$.
\end{proof}

Fix any $\delta \in (0, 1/3]$ and set $R := \lfloor 1/\delta \rfloor$ and $A_R := 6R + 30$. On the range $p > n^\delta$, we have $p^{R+1} > n$, so Legendre's formula has at most $R$ non-zero terms:
\begin{equation}\label{eq:2.10}
v_p(n!) = \sum_{r=1}^R \left\lfloor \frac{n}{p^r} \right\rfloor \qquad \forall\, p > n^\delta.
\end{equation}
Near each root transition $p = n^{1/r}$ ($2 \le r \le R$), we define the transition collar
\begin{equation}\label{eq:2.11}
\mathcal{C}_r(n) := \left(\frac{n^{1/r}}{(\log n)^{A_R/r}},\, n^{1/r}\right] \qquad (2 \le r \le R),
\end{equation}
on which $1 \le \frac{n}{p^r} < (\log n)^{A_R}$, so $\lfloor n/p^r \rfloor = m_r$ takes at most $M_r(n) := \lfloor (\log n)^{A_R} \rfloor$ constant integer values on the disjoint subintervals
\begin{equation}\label{eq:2.12}
J_{r,m_r} := \left(\left(\frac{n}{m_r + 1}\right)^{1/r},\, \left(\frac{n}{m_r}\right)^{1/r}\right] \cap \mathcal{C}_r(n), \qquad 1 \le m_r \le M_r(n).
\end{equation}
On each $J_{r,m_r}$, $(-1)^{\lfloor n/p^r \rfloor} = (-1)^{m_r}$ is constant, reducing the active floors to $1 \le s \le r - 1$ (each with $n/p^s \ge n^{1-s/r} \ge n^{1/r} \ge n^\delta$), whereas on the complementary inter-collar intervals $\mathcal{D}_r(n) := \bigl(n^{1/(r+1)}, \frac{n^{1/r}}{(\log n)^{A_R/r}}\bigr]$ every active phase $n/p^s$ ($1 \le s \le r$) on a dyadic block $(P, 2P] \subset \mathcal{D}_r(n)$ satisfies $n/P^s \ge 2^{-R}(\log n)^{A_R}$.

\begin{lemma}[{Prime Exponential Sum Bound on {$[n^\delta, n/(\log n)^{12}]$}}]\label{lem:2.4}
Fix $\delta \in (0, 1/3]$ and set $R := \lfloor 1/\delta \rfloor$, $A_R := 6R + 30$, and $H := \lfloor (\log n)^{R+2} \rfloor$. Let $P \in \bigl[n^\delta, \frac{n}{(\log n)^{12}}\bigr]$, let $I \subset (P, 2P]$ be any subinterval on which the set of active floor indices $\mathcal{A}(I) \subset \{1, \dots, R\}$ is fixed with $\frac{n}{P^r} \ge 2^{-R}(\log n)^{A_R}$ for all $r \in \mathcal{A}(I)$, and let
\[
F(t) = \sum_{r \in \mathcal{A}(I)} \frac{h_r n}{2t^r}, \qquad \mathbf{h} = (h_r)_{r \in \mathcal{A}(I)} \ne \mathbf{0}, \quad |h_r| \le 2H + 1.
\]
Then for every smooth weight $w(t)$ on $[P, 2P]$ satisfying $|w(t)| \le W_0$ and $\int_P^{2P} |w'(t)|\,dt \le W_0$, we have unconditionally:
\begin{enumerate}
\item[(i)] In general, for $T_0 := \frac{|h_{r_0}|n}{P^{r_0}} \ge 2^{-R}(\log n)^{A_R}$ (where $r_0 := \min\{r \in \mathcal{A}(I) : h_r \ne 0\}$), there exists a constant $\eta'_\delta > 0$ such that
\begin{equation}\label{eq:2.13}
\left|\sum_{p \in I} w(p)(\log p)\,e\bigl(F(p)\bigr)\right| \ll_\delta W_0 P\bigl(T_0^{-1/6} + P^{-\eta'_\delta}\bigr)\log^3 P \ll_\delta W_0 P(\log n)^{-(R+2)}.
\end{equation}
\item[(ii)] Moreover, whenever $T_0 \ge n^{\delta/2}$ (which holds in particular on every collar subinterval $J_{r,m_r} \subset \mathcal{C}_r(n)$ where $\mathcal{A}(J_{r,m_r}) = \{1, \dots, r - 1\}$ and $T_0 \ge n^{1-r_0/r} \ge n^{1/r} \ge n^\delta$), there exists a constant $\eta_\delta > 0$ such that
\begin{equation}\label{eq:2.14}
\left|\sum_{p \in I} w(p)(\log p)\,e\bigl(F(p)\bigr)\right| \ll_\delta W_0 P\,n^{-\eta_\delta}.
\end{equation}
\end{enumerate}
\end{lemma}

\begin{proof}
By partial summation against $w(t)$ and discarding prime powers $p^k$ ($k \ge 2$) at cost $O(P^{1/2})$, it suffices to bound $S(I') := \sum_{k \in I'} \Lambda(k)\,e(F(k))$ for any subinterval $I' \subset I \subset (P, 2P]$. Let $r_0 := \min\{r \in \mathcal{A}(I) : h_r \ne 0\}$. Because $t \in [P, 2P]$ with $P \ge n^\delta \gg H = \lfloor (\log n)^{R+2} \rfloor$ and $|h_r| \le 2H + 1$, for every derivative order $j \ge 1$ the $r_0$-th term $\frac{h_{r_0} n}{2t^{r_0}}$ strictly dominates all higher terms $r > r_0$:
\[
\bigl|F^{(j)}(t)\bigr| = \frac{(j + r_0 - 1)!}{(r_0 - 1)!} \frac{|h_{r_0}|n}{2t^{j+r_0}} \Bigl(1 + O_R\bigl(HP^{-1}\bigr)\Bigr) \asymp_{j,R} \frac{T_0}{P^j} \qquad (\forall\, t \in [P, 2P],\; j \ge 1),
\]
where $T_0 := \frac{|h_{r_0}|n}{P^{r_0}} \in \bigl[2^{-R}(\log n)^{6R+30}, (2H + 1)n^{1-r_0\delta}\bigr]$. Apply Vaughan's identity \cite[Proposition~13.4]{IwaniecKowalski2004} with parameters $U = V = P^{1/5}$ to decompose $S(I') = S_{\mathrm{I}} - S_{\mathrm{I}}' + S_{\mathrm{II}} + O(P^{1/5})$ into:
\begin{itemize}
\item \textbf{Type I sums} ($S_{\mathrm{I}}$ with $d \le U = P^{1/5}$ and $S_{\mathrm{I}}'$ with $d \le UV = P^{2/5}$): sums of the form $\sum_{d \le D} a_d \sum_{m \in I'/d} e(F(dm))$ (and its smooth $\log m$-weighted analogue), where $|a_d| \le \tau(d)\log P$ and the inner variable $m \asymp M := P/d \ge P^{3/5}$ ranges over a subinterval of $(P/d, 2P/d]$ with smooth phase $g_d(m) := F(dm)$ satisfying $|g_d^{(j)}(m)| = d^j |F^{(j)}(dm)| \asymp_{j,R} T_0 / M^j$.
\item \textbf{Type II bilinear sums} ($U < d \le P/V$, so $P^{1/5} < d \le P^{4/5}$): dyadic bilinear sums of the form $\sum_M \sum_{d \sim P/M} \alpha_d \sum_{m \sim M,\, dm \in I'} \beta_m e(F(dm))$ with $|\alpha_d|, |\beta_m| \le \tau(d)\log P$. Applying Cauchy--Schwarz in $d \sim D := P/M \in [P^{1/5}, P^{4/5}]$ and interchanging summation gives
\[
\bigl|S_{\mathrm{II}}(D, M)\bigr|^2 \ll (D\log^3 P)\left(DM\log^2 P + \sum_{m_1 \ne m_2 \sim M} |\beta_{m_1}\beta_{m_2}| \left|\sum_{d \in I(m_1,m_2)} e\bigl(G_{m_1,m_2}(d)\bigr)\right|\right),
\]
where the differenced phase $G_{m_1,m_2}(d) := F(dm_1) - F(dm_2)$ has $j$-th derivative
\[
\bigl|G_{m_1,m_2}^{(j)}(d)\bigr| \asymp_{j,R} \frac{|m_1 - m_2|}{M} \frac{T_0}{D^j} \qquad (\forall\, d \sim D,\; j \ge 1).
\]
\end{itemize}
We distinguish two regimes for $T_0$:

\noindent\textbf{Regime A ($2^{-R}(\log n)^{6R+30} \le T_0 \le P^{3/2}$):}
In Type I, $M \ge P^{3/5}$. If $T_0 \le M$, then $|g_d''(m)| \asymp T_0 / M^2$, so van der Corput's second-derivative bound \cite[Corollary~8.13]{IwaniecKowalski2004} gives $\sum_m e(g_d(m)) \ll M(T_0/M^2)^{1/2} + (M^2/T_0)^{1/2} \ll T_0^{1/2} + MT_0^{-1/2} \ll M(P^{-1/10} + T_0^{-1/2})$. If $M < T_0 \le P^{3/2}$, then $|g_d^{(3)}(m)| \asymp T_0/M^3$ and the third-derivative van der Corput bound ($j = 3$, \cite[Corollary~8.13]{IwaniecKowalski2004}) gives $\ll M(T_0/M^3)^{1/6} + M^{1/2}(M^3/T_0)^{1/6} = T_0^{1/6}M^{1/2} + MT_0^{-1/6} \ll M(P^{-1/20} + T_0^{-1/6})$. In Type II, for the terms with $|m_1 - m_2| \le MT_0^{-1/3}$ there are $\ll M^2 T_0^{-1/3}$ pairs contributing $\ll PT_0^{-1/6}\log^3 P$ after taking the square root; for $|m_1 - m_2| > MT_0^{-1/3}$, setting $T_\Delta := \frac{|m_1 - m_2|}{M} T_0 \in (T_0^{2/3}, T_0]$ and applying the Kusmin--Landau bound \cite[Corollary~8.11]{IwaniecKowalski2004} when $T_\Delta \le \frac{1}{2}D$ and the $j$-th derivative van der Corput bound ($j \in \{2, \dots, 9\}$, using $D \ge P^{1/5}$ and $T_\Delta \le P^{3/2} < D^8$) when $T_\Delta > \frac{1}{2}D$ bounds the inner $d$-sum by $\ll D(T_0^{-1/3} + P^{-\eta'_\delta})$ for some $\eta'_\delta > 0$. Summing over $d$ and $m_1, m_2$ gives $|S(I')| \ll_\delta P(T_0^{-1/6} + P^{-\eta'_\delta})\log^3 P$. Since $T_0 \ge 2^{-R}(\log n)^{6R+30}$, $T_0^{-1/6}\log^3 P \ll_\delta (\log n)^{-(R+5)}\log^3 P \ll (\log n)^{-(R+2)}$.

\noindent\textbf{Regime B ($T_0 > P^{3/2}$):}
Choose the integer derivative order $j := \lfloor \frac{\log(T_0/M)}{\log M} \rfloor + 2 \ge 3$ in Type I (and similarly with $D$ and $T_\Delta$ in Type II), where $j \le 5R + 2$ since $T_0 \le (2H + 1)n$ and $M, D \ge P^{1/5} \ge n^{\delta/5}$. Then $M^{j-2} \le T_0/M < M^{j-1}$, so $\lambda_j := T_0/M^j \in [M^{-2}, M^{-1})$, and the $j$-th derivative Weyl--van der Corput inequality \cite[Theorem~8.25]{IwaniecKowalski2004} yields a polynomial power saving $\ll_\delta P\,n^{-\eta_\delta}$ with $\eta_\delta > 0$, proving both \eqref{eq:2.13} and \eqref{eq:2.14}.
\end{proof}

We now state and prove the principal valuation-parity equidistribution theorem.

\begin{theorem}[{Valuation-Parity Equidistribution and Second-Order Discrepancy Law}]\label{thm:2.5}
For every integer $n \ge 4$, the following statements hold unconditionally:
\begin{enumerate}
\item[(i)] \textbf{Second-Order Discrepancy Formula:}
\begin{equation}\label{eq:2.15}
\Dparity(n) = \bigl(\log 2 - 1 - \Sparitys(n)\bigr)\,n + O\!\left(\frac{n}{\log n}\right).
\end{equation}
\item[(ii)] \textbf{Collar-Decomposed Cancellation and Full-Range 50--50 Weight Split:} For every fixed $\delta \in (0, 1/3]$,
\begin{equation}\label{eq:2.16}
\sum_{n^\delta < p \le \sqrt{n}} (-1)^{v_p(n!)-1} \frac{\log p}{p - 1} \ll_\delta \frac{1}{\log n},
\end{equation}
and across the full range $3 \le p \le n$, $\Sparitys(n) = o(\log n)$, $\Dparity(n) = o(n\log n)$, and
\begin{equation}\label{eq:2.17}
\Sodd(n) = \frac{1}{2}n\log n + o(n\log n), \qquad \Seven(n) = \frac{1}{2}n\log n + o(n\log n).
\end{equation}
\item[(iii)] \textbf{Multi-Modulus Power-Band Equidistribution:} For every integer modulus $h \ge 2$, every residue class $c \bmod h$, and every exponent band $0 \le a < b \le 1$,
\begin{equation}\label{eq:2.18}
\sum_{\substack{n^a < p \le n^b \\ v_p(n!) \equiv c \pmod h}} v_p(n!)\log p = \left(\frac{b - a}{h} + o(1)\right)n\log n.
\end{equation}
\end{enumerate}
\end{theorem}

\begin{proof}
(i) Set $R_1 := \lfloor (\log n)^{12} \rfloor$. On $(\frac{n}{R_1}, n]$, Lemma~\ref{lem:2.1}(i) and \eqref{eq:2.5} give
\begin{align*}
\sum_{n/R_1 < p \le n} (-1)^{v_p(n!)} v_p(n!)\log p &= \sum_{r=1}^{R_1-1} (-1)^r r\left(\vartheta\!\left(\frac{n}{r}\right) - \vartheta\!\left(\frac{n}{r + 1}\right)\right) \\
&= n \sum_{r=1}^{R_1-1} \frac{(-1)^r}{r + 1} + O\!\left(\frac{n}{(\log n)^{12}}\right) = (\log 2 - 1)\,n + O\!\left(\frac{n}{(\log n)^{12}}\right).
\end{align*}
On $(\sqrt{n}, \frac{n}{R_1}]$, writing $v_p(n!) = \lfloor n/p \rfloor = \frac{n}{p} - \{n/p\}$ (where $\{n/p\}$ contributes $\ll \vartheta(n/R_1) \ll \frac{n}{(\log n)^{12}}$) and applying Lemma~\ref{lem:2.3} with $H = \lfloor (\log n)^5 \rfloor$ and Lemma~\ref{lem:2.4}(i) ($\delta = 1/3$, $T_0 \ge R_1 \gg (\log n)^{12}$) with $w(t) = \frac{n}{t}$ ($W_0 = \frac{n}{P}$) across $O(\log n)$ dyadic intervals $(P, 2P]$ bounds the intermediate sum by $O\bigl(\frac{n}{\log n}\bigr)$. Combining this with Lemma~\ref{lem:2.2} establishes \eqref{eq:2.15}.

(ii)--(iii) Fix any $\delta \in (0, 1/3]$ and set $R := \lfloor 1/\delta \rfloor$, $A_R := 6R + 30$, and $H := \lfloor (\log n)^{R+2} \rfloor$. Partition $(n^\delta, \sqrt{n}]$ into the $R - 1$ transition collars $\mathcal{C}_r(n) = \bigl(\frac{n^{1/r}}{(\log n)^{A_R/r}}, n^{1/r}\bigr]$ ($2 \le r \le R$) of \eqref{eq:2.11} and the complementary inter-collar intervals $\mathcal{D}_r(n) := \bigl(n^{1/(r+1)}, \frac{n^{1/r}}{(\log n)^{A_R/r}}\bigr]$ ($2 \le r \le R - 1$, plus $(n^\delta, n^{1/R}]$ if $\delta < 1/R$).

On each collar $\mathcal{C}_r(n)$, partitioning into the $M_r(n) \le (\log n)^{A_R}$ subintervals $J_{r,m_r}$ of \eqref{eq:2.12} fixes $\lfloor n/p^r \rfloor = m_r$, while $\lfloor n/p^s \rfloor = 0$ for $s > r$. Expanding the $r - 1$ active floors $(-1)^{\lfloor n/p^s \rfloor}$ ($1 \le s \le r - 1$) via Lemma~\ref{lem:2.3} produces a main product $\prod_{s=1}^{r-1} \widetilde{\Psi}_H(n/p^s)$ and an error product whose zero-frequency constant term is bounded by $\ll_R \frac{1}{H}$ (since $|\widetilde{\Psi}_H(x)| \le 1 + \widetilde{B}_H(x)$). Summing the zero-frequency error over $p \in \mathcal{C}_r(n)$ contributes $\ll_R \frac{1}{H} \sum_{p \in \mathcal{C}_r(n)} \frac{\log p}{p-1} \ll_R \frac{\log\log n}{(\log n)^{R+2}} \ll (\log n)^{-1}$. For every non-zero frequency tuple $\mathbf{h} = (h_1, \dots, h_{r-1}) \in \Z^{r-1} \setminus \{\mathbf{0}\}$ ($|h_s| \le 2H + 1$), the active index $r_0 := \min\{1 \le s \le r - 1 : h_s \ne 0\}$ satisfies $r_0 \le r - 1$, so for $p \in \mathcal{C}_r(n)$ the phase parameter $T_0 = \frac{|h_{r_0}|n}{P^{r_0}}$ obeys the polynomial lower bound
\[
T_0 \ge \frac{n}{(n^{1/r})^{r_0}} = n^{1 - r_0/r} \ge n^{1/r} \ge n^\delta.
\]
Hence every non-zero frequency tuple on $J_{r,m_r}$ lies in the polynomial power-saving regime of Lemma~\ref{lem:2.4}(ii) with $w(t) = \frac{1}{t-1}$ ($W_0 \asymp P^{-1}$), contributing $\ll_\delta n^{-\eta_\delta}$ per subinterval and per tuple. Summing over all $M_r(n) \le (\log n)^{A_R}$ subintervals and all $(4H + 3)^{r-1} \ll_R (\log n)^{(R-1)(R+2)}$ tuples gives a total oscillatory contribution of $\ll_\delta (\log n)^{A_R + (R-1)(R+2)} n^{-\eta_\delta} \ll_\delta (\log n)^{-1}$.

On each inter-collar interval $\mathcal{D}_r(n)$, covered by $O(\log n)$ dyadic intervals $(P, 2P]$, we expand $\prod_{s=1}^r (-1)^{\lfloor n/p^s \rfloor}$ via Lemma~\ref{lem:2.3}. The zero-frequency constant term of the error product is bounded by $\ll_R \frac{1}{H}$, which summed over $p \in \mathcal{D}_r(n)$ contributes $\ll_R \frac{\log n}{(\log n)^{R+2}} \le (\log n)^{-3}$. For the non-zero frequency tuples $\mathbf{h} \in \Z^r \setminus \{\mathbf{0}\}$, the $L^1$ norm of the coefficients in $\prod_{s=1}^r \widetilde{\Psi}_H(n/p^s)$ (and in the mixed $\widetilde{B}_H$ terms) is $\ll_R (\log H)^r \ll_R (\log\log n)^r$. Applying Lemma~\ref{lem:2.4}(i) ($W_0 P \asymp 1$) to each non-zero tuple $\mathbf{h}$ and multiplying by the $L^1$ coefficient norm and the $O(\log n)$ dyadic blocks gives $\ll_\delta (\log\log n)^R (\log n)(\log n)^{-(R+2)} \ll_\delta (\log n)^{-1}$, proving \eqref{eq:2.16}. Combining \eqref{eq:2.16} with Mertens' bound $\bigl|\sum_{3 \le p \le n^\delta} (-1)^{v_p(n!)-1} \frac{\log p}{p-1}\bigr| \le \delta\log n + O(1)$ and letting $\delta \downarrow 0$ yields $\Sparitys(n) = o(\log n)$, $\Dparity(n) = o(n\log n)$, and \eqref{eq:2.17}. Replacing $(-1)^{v_p(n!)}$ by the additive characters $e(z v_p(n!)/h)$ ($1 \le z \le h - 1$) via $\mathbf{1}_{v_p(n!) \equiv c \pmod h} = \frac{1}{h} \sum_{z=0}^{h-1} e(-zc/h)e(zv_p(n!)/h)$ on $(n^a, n^b]$ identically establishes \eqref{eq:2.18}.
\end{proof}

\section{Chebyshev--Lucas Valuation Absorption and Proofs of Theorems~\ref{thm:1.1} and \ref{thm:1.2}}\label{sec:3}

We now apply the valuation-parity equidistribution of Theorem~\ref{thm:2.5} to the Pell embedding $m + s_n\sqrt{q_n} = \varepsilon_n^k$ in $K_n = \Q(\sqrt{q_n})$ and prove Theorem~\ref{thm:1.1} and Theorem~\ref{thm:1.2}.

\subsection{Basic Invariants of $K_n = \Q(\sqrt{q_n})$ and Chebyshev--Lucas Valuation Laws}

\begin{lemma}[{Fundamental Unit Norm and Growth of $q_n, s_n$}]\label{lem:3.1}
For every integer $n \ge 3$ with $n \ne 6$:
\begin{enumerate}
\item[(i)] The square-free integer $q_n$ has at least one prime divisor $p \equiv 3 \pmod 4$. Consequently, the negative Pell equation $x^2 - q_n y^2 = -1$ has no rational or integer solutions, and both the fundamental unit $\varepsilon_n = X_1(q_n) + Y_1(q_n)\sqrt{q_n} > 1$ of $\Z[\sqrt{q_n}]$ and the fundamental unit $\eta_n > 1$ of $\mathcal{O}_{K_n}$ have norm $\mathrm{N}_{K_n/\Q}(\varepsilon_n) = \mathrm{N}_{K_n/\Q}(\eta_n) = +1$, with $[\mathcal{O}_{K_n}^\times : \Z[\sqrt{q_n}]^\times] = d_n \in \{1, 3\}$ ($d_n = 1$ unless $q_n \equiv 5 \pmod 8$). Moreover, for every $n \ge 116$ (and for $8 \le n \le 115$ with $n \notin \{10, 11, 28\}$), $\log q_n > 0.4\,n$, while for all $n \ge 8$, $\log \varepsilon_n > 0.2\,n$ and any solution $(n, m)$ to $n! + 1 = m^2$ satisfies the explicit bound $1 \le k < 2.5\log n$ (with $2.5\log n < \frac{n}{5}$ for all $n \ge 116$).
\item[(ii)] As $n \to \infty$, $\log(s_n\sqrt{q_n}) = \frac{1}{2}\log(n!) = \frac{1}{2}n\log n - \frac{1}{2}n + \frac{1}{4}\log(2\pi n) + O(n^{-1})$ and $\log s_n = \frac{1}{2}n\log n - \frac{1+\log 2}{2}n + O\bigl(ne^{-c_0\sqrt{\log n}}\bigr)$.
\end{enumerate}
\end{lemma}

\begin{proof}
(i) For $n \in \{3, 4, 5\}$, $v_3(n!) = 1$, so $3 \mid q_n$; for $7 \le n \le 13$, $v_7(n!) = 1$, so $7 \mid q_n$. For $n \ge 14$, we have $x := n/2 \ge 7$, so Breusch's generalization of Bertrand's postulate to arithmetic progressions modulo 4 \cite{Breusch1932} guarantees the existence of a prime $p \in (n/2, n]$ with $p \equiv 3 \pmod 4$; for any such prime, $v_p(n!) = 1$, so $p \mid q_n$. Reducing $x^2 - q_n y^2 = -1$ modulo $p \equiv 3 \pmod 4$ ($p \mid q_n$) gives $\bigl(\frac{-1}{p}\bigr) = +1$, a contradiction, so $\mathrm{N}(\varepsilon_n) = \mathrm{N}(\eta_n) = +1$. Furthermore, since every prime $p \in (n/2, n]$ divides $q_n$, Rosser and Schoenfeld's explicit Chebyshev bounds \cite[Theorem~10]{RosserSchoenfeld1962} give $\log q_n \ge \vartheta(n) - \vartheta(n/2) > 0.4\,n$ for all $n \ge 116$ (and direct check verifies $\log q_n > 0.4\,n$ for $8 \le n \le 115$ with $n \notin \{10, 11, 28\}$, while at $n \in \{10, 11, 28\}$ direct evaluation gives $\log \varepsilon_{10} \approx 2.7687 > 2.0$, $\log \varepsilon_{11} \approx 6.5539 > 2.2$, and $\log \varepsilon_{28} \approx 31.7026 > 5.6$). Thus $\log \varepsilon_n > 0.2\,n$ and $k = \frac{\log(m + s_n\sqrt{q_n})}{\log \varepsilon_n} < \frac{\frac{1}{2}n\log n}{0.2\,n} = 2.5\log n$ for all $n \ge 8$ (with $2.5\log n < \frac{n}{5}$ for all $n \ge 116$). Part (ii) follows from Stirling's formula and Lemma~\ref{lem:2.1}(ii).
\end{proof}

Define the Lucas--Pell sequence $(X_r(q_n), Y_r(q_n))_{r \ge 1}$ by $X_r(q_n) + Y_r(q_n)\sqrt{q_n} := \varepsilon_n^r$.

\begin{lemma}[{Chebyshev--Lucas Factorization and $p$-Adic Lifting the Exponent}]\label{lem:3.2}
For every integer $k \ge 1$, we have $X_k(q_n) = T_k(X_1(q_n))$ and $Y_k(q_n) = Y_1(q_n)U_{k-1}(X_1(q_n))$, where $U_{k-1}(x) \in \Z[x]$ is the Chebyshev polynomial of the second kind:
\begin{equation}\label{eq:3.1}
U_{k-1}\bigl(X_1(q_n)\bigr) = \sum_{j=0}^{\lfloor (k-1)/2 \rfloor} \binom{k}{2j + 1} X_1(q_n)^{k-1-2j}\bigl(q_n Y_1(q_n)^2\bigr)^j.
\end{equation}
Moreover, $\gcd\bigl(X_1(q_n)^2 - 1, U_{k-1}(X_1(q_n))^2\bigr)$ divides $k^2$, and for every prime $p$ dividing $q_n Y_1(q_n)$ satisfying $p \ne 3$, $3 \nmid k$, or $3 \mid Y_1(q_n)$,
\begin{equation}\label{eq:3.2}
v_p\Bigl(U_{k-1}\bigl(X_1(q_n)\bigr)\Bigr) = v_p(k).
\end{equation}
In particular, for every prime $p \in \Podd(n) \setminus \{3\}$ in $m + s_n\sqrt{q_n} = \varepsilon_n^k$,
\begin{equation}\label{eq:3.3}
v_p\bigl(q_n Y_1(q_n)^2\bigr) = v_p(n!) - 2v_p(k).
\end{equation}
\end{lemma}

\begin{proof}
Expanding $\varepsilon_n^k = (X_1(q_n) + Y_1(q_n)\sqrt{q_n})^k$ by the binomial theorem and substituting $q_n Y_1(q_n)^2 = X_1(q_n)^2 - 1$ gives \eqref{eq:3.1}. Let $p$ be any prime dividing $M := q_n Y_1(q_n)^2 = X_1(q_n)^2 - 1$, so $\gcd(p, X_1(q_n)) = 1$. In \eqref{eq:3.1}, the $j = 0$ term is $k X_1(q_n)^{k-1}$ of $p$-adic valuation $v_p(k)$. For every $j \ge 1$, the identity $(2j + 1)\binom{k}{2j+1} = k\binom{k-1}{2j}$ implies
\[
v_p\!\left(\binom{k}{2j + 1}M^j\right) \ge v_p(k) - v_p(2j + 1) + j\,v_p(M) \ge v_p(k) + \bigl(j - v_p(2j + 1)\bigr).
\]
Whenever $p \ne 3$ or $3 \mid Y_1(q_n)$, we have $j\,v_p(M) > v_p(2j + 1)$ for all $j \ge 1$ (since $v_2(2j + 1) = 0 < j$, $p^j \ge 5^j > 2j + 1$ for $p \ge 5$, and $v_3(M) \ge 2$ gives $2j > v_3(2j + 1)$ when $p = 3$ and $3 \mid Y_1(q_n)$), so every $j \ge 1$ term has valuation strictly exceeding $v_p(k)$; and when $p = 3$ and $3 \nmid k$, we have $3 \mid (k - 1)(k - 2)$ so $v_3(\binom{k}{3}M) \ge 1 > 0 = v_3(k)$, while $j > v_3(2j + 1)$ for all $j \ge 2$. This proves \eqref{eq:3.2}. For $p \in \Podd(n) \setminus \{3\}$, $v_p(q_n) = 1$, so $p \mid q_n Y_1(q_n)$ and taking valuations in $n! = q_n Y_1(q_n)^2 U_{k-1}(X_1(q_n))^2$ yields \eqref{eq:3.3}.
\end{proof}

\subsection{Proof of Theorem~\ref{thm:1.1}}

\begin{proof}[Proof of Theorem~\ref{thm:1.1}]
For $n \in \{4, 5, 7\}$, direct evaluation gives $q_4 = 6$, $\varepsilon_4 = 5 + 2\sqrt{6}$, $k(4) = 1$; $q_5 = 30$, $\varepsilon_5 = 11 + 2\sqrt{30}$, $k(5) = 1$; and $q_7 = 35$, $\varepsilon_7 = 6 + \sqrt{35}$, $71 + 12\sqrt{35} = \varepsilon_7^2$, $k(7) = 2$.

Now let $(n, m)$ be a solution with $n \ge 8$. Since $\varepsilon_n = X_1(q_n) + Y_1(q_n)\sqrt{q_n} > 2Y_1(q_n)\sqrt{q_n}$ and $\varepsilon_n^k = m + s_n\sqrt{q_n} = (2 + O(1/n!))\sqrt{n!}$, taking logarithms gives
\begin{equation}\label{eq:3.4}
\log\bigl(q_n Y_1(q_n)^2\bigr) = 2\log\bigl(Y_1(q_n)\sqrt{q_n}\bigr) < 2\log \varepsilon_n = \frac{1}{k}\log(n!) + O(1).
\end{equation}
On the other hand, summing \eqref{eq:3.3} over all $p \in \Podd(n) \setminus \{3\}$ (and noting $v_3(n!)\log 3 = O(n)$) gives
\begin{equation}\label{eq:3.5}
\log\bigl(q_n Y_1(q_n)^2\bigr) \ge \Sodd(n) - v_3(n!)\log 3 - 2\log k = \Sodd(n) - O(n).
\end{equation}
Since $k < 2.5\log n$ by Lemma~\ref{lem:3.1}(i), Theorem~\ref{thm:2.5}(ii) ($\Sodd(n) = (\frac{1}{2} + o(1))\log(n!)$) combined with \eqref{eq:3.4}--\eqref{eq:3.5} yields $\frac{1}{2} - o(1) \le \frac{1}{k}$, forcing $k \in \{1, 2\}$ for all sufficiently large $n$. Taking logarithms of $\varepsilon_n^k = (2 + o(1))(n!)^{1/2}$ gives \eqref{eq:1.7}.

Finally, suppose that $q_{n_1} = q_{n_2} = D$ for two solutions $4 \le n_1 < n_2$. Then $\frac{n_2!}{n_1!} = \prod_{j=n_1+1}^{n_2} j = (s_{n_2}/s_{n_1})^2$ is a perfect square. By the Erd\H{o}s--Selfridge theorem \cite{ErdosSelfridge1975}, a product of two or more consecutive positive integers is never a square, forcing $n_2 = n_1 + 1$ and $n_1 + 1 = t^2$ for some integer $t \ge 3$. However, $(n_1 + 1)n_1! + 1 = m_2^2$ and $n_1! + 1 = m_1^2$ give $(tm_1 - m_2)(tm_1 + m_2) = t^2 - 1$, which is impossible since $tm_1 - m_2 \ge 1$ forces $tm_1 + m_2 \le t^2 - 1$ whereas $tm_1 + m_2 > t\sqrt{n_1! + 1} > t^2 - 1$. (For sufficiently large $n_1$, where $k_{n_1} \in \{1, 2\}$, one also obtains $n_2 - n_1 > 0.49\,n_1 > \frac{n_1}{16597}$ directly from $2(k_{n_2} - k_{n_1})\log \varepsilon_D \ge \frac{1}{2}\log(n_1!)$, contradicting Schoenfeld's prime gap theorem \cite{Schoenfeld1976}.) Thus $q_{n_1} \ne q_{n_2}$ for all $4 \le n_1 < n_2$.
\end{proof}

\subsection{Proof of Theorem~\ref{thm:1.2} and Imaginary Quadratic Picard 2-Torsion}

\begin{proof}[Proof of Theorem~\ref{thm:1.2}]
Let $(n, m)$ be a solution to $n! + 1 = m^2$ with $n \ge 116$. Let $\eta_n > 1$ be the fundamental unit of $\mathcal{O}_{K_n}$ ($\varepsilon_n = \eta_n^{d_n}$ with $d_n \in \{1, 3\}$), and let $\widetilde{R}_p(K_n) := \frac{\log_p(\eta_n)}{\sqrt{D_n}} \in \Z_p$ be the normalized $p$-adic regulator of $K_n = \Q(\sqrt{q_n})$. By Nagura's theorem \cite{Nagura1952} (which guarantees a prime in $(x, \frac{6}{5}x)$ for all $x \ge 25$, applied at $x = n/5$ for $n \ge 125$, and verified directly for $116 \le n \le 124$), the interval $I_4 = (n/5, n/4] \subset (n/5, n/3]$ is non-empty for every $n \ge 116$, and every prime $p \in (n/5, n/3]$ satisfies $p > n/5 \ge 23.2 > \sqrt{n}$, with $v_p(n!) = 4$ on $I_4 = (n/5, n/4]$ and $v_p(n!) = 3$ on $I_3 = (n/4, n/3]$. Moreover, by Lemma~\ref{lem:3.1}(i), $1 \le k < 2.5\log n < \frac{n}{5} < p$, so $p \nmid k$ unconditionally for every $n \ge 116$ and every $p \in (n/5, n/3]$ (and for all sufficiently large $n$, $k \in \{1, 2\}$ by Theorem~\ref{thm:1.1}, so $p \nmid k$ for all primes $p \ge 5$).

(i) For any prime $p \in [5, n/2]$ with $p \nmid k$, we have $v_p(2m d_n k) = 0$. In $\Q_p(\sqrt{q_n})$,
\[
\eta_n^{2d_n k} = \varepsilon_n^{2k} = \bigl(m + s_n\sqrt{q_n}\bigr)^2 = \bigl(1 + 2q_n s_n^2\bigr) + 2ms_n\sqrt{q_n}.
\]
Since $p \le n/2 \implies v_p(n!) \ge 2 \implies v_p(s_n) = \lfloor v_p(n!)/2 \rfloor \ge 1$ and $p \ge 5$, the $p$-adic logarithm $\log_p(1 + z) = \sum_{j=1}^\infty \frac{(-1)^{j-1}}{j}z^j$ is valuation-preserving on $z = 2q_n s_n^2 + 2ms_n\sqrt{q_n} \in p\mathcal{O}_{\Q_p(\sqrt{q_n})}$. Since $v_p(2q_n s_n^2) = v_p(n!) > \lfloor v_p(n!)/2 \rfloor + \frac{1}{2}v_p(q_n) = v_p(2ms_n\sqrt{q_n})$, the ultrametric triangle inequality gives
\begin{equation}\label{eq:3.6}
v_p\bigl(\widetilde{R}_p(K_n)\bigr) = v_p(s_n) = \left\lfloor \frac{v_p(n!)}{2} \right\rfloor \qquad \forall\, p \in [5, n/2],\; p \nmid k.
\end{equation}
In particular, every prime $p \in [5, n/4]$ with $p \nmid k$ satisfies $v_p(n!) \ge 4 \implies v_p(\widetilde{R}_p(K_n)) \ge 2 \implies \eta_n^{p^2-1} \equiv 1 \pmod{p^2\mathcal{O}_{K_n}}$ (so $p$ is a Wieferich prime for $\eta_n$).

By Leopoldt's $p$-adic analytic class number formula \cite[Theorem~5.24]{Washington1997},
\begin{equation}\label{eq:3.7}
\left(1 - \frac{\chi_{D_n}(p)}{p}\right)^{-1} L_p(1, \chi_{D_n}) = \frac{2h(q_n)R_p(K_n)}{\sqrt{D_n}} = 2h(q_n)\,\widetilde{R}_p(K_n).
\end{equation}
Simultaneously, by Coates's formula \cite[Appendix, Lemma~8]{Coates1977} (see also \cite[Proposition~1]{Taya1996} and \cite[Theorem~3.1]{QiStokes2024}), where $M_p(F)$ denotes the maximal abelian pro-$p$ extension of a field $F$ unramified outside $p$, the abelian $p$-ramification torsion group $\mathcal{T}_{K_n, p} = \Gal(M_p(K_n)/K_{n,\infty}) \simeq X_p(K_{n,\infty})_\Gamma$ has $p$-adic valuation
\[
v_p\bigl(\#\mathcal{T}_{K_n, p}\bigr) = v_p\!\left(w_1\bigl(K_n(\mu_p)\bigr)\,h(q_n)\,\widetilde{R}_p(K_n)\prod_{\mathfrak{p} \mid p}\bigl(1 - \mathrm{N}(\mathfrak{p})^{-1}\bigr)\right),
\]
where $v_p(w_1(K_n(\mu_p))) = 1$ for $p \ge 5$. We now evaluate both formulas according to the parity of $v_p(n!)$:
\begin{itemize}
\item When $p \in \Peven(n) \cap [5, n/2]$ ($2 \mid v_p(n!)$), $p \nmid D_n$ is unramified in $K_n$, so $\chi_{D_n}(p) \in \{\pm 1\}$, $v_p(1 - \chi_{D_n}(p)p^{-1}) = -1$, and $v_p(\prod_{\mathfrak{p} \mid p}(1 - \mathrm{N}(\mathfrak{p})^{-1})) = -2$ (both when $p$ splits with $\mathrm{N}(\mathfrak{p}) = p$ and when $p$ is inert with $\mathrm{N}(\mathfrak{p}) = p^2$). Thus \eqref{eq:3.6} and \eqref{eq:3.7} give
\[
v_p\bigl(L_p(1, \chi_{D_n})\bigr) = v_p\bigl(\#\mathcal{T}_{K_n, p}\bigr) = v_p\bigl(h(q_n)\bigr) + \frac{v_p(n!)}{2} - 1 = v_p\bigl(h(q_n)\bigr) + \left\lceil \frac{v_p(n!)}{2} \right\rceil - 1.
\]
\item When $p \in \Podd(n) \cap [5, n/2]$ ($2 \nmid v_p(n!)$), $p \mid D_n$ is ramified in $K_n$, so $\chi_{D_n}(p) = 0$, $v_p(1 - \chi_{D_n}(p)p^{-1}) = 0$, and $p\mathcal{O}_{K_n} = \mathfrak{p}^2$ gives $v_p(\prod_{\mathfrak{p} \mid p}(1 - \mathrm{N}(\mathfrak{p})^{-1})) = v_p(1 - p^{-1}) = -1$. Thus \eqref{eq:3.6} and \eqref{eq:3.7} give
\[
v_p\bigl(L_p(1, \chi_{D_n})\bigr) = v_p\bigl(\#\mathcal{T}_{K_n, p}\bigr) = v_p\bigl(h(q_n)\bigr) + \frac{v_p(n!) - 1}{2} = v_p\bigl(h(q_n)\bigr) + \left\lceil \frac{v_p(n!)}{2} \right\rceil - 1.
\]
\end{itemize}
This establishes the unified identity \eqref{eq:1.8} across all primes $p \in [5, n/2]$ with $p \nmid k$. Whenever $p \in [5, n/3]$ with $p \nmid k$, we have $v_p(n!) \ge 3$, so $\lceil v_p(n!)/2 \rceil - 1 \ge 1$ and $v_p(L_p(1, \chi_{D_n})) = v_p(\#\mathcal{T}_{K_n, p}) \ge 1$, proving that $K_n$ is non-$p$-rational ($\mathcal{T}_{K_n, p} \ne 0$). By Iwasawa's construction \cite[Theorem~7.10]{Washington1997}, $L_p(1 - s, \chi_{D_n}) = G_p((1 + p)^s - 1, \chi_{D_n})$ for $G_p(T, \chi_{D_n}) \in \Z_p[[T]]$. By the Ferrero--Washington theorem \cite{FerreroWashington1979}, $\mu_p(\chi_{D_n}) = 0$, so $G_p(T, \chi_{D_n}) = U(T)P_p(T)$ with $U(T) \in \Z_p[[T]]^\times$ and a distinguished polynomial $P_p(T) \in \Z_p[T]$ of degree $\lambda_p(\chi_{D_n})$. By the Mazur--Wiles Iwasawa Main Conjecture \cite{MazurWiles1984} (see also \cite{OzakiTaya1998,Taya1996}), $P_p((1+p)(1+T)^{-1}-1)$ generates the characteristic ideal of the $p$-ramified Iwasawa module $X_p(K_{n,\infty}) = \Gal(M_p(K_{n,\infty})/K_{n,\infty})$. Since Leopoldt's conjecture holds for the abelian field $K_n$, Greenberg's theorem \cite[Proposition~3]{Greenberg1978} (cf.\ \cite[Proposition~4.7]{Yanagisawa2026}) ensures that $X_p(K_{n,\infty})$ has no non-trivial finite $\Lambda$-submodule, so $X_p(K_{n,\infty})$ is a free $\Z_p$-module of rank $\lambda_p^*(K_n) = \lambda_p(\chi_{D_n}) = \deg P_p$. Since $X_p(K_{n,\infty})_\Gamma \simeq \mathcal{T}_{K_n, p} \ne 0$ (or equivalently, since evaluating $L_p(1-s, \chi_{D_n}) = G_p((1+p)^s - 1, \chi_{D_n})$ at $s = 0$, i.e.\ $T = 0$, gives $v_p(P_p(0)) = v_p(L_p(1, \chi_{D_n})) \ge 1$), we have $\lambda_p^*(K_n) = \deg P_p \ge 1$.

(ii)--(iii) Suppose now that $p \in \Peven(n) \cap [5, n/4]$ with $p \nmid k$ splits in $K_n$ ($\chi_{D_n}(p) = +1$). Because the primes $\mathfrak{p}, \bar{\mathfrak{p}} \mid p$ of $K_n$ have completion $K_{n,\mathfrak{p}} \simeq \Q_p$ of degree $1$ and $p \ge 5$, the local principal unit group $U_{\mathfrak{p}}^{(1)} = 1 + p\Z_p \simeq \Z_p$ is torsion-free, so the inertia subgroup of $\mathfrak{p}$ in $M_p(K_n)/K_n$ injects into $\Gamma = \Gal(K_{n,\infty}/K_n)$ and $M_p(K_n) \subseteq L(K_{n,\infty})$ is unramified over $K_{n,\infty}$ \cite[Proposition~1]{Ozaki1997}, \cite[Lemmas~6.1--6.2]{Yanagisawa2026}. Here $L(K_{n,\infty})$ denotes the maximal unramified abelian pro-$p$ extension of $K_{n,\infty}$, $K_{n,m}$ the $m$-th layer of $K_{n,\infty}/K_n$ of degree $p^m$, $A_m$ the $p$-Sylow subgroup of $\Cl(K_{n,m})$, and $D_m \subseteq A_m$ the subgroup generated by the prime ideals above $p$. Together with Taya's ambiguous class number theorem \cite[Corollary~2 and Proposition~1]{Taya1996} (with $n_2 = v_p(\widetilde{R}_p(K_n)) = \frac{1}{2}v_p(n!) \ge 2$), this gives
\[
\bigl|X(K_{n,\infty})_\Gamma\bigr| = \bigl|A_m^\Gamma\bigr| = \bigl|\mathcal{T}_{K_n, p}\bigr| = p^{v_p(h(q_n)) + \frac{1}{2}v_p(n!) - 1} \ge p^{\frac{1}{2}v_p(n!) - 1} \ge p \qquad \forall\, m \ge \frac{1}{2}v_p(n!) - 1.
\]
Hence $X(K_{n,\infty}) \ne 0$ and $\lambda_p(K_n) + \nu_p(K_n) \ge \frac{1}{2}v_p(n!) - 1 \ge 1$ (cf.\ \cite[Theorem~1]{FukudaKomatsu1986}, \cite[Theorem~6]{Jaulent2024}, \cite[Corollary~6.4]{Yanagisawa2026}). Moreover, letting $n_0^{(r)}(K_n)$ and $n_2^{(r)}(K_n)$ denote Taya's $r$-th layer $p$-unit and unit invariants \cite[§2]{Taya1996}, Taya's criterion \cite[Corollary~3 and Theorem~4]{Taya1996} gives $|D_m| \ge |D_r|\,p^{n_2^{(r)}(K_n) - n_0^{(r)}(K_n)}$ for $m \ge n_2^{(r)} - 1$, and $\lambda_p(K_n) = 0$ holds if and only if $A_0$ capitulates in $K_{n,\infty}$ and $n_0^{(r)}(K_n) = r + 1$ for some $r \ge 0$; thus $n_0^{(r)}(K_n) > r + 1$ for all $r \ge 0$ forces $\lambda_p(K_n) \ge 1$. Finally, if any of the hypotheses \textup{(a)}, \textup{(b)}, or \textup{(c)} in part~(iii) holds for $n \ge 116$, it contradicts \eqref{eq:1.8} or part~(ii), ruling out all solutions with $n \ge 116$ and leaving only $(4, 5)$, $(5, 11)$, and $(7, 71)$.
\end{proof}

We conclude this section by showing that the dichotomy $k \in \{1, 2\}$ of Theorem~\ref{thm:1.1} admits a dual characterization in the Picard group of the imaginary quadratic order $\mathcal{O}_n^- := \Z[\sqrt{-q_n}]$.

\begin{proposition}[{Imaginary Quadratic Picard 2-Torsion Dichotomy}]\label{prop:3.3}
Let $(n, m)$ be a solution to $n! + 1 = m^2$ with $n \ge 4$ and $k \in \{1, 2\}$ (which holds for $n \in \{4, 5, 7\}$ and all sufficiently large $n$ by Theorem~\ref{thm:1.1}). In the imaginary quadratic order $\mathcal{O}_n^- = \Z[\sqrt{-q_n}]$, the principal ideal $(1 + s_n\sqrt{-q_n})$ is the square of a unique invertible integral ideal $\mathfrak{J}_n \subset \mathcal{O}_n^-$ of norm $m$, defining a 2-torsion class $[\mathfrak{J}_n] \in \Pic(\mathcal{O}_n^-)[2]$ such that
\begin{equation}\label{eq:3.8}
[\mathfrak{J}_n] = 1 \text{ in } \Pic(\mathcal{O}_n^-) \iff k = 2, \qquad [\mathfrak{J}_n] = [\mathfrak{d}_2] \ne 1 \text{ in } \Pic(\mathcal{O}_n^-) \iff k = 1,
\end{equation}
where $\mathfrak{d}_2 = (\delta_2(n), \sqrt{-q_n})$ is the ambiguous ramified ideal above $\delta_2(n) = \core((m + 1)/2)$. Moreover, in the $k = 2$ regime, the second-stage ideal $\mathfrak{I}_n \subset \mathcal{O}_{\Q(\sqrt{-q_n})}$ satisfying $(1 + Y_1(q_n)\sqrt{-q_n})_{\mathrm{odd}} = \mathfrak{I}_n^2$ has exact order 2 ($[\mathfrak{I}_n] \ne 1$) in $\Cl(\Q(\sqrt{-q_n}))[2]$ unconditionally.
\end{proposition}

\begin{proof}
In $\mathcal{O}_n^- = \Z[\sqrt{-q_n}]$, $(1 + s_n\sqrt{-q_n})(1 - s_n\sqrt{-q_n}) = 1 + q_n s_n^2 = m^2$. Since $m$ is odd and $(1 + s_n\sqrt{-q_n}) + (1 - s_n\sqrt{-q_n}) = 2$, the conjugate ideals $(1 + s_n\sqrt{-q_n})$ and $(1 - s_n\sqrt{-q_n})$ are coprime and prime to the conductor of $\mathcal{O}_n^-$, so $(1 + s_n\sqrt{-q_n}) = \mathfrak{J}_n^2$ for a unique invertible ideal $\mathfrak{J}_n \subset \mathcal{O}_n^-$ of norm $m$. Since $q_n \ge 6$ is square-free, $\mathcal{O}_n^{-\times} = \{\pm 1\}$, so $\mathfrak{J}_n = (u + v\sqrt{-q_n})$ is principal if and only if $1 + s_n\sqrt{-q_n} = \pm((u^2 - q_n v^2) + 2uv\sqrt{-q_n})$. By Lemma~\ref{lem:3.1}(i), $u^2 - q_n v^2 = -1$ is impossible, so the sign is $+$, giving $u^2 - q_n v^2 = 1$, $s_n = 2|uv|$, and $m + s_n\sqrt{q_n} = (|u| + |v|\sqrt{q_n})^2 = \varepsilon_n^2$, i.e., $k = 2$. In the $k = 1$ regime, writing $\frac{m - 1}{2} = \delta_1(n)z_1^2$ and $\frac{m + 1}{2} = \delta_2(n)z_2^2$ with $\delta_1(n), \delta_2(n)$ square-free (so that $\delta_1(n)\delta_2(n) = q_n$, $2z_1 z_2 = s_n$, and $\delta_1(n)z_1^2 + \delta_2(n)z_2^2 = m$), multiplying $1 + s_n\sqrt{-q_n}$ by $\delta_2(n)$ gives $\delta_2(n)(1 + s_n\sqrt{-q_n}) = (\delta_2(n)z_2 + z_1\sqrt{-q_n})^2$, so $(\mathfrak{d}_2\mathfrak{J}_n)^2 = (\delta_2(n)z_2 + z_1\sqrt{-q_n})^2 \implies [\mathfrak{J}_n] = [\mathfrak{d}_2] \ne 1$. Finally, in the $k = 2$ regime, if $[\mathfrak{I}_n] = 1$, then $X_1(q_n) + Y_1(q_n)\sqrt{q_n}$ would be the square of a smaller unit in $\mathcal{O}_{K_n}^\times$, contradicting the minimality of $\varepsilon_n$.
\end{proof}

\section{Arithmetic Structure of the Critical Regime $k = 2$}\label{sec:4}

Theorem~\ref{thm:1.1} bifurcates the study of \eqref{eq:1.1} into two arithmetic regimes: the critical-capacity regime $k = 2$, where $\log \varepsilon_n = (\frac{1}{4} + o(1))n\log n$ saturates the Chebyshev--Lucas capacity bound \eqref{eq:3.5}, and the fundamental regime $k = 1$, where $m = X_1(q_n)$ and $s_n = Y_1(q_n)$. In this section we analyze the critical regime $k = 2$, establishing the $(n/2)$-smoothness asymmetry of $m + 1$, the first- and second-stage Euclidean division quartic laws, and the stretched-exponential lower bounds on the midpoint cores $\delta_i(n) = \core(X_1(q_n) \mp 1)$.

When $k = 2$, expanding $m + s_n\sqrt{q_n} = \varepsilon_n^2 = (X_1(q_n) + Y_1(q_n)\sqrt{q_n})^2$ and substituting $X_1(q_n)^2 - q_n Y_1(q_n)^2 = 1$ yields
\begin{equation}\label{eq:4.1}
m = X_1(q_n)^2 + q_n Y_1(q_n)^2 = 2X_1(q_n)^2 - 1 = 2q_n Y_1(q_n)^2 + 1, \qquad s_n = 2X_1(q_n)Y_1(q_n).
\end{equation}
Consequently, the factors $m - 1$ and $m + 1$ in $(m - 1)(m + 1) = n!$ take the form
\begin{equation}\label{eq:4.2}
m + 1 = 2X_1(q_n)^2, \qquad m - 1 = 2q_n Y_1(q_n)^2 = 2\bigl(X_1(q_n) - 1\bigr)\bigl(X_1(q_n) + 1\bigr),
\end{equation}
and Brocard's equation \eqref{eq:1.1} becomes the three-consecutive-integer factorial equation
\begin{equation}\label{eq:4.3}
n! = 4\bigl(X_1(q_n) - 1\bigr)\,X_1(q_n)^2\,\bigl(X_1(q_n) + 1\bigr).
\end{equation}
We begin by recording the prime-support separation and biquadratic splitting laws forced by \eqref{eq:4.3}.

\begin{proposition}[{$(n/2)$-Smoothness of $m + 1$, Biquadratic Splitting, and Fermat-Quotient Lock for $k = 2$}]\label{prop:4.1}
Suppose $(n, m)$ is a solution to $n! + 1 = m^2$ in the $k = 2$ regime with $n \ge 8$. Then:
\begin{enumerate}
\item[(i)] Every odd prime $p \le n$ divides at most one of the three consecutive integers $X_1(q_n) - 1$, $X_1(q_n)$, $X_1(q_n) + 1$. Specifically, every odd prime $p \in \Podd(n)$ divides $(X_1(q_n) - 1)(X_1(q_n) + 1) = q_n Y_1(q_n)^2$ to its full factorial power $p^{v_p(n!)}$, whereas every odd prime $p \mid X_1(q_n)$ belongs to $\Peven(n)$ with $p^{v_p(n!)/2} \parallel X_1(q_n)$ and $p \le n/2$. In particular, $m + 1 = 2X_1(q_n)^2$ is strictly $(n/2)$-smooth, and as $n \to \infty$,
\begin{equation}\label{eq:4.4}
\log \rad(m - 1) \ge \sum_{p \in \Podd(n)} \log p = (\log 2 + o(1))\,n, \qquad \log \rad(m + 1) \le (1 - \log 2 + o(1))\,n.
\end{equation}
\item[(ii)] Write $X_1(q_n) - 1 = \delta_1(n)u_n^2$ and $X_1(q_n) + 1 = \delta_2(n)v_n^2$ with $\delta_1(n), \delta_2(n)$ square-free. Every odd prime $p \mid X_1(q_n)$ splits completely in the biquadratic field $L_n := \Q\bigl(\sqrt{-\delta_1(n)}, \sqrt{\delta_2(n)}\bigr)$:
\begin{equation}\label{eq:4.5}
\left(\frac{-\delta_1(n)}{p}\right) = +1, \qquad \left(\frac{\delta_2(n)}{p}\right) = +1, \qquad \left(\frac{-q_n}{p}\right) = +1 \qquad \forall\, p \mid X_1(q_n),\; p \ge 3,
\end{equation}
and satisfies the modulo-$p^{v_p(n!)}$ and Fermat-quotient ($q_p(a) := \frac{a^{p-1} - 1}{p}$) congruences
\begin{equation}\label{eq:4.6}
q_n Y_1(q_n)^2 \equiv -1 \pmod{p^{v_p(n!)}}, \qquad q_p(-q_n) + 2\,q_p\bigl(Y_1(q_n)\bigr) \equiv 0 \pmod p.
\end{equation}
\end{enumerate}
\end{proposition}

\begin{proof}
(i) Since $\gcd(X_1(q_n), X_1(q_n) \mp 1) = 1$ and $\gcd(X_1(q_n) - 1, X_1(q_n) + 1) \mid 2$, the odd prime factors of $X_1(q_n) - 1$, $X_1(q_n)$, and $X_1(q_n) + 1$ are pairwise disjoint. By \eqref{eq:4.3}, any odd prime $p \mid X_1(q_n)$ satisfies $v_p(n!) = 2v_p(X_1(q_n))$, so $p \in \Peven(n)$; by Lemma~\ref{lem:2.1}(i), every prime $p \in I_1 = (n/2, n]$ has $v_p(n!) = 1$ and thus lies in $\Podd(n)$, so $X_1(q_n)$ is $(n/2)$-smooth. Summing $\log p$ over $p \in \Podd(n)$ via Lemma~\ref{lem:2.1}(ii) gives \eqref{eq:4.4}.

(ii) Let $p \ge 3$ be a prime dividing $X_1(q_n)$. Since $p \in \Peven(n)$ and $p \le n/2$, we have $v_p(n!) \ge 2$ and $p^{v_p(n!)/2} \parallel X_1(q_n)$, so $X_1(q_n)^2 \equiv 0 \pmod{p^{v_p(n!)}}$. Reducing $q_n Y_1(q_n)^2 = X_1(q_n)^2 - 1$ modulo $p^{v_p(n!)}$ gives $q_n Y_1(q_n)^2 \equiv -1 \pmod{p^{v_p(n!)}}$, while reducing $\delta_1(n)u_n^2 = X_1(q_n) - 1$ and $\delta_2(n)v_n^2 = X_1(q_n) + 1$ modulo $p$ gives $-\delta_1(n)u_n^2 \equiv +1 \pmod p$ and $\delta_2(n)v_n^2 \equiv +1 \pmod p$, proving \eqref{eq:4.5}. Since $v_p(n!) \ge 2$ and $\bigl(\frac{-q_n}{p}\bigr) = +1$, raising $(-q_n)Y_1(q_n)^2 \equiv 1 \pmod{p^2}$ to the $(p - 1)$-st power gives $q_p(-q_n) + 2\,q_p(Y_1(q_n)) \equiv 0 \pmod p$, proving \eqref{eq:4.6}.
\end{proof}

Proposition~\ref{prop:4.1}(i) shows that in the $k = 2$ regime, the odd prime factors of $X_1(q_n)$ are drawn exclusively from $\Peven(n)$, while every prime in $\Podd(n)$ is confined to $q_n Y_1(q_n)^2 = X_1(q_n)^2 - 1$. To turn this prime-support separation into an exact Diophantine equation, define the explicit even-valuation and odd-valuation factorial integers
\begin{equation}\label{eq:4.7}
A_n^* := \prod_{\substack{3 \le p \le n \\ p \in \Peven(n)}} p^{v_p(n!)}, \qquad B_n^* := \prod_{\substack{3 \le p \le n \\ p \in \Podd(n)}} p^{v_p(n!)}, \qquad Z := \prod_{\substack{3 \le p \le n \\ p \in \Peven(n),\, p \mid Y_1(q_n)}} p^{v_p(n!)/2},
\end{equation}
and set $(A_n, B_n) := (A_n^*, 2^{v_2(n!)-2} B_n^*)$ if $X_1(q_n)$ is odd, and $(A_n, B_n) := (2^{v_2(n!)-2} A_n^*, B_n^*)$ if $X_1(q_n)$ is even (which requires $2 \mid v_2(n!)$ and $q_n \equiv 7 \pmod 8$). The following theorem establishes the resulting two-stage Euclidean division quartic law and stretched-exponential lower bounds on the midpoint cores $\delta_1(n)$ and $\delta_2(n)$.

\begin{theorem}[{Two-Stage Euclidean Division Quartic Law and Stretched-Exponential Midpoint Core Bounds}]\label{thm:4.2}
Suppose $(n, m)$ is a solution to $n! + 1 = m^2$ in the $k = 2$ regime with $n \ge 8$. Then:
\begin{enumerate}
\item[(i)] \textbf{First-Stage Euclidean Division Quartic Law:} The integers $A_n, B_n, Z$ satisfy $\gcd(A_n, B_n) = 1$, $X_1(q_n) = \frac{\sqrt{A_n}}{Z}$, $q_n Y_1(q_n)^2 = B_n Z^2$, and the exact quartic Diophantine identity
\begin{equation}\label{eq:4.8}
A_n = B_n Z^4 + Z^2, \qquad 4\log Z = \bigl(c_2(n) - \Sparitys(n)\bigr)n + O\!\left(\frac{n}{\log n}\right) = o(n\log n),
\end{equation}
where $c_2(n) = -1$ if $X_1(q_n)$ is odd and $c_2(n) = 2\log 2 - 1 \approx 0.38629$ if $X_1(q_n)$ is even. For all sufficiently large $n$, $0 < Z^2 < B_n$, so Euclidean division $A_n = Q_n B_n + R_n$ ($0 \le R_n < B_n$) satisfies the simultaneous quotient-remainder lock
\begin{equation}\label{eq:4.9}
Q_n = Z^4, \qquad R_n = Z^2, \qquad R_n^2 = Q_n.
\end{equation}
Moreover, splitting $Z = Z_{\le \sqrt{n}} Z_{> \sqrt{n}}$ into its $(\sqrt{n})$-smooth and $(\sqrt{n})$-rough parts flips the sign of every prime $p \le \sqrt{n}$ dividing $Z$:
\begin{equation}\label{eq:4.10}
4\log Z_{> \sqrt{n}} = \left(c_2(n) + \sum_{3 \le p \le \sqrt{n}} \varepsilon_X(p)\frac{\log p}{p - 1}\right)n + O\!\left(\frac{n}{\log n}\right),
\end{equation}
where $\varepsilon_X(p) = +1$ if $p \mid X_1(q_n)$ and $\varepsilon_X(p) = -1$ if $p \nmid X_1(q_n)$.
\item[(ii)] \textbf{Second-Stage Consecutive Splitting and Dual Quartic Locks:} Let $c_0 := \gcd(X_1(q_n) - 1, X_1(q_n) + 1) \in \{1, 2\}$. There exist unique coprime factorizations $B_- B_+ = B_n / c_0^2$ and $Z_- Z_+ = Z$ such that
\begin{equation}\label{eq:4.11}
\frac{X_1(q_n) - 1}{c_0} = B_- Z_-^2, \qquad \frac{X_1(q_n) + 1}{c_0} = B_+ Z_+^2, \qquad B_+ Z_+^2 - B_- Z_-^2 = \frac{2}{c_0},
\end{equation}
and for all sufficiently large $n$, the first-stage remainder $R_n = Z^2 = Z_-^2 Z_+^2$ satisfies the dual quartic equations
\begin{equation}\label{eq:4.12}
B_+ R_n = B_- Z_-^4 + \frac{2}{c_0}Z_-^2, \qquad B_- R_n = B_+ Z_+^4 - \frac{2}{c_0}Z_+^2,
\end{equation}
with unconditional Euclidean quotient-remainder locks:
\begin{equation}\label{eq:4.13}
\left\lfloor \frac{B_+ R_n}{B_-} \right\rfloor = Z_-^4, \qquad \left\lceil \frac{B_- R_n}{B_+} \right\rceil = Z_+^4.
\end{equation}
\item[(iii)] \textbf{Stretched-Exponential Lower Bounds on the Midpoint Cores $\delta_1(n), \delta_2(n)$:} Let $\delta_1(n) := \core(X_1(q_n) - 1)$ and $\delta_2(n) := \core(X_1(q_n) + 1)$. For both $i \in \{1, 2\}$, the explicit Mertens bound gives $\omega(\delta_i(n)) \ge (8e^\gamma - o(1))\frac{n^{1/4}}{\log n}$ and $\delta_i(n) \ge \exp((2e^\gamma - o(1))n^{1/4})$, and applying the power-band equidistribution \eqref{eq:2.18} upgrades these bounds unconditionally for both $i \in \{1, 2\}$ to
\begin{equation}\label{eq:4.14}
P^+\bigl(\delta_i(n)\bigr) \ge n^{1/2-o(1)}, \qquad \omega\bigl(\delta_i(n)\bigr) \ge n^{1/2-o(1)}, \qquad \delta_i(n) \ge \exp\bigl(n^{1/2-o(1)}\bigr).
\end{equation}
\end{enumerate}
\end{theorem}

\begin{proof}
\textbf{Part (i):} Since $\Podd(n)$ and $\Peven(n)$ are disjoint sets of odd primes whose union is the set of all odd primes $3 \le p \le n$, we have $\gcd(A_n^*, B_n^*) = 1$ and $A_n^* B_n^* = \frac{n!}{2^{v_2(n!)}}$. By Proposition~\ref{prop:4.1}(i), every odd prime $p \in \Podd(n)$ satisfies $p^{v_p(n!)} \parallel q_n Y_1(q_n)^2$, while every odd prime $p \in \Peven(n)$ divides $s_n = 2X_1(q_n)Y_1(q_n)$ with $\gcd(X_1(q_n), Y_1(q_n)) = 1$, so $p^{v_p(n!)/2} \parallel Y_1(q_n)$ if $p \mid Y_1(q_n)$ (giving $Z = \prod_{p \in \Peven(n),\, p \mid Y_1} p^{v_p(n!)/2}$) and $p^{v_p(n!)/2} \parallel X_1(q_n)$ if $p \mid X_1(q_n)$ (giving the odd part of $X_1(q_n)$ equal to $\frac{\sqrt{A_n^*}}{Z}$).

For the prime $p = 2$, since $4X_1(q_n)^2(X_1(q_n)^2 - 1) = n!$ and $\gcd(X_1(q_n)^2, X_1(q_n)^2 - 1) = 1$, there are two mutually exclusive cases:
\begin{itemize}
\item \emph{Type I ($X_1(q_n)$ odd):} Here $2 \nmid X_1(q_n)^2$, so $2^{v_2(n!)-2} \parallel (X_1(q_n)^2 - 1) = q_n Y_1(q_n)^2$. Setting $A_n := A_n^*$ and $B_n := 2^{v_2(n!)-2} B_n^*$ gives $X_1(q_n) = \frac{\sqrt{A_n}}{Z}$ and $q_n Y_1(q_n)^2 = B_n Z^2$.
\item \emph{Type II ($X_1(q_n)$ even):} Here $2 \nmid (X_1(q_n)^2 - 1)$, so $2^{v_2(n!)-2} \parallel X_1(q_n)^2$, which requires $2 \mid v_2(n!)$ and $v_2(X_1(q_n)) = \frac{v_2(n!)-2}{2} \ge 3$ for $n \ge 8$; reducing $q_n Y_1(q_n)^2 = X_1(q_n)^2 - 1 \equiv -1 \pmod 8$ forces $q_n \equiv 7 \pmod 8$ and $2 \nmid Y_1(q_n)$. Setting $A_n := 2^{v_2(n!)-2} A_n^*$ and $B_n := B_n^*$ again gives $X_1(q_n) = \frac{\sqrt{A_n}}{Z}$ and $q_n Y_1(q_n)^2 = B_n Z^2$.
\end{itemize}
Substituting $X_1(q_n) = \frac{\sqrt{A_n}}{Z}$ and $q_n Y_1(q_n)^2 = B_n Z^2$ into the fundamental Pell equation $X_1(q_n)^2 - q_n Y_1(q_n)^2 = 1$ gives $\frac{A_n}{Z^2} - B_n Z^2 = 1$, and multiplying by $Z^2$ yields the exact quartic identity
\[
A_n = B_n Z^4 + Z^2.
\]
Taking logarithms of $Z^4 = \frac{A_n}{B_n}\bigl(1 + \frac{1}{B_n Z^2}\bigr)^{-1}$ and using $\log A_n^* - \log B_n^* = \Seven(n) - \Sodd(n) = \Dparity(n)$ together with $v_2(n!)\log 2 = (\log 2)n + O(\log n)$ and Theorem~\ref{thm:2.5}(i)--(ii) yields
\begin{align*}
4\log Z &= \Dparity(n) \mp \bigl(v_2(n!) - 2\bigr)\log 2 - \log\!\left(1 + \frac{1}{B_n Z^2}\right) \\
&= \bigl(c_2(n) - \Sparitys(n)\bigr)n + O\!\left(\frac{n}{\log n}\right) = o(n\log n),
\end{align*}
where $c_2(n) = (\log 2 - 1) - \log 2 = -1$ when $X_1(q_n)$ is odd and $c_2(n) = (\log 2 - 1) + \log 2 = 2\log 2 - 1$ when $X_1(q_n)$ is even. Since $\log B_n \ge \Sodd(n) = (\frac{1}{2} - o(1))n\log n$ while $\log(Z^2) = o(n\log n)$, we have $0 < Z^2 < B_n$ for all sufficiently large $n$. By uniqueness of Euclidean division $A_n = Q_n B_n + R_n$ ($0 \le R_n < B_n$), we obtain $Q_n = Z^4$, $R_n = Z^2$, and $R_n^2 = Q_n$, proving \eqref{eq:4.9}. Finally, writing $Z = Z_{\le \sqrt{n}} Z_{> \sqrt{n}}$, every prime $p \le \sqrt{n}$ in $\Peven(n)$ dividing $Z$ contributes $2v_p(n!)\log p = \frac{2\log p}{p-1}n + O(\log n)$ to $4\log Z_{\le \sqrt{n}}$ and $+\frac{\log p}{p-1}n$ to $-\Sparitys(n)n$, so subtracting $4\log Z_{\le \sqrt{n}}$ from both sides flips the sign of $p \mid Z$ from $+\frac{\log p}{p-1}n$ to $-\frac{\log p}{p-1}n = \varepsilon_X(p)\frac{\log p}{p-1}n$, establishing \eqref{eq:4.10}.

\textbf{Part (ii):} Factor $X_1(q_n)^2 - 1 = (X_1(q_n) - 1)(X_1(q_n) + 1) = B_n Z^2$. Setting $c_0 := \gcd(X_1(q_n) - 1, X_1(q_n) + 1) \in \{1, 2\}$, the positive integers $\frac{X_1(q_n)-1}{c_0}$ and $\frac{X_1(q_n)+1}{c_0}$ are coprime and have product $\frac{B_n}{c_0^2}Z^2$. Since $\gcd(B_n, Z) = 1$, each prime power dividing $\frac{B_n}{c_0^2}$ or $Z^2$ divides either $\frac{X_1(q_n)-1}{c_0}$ or $\frac{X_1(q_n)+1}{c_0}$ to its full valuation, yielding unique coprime factorizations $B_- B_+ = B_n / c_0^2$ and $Z_- Z_+ = Z$ satisfying \eqref{eq:4.11}. Multiplying $B_+ Z_+^2 - B_- Z_-^2 = \frac{2}{c_0}$ by $Z_-^2$ and by $Z_+^2$ and substituting $Z_-^2 Z_+^2 = Z^2 = R_n$ gives \eqref{eq:4.12}. Since $\log(Z_\mp^4) \le 4\log Z = o(n\log n)$ and $\log(B_\mp Z_\mp^2) = \log(\frac{X_1(q_n)\mp 1}{c_0}) = (\frac{1}{4} + o(1))n\log n$, we have
\begin{equation}\label{eq:4.15}
\log B_- = \left(\frac{1}{4} + o(1)\right)n\log n, \qquad \log B_+ = \left(\frac{1}{4} + o(1)\right)n\log n.
\end{equation}
Hence $0 < \frac{2}{c_0}Z_\mp^2 < B_\mp$ for all sufficiently large $n$, and Euclidean division in \eqref{eq:4.12} yields \eqref{eq:4.13}.

\textbf{Part (iii):} The odd parts of the square-free cores $\delta_1(n) = \core(X_1(q_n) - 1)$ and $\delta_2(n) = \core(X_1(q_n) + 1)$ are $\delta_1^*(n) = \prod_{p \in \Podd(n),\, p \mid B_-} p$ and $\delta_2^*(n) = \prod_{p \in \Podd(n),\, p \mid B_+} p$. Fix $i \in \{1, 2\}$ and let $t := \omega(\delta_i^*(n))$. Since $v_2(B_\mp)\log 2 \le v_2(n!)\log 2 = O(n)$, equation \eqref{eq:4.15} gives
\begin{equation}\label{eq:4.16}
\sum_{p \mid \delta_i^*(n)} v_p(n!)\log p = \left(\frac{1}{4} + o(1)\right)n\log n.
\end{equation}
Since $v_p(n!) \le \frac{n}{p-1}$ and $p \mapsto \frac{\log p}{p-1}$ is strictly decreasing for $p \ge 3$, the left-hand side of \eqref{eq:4.16} is bounded above by the sum over the first $t$ odd primes $3 = p_2 < p_3 < \dots < p_{t+1}$. By Mertens' first theorem \cite[Chapter~7]{Davenport2000}, $\sum_{j=2}^{t+1} \frac{\log p_j}{p_j-1} = \log p_{t+1} - \gamma - \log 2 + o(1)$, so comparing with \eqref{eq:4.16} gives $\log p_{t+1} \ge \frac{1}{4}\log n + \gamma + \log 2 - o(1)$, i.e., $p_{t+1} \ge (2e^\gamma - o(1))n^{1/4}$, $\omega(\delta_i(n)) \ge (8e^\gamma - o(1))\frac{n^{1/4}}{\log n}$, and $\delta_i(n) \ge \exp(\vartheta(p_{t+1}) - \log 2) \ge \exp((2e^\gamma - o(1))n^{1/4})$.

To upgrade the exponent from $1/4$ to $1/2 - o(1)$, fix any $\varepsilon \in (0, 1/6)$. Since every prime $p \mid \delta_i^*(n)$ belongs to $\Podd(n)$, Theorem~\ref{thm:2.5}(iii) (equation \eqref{eq:2.18} with $h = 2$, $c = 1$, $a = 0$, $b = \frac{1}{2} - \varepsilon$) gives
\[
\sum_{\substack{p \mid \delta_i^*(n) \\ p \le n^{1/2-\varepsilon}}} v_p(n!)\log p \le \sum_{\substack{3 \le p \le n^{1/2-\varepsilon} \\ p \in \Podd(n)}} v_p(n!)\log p = \left(\frac{1}{4} - \frac{\varepsilon}{2} + o(1)\right)n\log n.
\]
Subtracting this from \eqref{eq:4.16} shows that the prime divisors of $\delta_i^*(n)$ exceeding $n^{1/2-\varepsilon}$ carry factorial weight
\[
\sum_{\substack{p \mid \delta_i^*(n) \\ p > n^{1/2-\varepsilon}}} v_p(n!)\log p \ge \left(\frac{\varepsilon}{2} - o(1)\right)n\log n.
\]
Since every prime $p > n^{1/2-\varepsilon}$ satisfies $v_p(n!) \le \frac{n}{p-1} \le 2n^{1/2+\varepsilon}$ and $\log p \le \log n$, dividing by $2n^{1/2+\varepsilon}\log n$ yields $P^+(\delta_i(n)) \ge n^{1/2-\varepsilon}$, $\omega(\delta_i(n)) \ge (\frac{\varepsilon}{4} - o(1))n^{1/2-\varepsilon}$, and $\log \delta_i(n) \ge (\frac{\varepsilon}{4} - o(1))n^{1/2-\varepsilon}\log n$. Since $\varepsilon > 0$ is arbitrary, this establishes \eqref{eq:4.14}.
\end{proof}

\section{The Fundamental Regime $k = 1$ and R\'edei--Reichardt Cohomology}\label{sec:5}

We now turn to the fundamental regime $k = 1$ (where $m = X_1(q_n)$ and $s_n = Y_1(q_n)$), together with the intrinsic continued-fraction midpoint structure of $\sqrt{q_n}$ valid for all $n \ge 8$. In this section we establish the continued-fraction midpoint bias, the reciprocal cross-symmetry and effective Landau--Siegel sub-exponential core bounds, the Galois cohomology / R\'edei--Reichardt 4-rank exact sequence, and the reduced dual quartic laws.

\subsection{Midpoint Core Bipartition, Cross-Symmetry, and Effective Sub-Exponential Core Bounds}
We first establish the properties of the continued-fraction midpoint of $\sqrt{q_n}$ and the one-sided character bias on $I_1 = (n/2, n]$.

\begin{proposition}[{Continued-Fraction Midpoint Denominator and Square-Root Conductor Bias}]\label{prop:5.1}
Let $n \ge 8$ be an integer, let $\varepsilon_n = X_1(q_n) + Y_1(q_n)\sqrt{q_n}$ be the fundamental unit of $\Z[\sqrt{q_n}]$, and write
\begin{equation}\label{eq:5.1}
X_1(q_n) - 1 = \delta_1^{\mathrm{mid}}(n)\,u_n^2, \qquad X_1(q_n) + 1 = \delta_2^{\mathrm{mid}}(n)\,v_n^2, \qquad \delta_2^{\mathrm{mid}}(n)\,v_n^2 - \delta_1^{\mathrm{mid}}(n)\,u_n^2 = 2,
\end{equation}
where $\delta_i^{\mathrm{mid}}(n) := \core(X_1(q_n) \mp 1)$ have odd parts $\delta_i^{\mathrm{mid}*}(n) = \delta_i^{\mathrm{mid}}(n)/\gcd(\delta_i^{\mathrm{mid}}(n), 2)$. Then:
\begin{enumerate}
\item[(i)] $\delta_1^{\mathrm{mid}*}(n)\,\delta_2^{\mathrm{mid}*}(n) = q_n^*$ is a coprime factorization of the odd part $q_n^* > 1$ of the square-free integer $q_n$ into two distinct factors, immediately forcing
\begin{equation}\label{eq:5.2}
\delta_{\min}(n) := \min\bigl(\delta_1^{\mathrm{mid}*}(n),\, \delta_2^{\mathrm{mid}*}(n)\bigr) < \sqrt{q_n^*} \le \sqrt{q_n}.
\end{equation}
Moreover, the period length $L(q_n)$ of the regular continued fraction of $\sqrt{q_n}$ is even, and at the half-period $M = L(q_n)/2$, the complete quotient denominator $Q_M^* < 2\sqrt{q_n}$ satisfies $Q_M^* = 2^a \delta_{i_n}^{\mathrm{mid}}(n)$ for some $a \in \{-1, 0, 1\}$ and $i_n \in \{1, 2\}$.
\item[(ii)] Let $\psi_1(p) := \bigl(\frac{-2\delta_1^{\mathrm{mid}}(n)}{p}\bigr)$ and $\psi_2(p) := \bigl(\frac{2\delta_2^{\mathrm{mid}}(n)}{p}\bigr)$, and let $\psi_{\min} \in \{\psi_1, \psi_2\}$ be the character corresponding to $\delta_{\min}(n)$, of conductor $Q_{\min} \le 8\delta_{\min}(n) < 8\sqrt{q_n}$. On $I_1 = (n/2, n]$, $\psi_{\min}$ satisfies the unconditional one-sided bias
\begin{equation}\label{eq:5.3}
\sum_{p \in (n/2, n]} \psi_{\min}(p)\log p = \vartheta(n) - \vartheta(n/2) - \sum_{\substack{p \in (n/2, n] \\ p \mid \delta_{\min}(n)}} \log p \ge \left(\frac{1 - \log 2}{2} - o(1)\right)n \approx 0.153426\,n.
\end{equation}
\end{enumerate}
\end{proposition}

\begin{proof}
(i) Since $\delta_1^{\mathrm{mid}}(n)\delta_2^{\mathrm{mid}}(n)(u_n v_n)^2 = X_1(q_n)^2 - 1 = q_n Y_1(q_n)^2$ and $\gcd(X_1(q_n) - 1, X_1(q_n) + 1) \mid 2$, taking $p$-adic valuations modulo 2 at every odd prime $p$ shows that $\delta_1^{\mathrm{mid}*}(n)\delta_2^{\mathrm{mid}*}(n) = q_n^*$ with $\gcd(\delta_1^{\mathrm{mid}*}(n), \delta_2^{\mathrm{mid}*}(n)) = 1$. Because $q_n^* > 1$ is square-free, $\delta_1^{\mathrm{mid}*}(n) \ne \delta_2^{\mathrm{mid}*}(n)$, so $\min(\delta_1^{\mathrm{mid}*}(n), \delta_2^{\mathrm{mid}*}(n)) < \sqrt{q_n^*} \le \sqrt{q_n}$, proving \eqref{eq:5.2}. Moreover, by Lemma~\ref{lem:3.1}(i), $x^2 - q_n y^2 = -1$ has no solutions, so $L(q_n)$ is even, and by Lagrange's midpoint symmetry theorem for regular continued fractions \cite[\S 25]{Perron1954}, the complete quotient denominator $Q_{L(q_n)/2}^*$ at the half-period divides $2q_n$, is square-free up to a factor of 2, satisfies $0 < Q_{L(q_n)/2}^* < 2\sqrt{q_n}$, and equals $2^a \delta_{i_n}^{\mathrm{mid}}(n)$ ($a \in \{-1, 0, 1\}$, $i_n \in \{1, 2\}$).

(ii) Every prime $p \in I_1 = (n/2, n]$ has $v_p(n!) = 1$ and therefore divides $q_n^* = \delta_1^{\mathrm{mid}*}(n)\delta_2^{\mathrm{mid}*}(n)$. Reducing $\delta_2^{\mathrm{mid}}(n)v_n^2 - \delta_1^{\mathrm{mid}}(n)u_n^2 = 2$ modulo $p \in I_1$ shows that if $p \mid \delta_1^{\mathrm{mid}}(n)$, then $\psi_1(p) = 0$ and $\psi_2(p) = \bigl(\frac{2\delta_2^{\mathrm{mid}}(n)}{p}\bigr) = \bigl(\frac{(2\delta_2^{\mathrm{mid}}(n)v_n)^2}{p}\bigr) = +1$; symmetrically, if $p \mid \delta_2^{\mathrm{mid}}(n)$, then $\psi_2(p) = 0$ and $\psi_1(p) = +1$. Thus $\psi_{\min}(p) = +1$ for all $p \in I_1$ with $p \nmid \delta_{\min}(n)$ and $\psi_{\min}(p) = 0$ for $p \in I_1$ with $p \mid \delta_{\min}(n)$. Since $\sum_{p \in I_1,\, p \mid \delta_{\min}(n)} \log p \le \log \delta_{\min}(n) < \frac{1}{2}\log q_n = (\frac{\log 2}{2} + o(1))n$ and $\vartheta(n) - \vartheta(n/2) = (\frac{1}{2} - o(1))n$, subtracting yields \eqref{eq:5.3}.
\end{proof}

Combining Proposition~\ref{prop:5.1} with Theorem~\ref{thm:2.5} and the explicit formula for $\psi(x, \psi_i)$ yields the reciprocal cross-symmetry laws, radical lower bounds, and effective Landau--Siegel sub-exponential lower bounds on $\delta_1^{\mathrm{mid}}(n)$ and $\delta_2^{\mathrm{mid}}(n)$.

\begin{theorem}[{Midpoint Core Bipartition, Cross-Symmetry, Radical Bounds, and Effective Landau--Siegel Lower Bounds}]\label{thm:5.2}
Let $n \ge 4$ be an integer, and for any solution $(n, m)$ to $n! + 1 = m^2$, write $u_1 := \frac{m-1}{2} = \delta_1(n)z_1^2$ and $u_2 := \frac{m+1}{2} = \delta_2(n)z_2^2$ with $\delta_1(n), \delta_2(n) \ge 1$ square-free. Then:
\begin{enumerate}
\item[(i)] \textbf{Midpoint Core Bipartition, Cross-Symmetry, and Radical Bounds:} We have $\delta_1(n)\delta_2(n) = q_n$, $\gcd(\delta_1(n), \delta_2(n)) = 1$, $z_1 z_2 = s_n/2$, $\delta_2(n)z_2^2 - \delta_1(n)z_1^2 = 1$, and $\delta_1(n) > 1$ for all $n \ge 3$ (with $\delta_2(n) = 1 \iff k = 2$ and $\delta_2(n) > 1 \iff k = 1$ whenever $k \in \{1, 2\}$, in particular for all sufficiently large $n$). Decomposing the odd parts $u_i^* = B^{(i)}(Z_i^*)^2$ into $B^{(i)} := \prod_{p \in \Podd(n),\, p \mid u_i^*} p^{v_p(n!)}$ and $Z_i^* := \prod_{q \in \Peven(n),\, q \mid u_i^*} q^{v_q(n!)/2}$ yields the exact reciprocal cross-symmetry laws
\begin{equation}\label{eq:5.4}
\begin{aligned}
\log B^{(2)} - \log\bigl((Z_1^*)^2\bigr) &= \frac{1}{2}\Bigl(\bigl(w_2^{(1)} - w_2^{(2)}\bigr)\log 2 - \Dparity(n)\Bigr) + O\bigl(e^{-\frac{1}{3}n\log n}\bigr) = O(n), \\
\log B^{(1)} - \log\bigl((Z_2^*)^2\bigr) &= \frac{1}{2}\Bigl(\bigl(w_2^{(2)} - w_2^{(1)}\bigr)\log 2 - \Dparity(n)\Bigr) + O\bigl(e^{-\frac{1}{3}n\log n}\bigr) = O(n),
\end{aligned}
\end{equation}
where $w_2^{(i)} = v_2(u_i)$, together with the unconditional radical bounds for both $m - 1$ and $m + 1$:
\begin{equation}\label{eq:5.5}
P^+(m \mp 1) \ge n^{1/2-o(1)}, \qquad \omega(m \mp 1) \ge n^{1/2-o(1)}, \qquad \rad(m \mp 1) \ge \exp\bigl(n^{1/2-o(1)}\bigr).
\end{equation}
\item[(ii)] \textbf{Effective Sub-Exponential Midpoint Core Bounds and Non-Solubility of $x^2 - q_n y^2 = \pm 2$:} For every integer $n \ge 8$ (regardless of whether $n! + 1 = m^2$ is soluble), there exists an absolute effective constant $c_0 > 0$ such that for all sufficiently large $n$,
\begin{equation}\label{eq:5.6}
\delta_1^{\mathrm{mid}}(n) > \exp\bigl(c_0\sqrt{\log n}\bigr), \qquad \delta_2^{\mathrm{mid}}(n) > \exp\bigl(c_0\sqrt{\log n}\bigr).
\end{equation}
Consequently, the Pell-type equations $x^2 - q_n y^2 = 2$ and $x^2 - q_n y^2 = -2$ have no integer solutions for all sufficiently large $n$.
\end{enumerate}
\end{theorem}

\begin{proof}
(i) For any solution $(n, m)$ to $n! + 1 = m^2$ with $n \ge 4$, $u_1 = \frac{m-1}{2}$ and $u_2 = \frac{m+1}{2}$ are coprime positive integers with $u_2 - u_1 = 1$ and $u_1 u_2 = \frac{n!}{4} = q_n(s_n/2)^2$. Writing $u_1 = \delta_1(n)z_1^2$ and $u_2 = \delta_2(n)z_2^2$ with $\delta_1(n), \delta_2(n) \ge 1$ square-free gives $\delta_1(n)\delta_2(n) = q_n$, $\gcd(\delta_1(n), \delta_2(n)) = 1$, $z_1 z_2 = s_n/2$, and $\delta_2(n)z_2^2 - \delta_1(n)z_1^2 = 1$. If $\delta_1(n) = 1$, then $z_1^2 - q_n z_2^2 = -1$, which is impossible by Lemma~\ref{lem:3.1}(i). Since $(z_2\sqrt{\delta_2(n)} + z_1\sqrt{\delta_1(n)})^2 = m + s_n\sqrt{q_n} = \varepsilon_n^k$ with $k \in \{1, 2\}$ by Theorem~\ref{thm:1.1}, we have $\delta_2(n) = 1 \iff k = 2$ and $\delta_2(n) > 1 \iff k = 1$.

For \eqref{eq:5.4}, note that $u_1 u_2 = n!/4$ gives $v_2(n!) = 2 + w_2^{(1)} + w_2^{(2)}$ and $\log B^{(1)} + \log((Z_1^*)^2) = \log u_1^* = \frac{1}{2}\log(n!) - (w_2^{(1)} + 1)\log 2 + O((n!)^{-1/2})$, while $\log B^{(1)} + \log B^{(2)} = \Sodd(n) = \frac{1}{2}\log(n!) - \frac{1}{2}\Dparity(n) - \frac{1}{2}v_2(n!)\log 2$. Subtracting these two identities eliminates $\log B^{(1)}$ and $\frac{1}{2}\log(n!)$ identically and uses $2w_2^{(1)} + 2 - v_2(n!) = w_2^{(1)} - w_2^{(2)}$, giving \eqref{eq:5.4}. For \eqref{eq:5.5}, let $S_\mp$ be the set of odd prime divisors of $m \mp 1$. Since $\sum_{p \in S_\mp} v_p(n!)\log p = (\frac{1}{2} + o(1))n\log n$ and Theorem~\ref{thm:2.5}(iii) bounds $\sum_{3 \le p \le n^{1/2-\varepsilon}} v_p(n!)\log p = (\frac{1}{2} - \varepsilon + o(1))n\log n$ for any $\varepsilon \in (0, 1/4)$, the primes $p \in S_\mp$ with $p > n^{1/2-\varepsilon}$ carry weight $\ge (\varepsilon - o(1))n\log n$ with $v_p(n!) \le 2n^{1/2+\varepsilon}$, forcing $P^+(m \mp 1) \ge n^{1/2-\varepsilon}$, $\omega(m \mp 1) \ge (\frac{\varepsilon}{2} - o(1))n^{1/2-\varepsilon}$, and $\log\rad(m \mp 1) \ge (\frac{\varepsilon}{2} - o(1))n^{1/2-\varepsilon}\log n$.

(ii) By Proposition~\ref{prop:5.1}(i)--(ii), for every $n \ge 7$ and $i \in \{1, 2\}$, the character $\psi_i(p)$ of conductor $Q_i \le 8\delta_i^{\mathrm{mid}}(n)$ is non-principal: $\psi_1(p) = \bigl(\frac{-2\delta_1^{\mathrm{mid}}(n)}{p}\bigr)$ is imaginary quadratic ($-2\delta_1^{\mathrm{mid}}(n) < 0$) and hence never principal, while $\psi_2(p) = \bigl(\frac{2\delta_2^{\mathrm{mid}}(n)}{p}\bigr)$ is principal if and only if $\delta_2^{\mathrm{mid}}(n) = 2$, which would give $X_1(q_n) + 1 = 2v_n^2$ and $X_1(q_n) - 1 = 2q_n^*(u_n')^2$ and hence $v_n^2 - q_n(u_n')^2 = 1$ with $0 < v_n < X_1(q_n)$, contradicting the minimality of the fundamental unit $\varepsilon_n$. Moreover, $\psi_i$ satisfies
\begin{equation}\label{eq:5.7}
\sum_{n/2 < p \le n} \psi_i(p)\log p = \vartheta(n) - \vartheta(n/2) - \sum_{\substack{p \in (n/2, n] \\ p \mid \delta_i^{\mathrm{mid}}(n)}} \log p \ge \frac{1}{2}n - O\!\left(ne^{-c_1\sqrt{\log n}}\right) - \log \delta_i^{\mathrm{mid}}(n).
\end{equation}
On the other hand, by the explicit formula for $\psi(x, \psi_i)$ and the de la Vall\'ee Poussin zero-free region \cite[Chapter~19]{Davenport2000}, if $Q_i \le \exp(2c_0\sqrt{\log n})$ for a small absolute constant $c_0 > 0$, then $L(s, \psi_i)$ has at most one simple real Landau--Siegel zero $\beta_1 \in (1 - \frac{c_2}{\log Q_i}, 1)$, and
\[
\sum_{n/2 < p \le n} \psi_i(p)\log p = -\frac{n^{\beta_1} - (n/2)^{\beta_1}}{\beta_1} + O\!\left(n\exp\bigl(-c_3\sqrt{\log n}\bigr)\right).
\]
Because $\beta_1 > 0 \implies n^{\beta_1} > (n/2)^{\beta_1}$, the Landau--Siegel zero contribution $-\frac{n^{\beta_1} - (n/2)^{\beta_1}}{\beta_1} < 0$ is strictly negative. Thus $\sum_{n/2 < p \le n} \psi_i(p)\log p \le O(n\exp(-c_3\sqrt{\log n})) = o(n)$ effectively, which contradicts \eqref{eq:5.7} unless $\delta_i^{\mathrm{mid}}(n) > \exp(c_0\sqrt{\log n})$. In particular, if $x^2 - q_n y^2 = \pm 2$ had an integer solution $(x, y)$, then $(x^2 \mp 1)^2 - q_n(xy)^2 = 1$ would be a unit in $\Z[\sqrt{q_n}]$ of the form $\varepsilon_n^{2r+1}$, forcing $\min(\delta_1^{\mathrm{mid}}(n), \delta_2^{\mathrm{mid}}(n)) \in \{1, 2\}$, contradicting \eqref{eq:5.6}.
\end{proof}

\subsection{R\'edei--Reichardt 4-Rank Cohomology and Reduced Dual Quartic Laws}
Let $\mathcal{L}_0(n) := \Podd(n) \setminus \{2\}$, and divide $\Peven(n) \setminus \{2\}$ into the split-even primes $\Peven^+(n) := \{q \in \Peven(n) \setminus \{2\} : \bigl(\frac{-q_n}{q}\bigr) = +1\}$ and the inert-even primes $\Peven^-(n) := \{q \in \Peven(n) \setminus \{2\} : \bigl(\frac{-q_n}{q}\bigr) = -1\}$. For any odd prime $p$ and integer $a$ with $p \nmid a$, define the additive Legendre index $\ell_p(a) \in \F_2$ by $(-1)^{\ell_p(a)} := \bigl(\frac{a}{p}\bigr)$. For any binary vector $y = (y_p)_{p \in \mathcal{L}_0(n)} \in \F_2^{\mathcal{L}_0(n)}$ encoding a coprime factorization $q_n = c_1(y)c_2(y)$ of $q_n$ into $c_2(y) = \prod_{p \in \mathcal{L}_0(n)} p^{y_p}$ (augmented at $p = 2$ when $2 \mid D_n$), reducing $c_2(y)z_2^2 - c_1(y)z_1^2 = 1$ modulo each $p \in \mathcal{L}_0(n)$ forces $\bigl(\frac{c_2(y)}{p}\bigr) = +1$ when $y_p = 0$ and $\bigl(\frac{-c_1(y)}{p}\bigr) = +1$ when $y_p = 1$. Since $c_1(y) = (q_n/p)/(c_2(y)/p)$ when $y_p = 1$, both cases combine into the single homogeneous $\F_2$-linear equation $\sum_{p' \in \mathcal{L}_0(n) \setminus \{p\}} \ell_p(p')y_{p'} + \ell_p(-q_n/p)y_p = 0$ over $\F_2$ (together with $c_2(y) \equiv 1 \pmod 8$, i.e., $\sum_{p \in \mathcal{L}_0(n)} \frac{p-1}{2}y_p = \sum_{p \in \mathcal{L}_0(n)} \frac{p^2-1}{8}y_p = 0$ in $\F_2$, when $2 \mid D_n$). Let $R_n^{\mathrm{ram}}$ denote the resulting ramified R\'edei symbol matrix over $\F_2$.

\begin{theorem}[{R\'edei--Reichardt 4-Rank Cohomology and Reduced Dual Quartic Laws}]\label{thm:5.3}
Let $n \ge 8$ be an integer.
\begin{enumerate}
\item[(i)] \textbf{Galois Cohomology Exact Sequence and R\'edei--Reichardt 4-Rank Identity:} Hilbert's Theorem~90 yields the exact sequence of $\F_2$-vector spaces
\begin{equation}\label{eq:5.8}
0 \longrightarrow H^1\bigl(\Gal(K_n/\Q),\, \mathcal{O}_{K_n}^{\times,+}\bigr) \longrightarrow \ker_{\F_2}\bigl(R_n^{\mathrm{ram}}\bigr) \xrightarrow{\;\Psi_+\;} \Cl^+(q_n)[2] \cap 2\Cl^+(q_n) \longrightarrow 0,
\end{equation}
where $H^1(\Gal(K_n/\Q), \mathcal{O}_{K_n}^{\times,+}) = \{\overline{1}, \overline{\eta_n}\} \simeq \F_2$ is spanned by $y = \mathbf{0}$ (the $k = 2$ branch) and the continued-fraction midpoint vector $y^{\mathrm{mid}}(n) \in \F_2^{\mathcal{L}_0(n)}$ encoding $c_2(y^{\mathrm{mid}}(n)) = \delta_2^{\mathrm{mid}}(n)$ (the $k = 1$ branch), so that
\begin{equation}\label{eq:5.9}
\dim_{\F_2} \ker_{\F_2}\bigl(R_n^{\mathrm{ram}}\bigr) = 1 + \rank_4 \Cl^+(q_n).
\end{equation}
\item[(ii)] \textbf{Deterministic $\Peven^-(n)$-Absorption and Reduced Dual Quartic Laws:} For any solution $(n, m)$ to $n! + 1 = m^2$ with $u_1 = \frac{m-1}{2} = \delta_1(n)z_1^2$ and $u_2 = \frac{m+1}{2} = \delta_2(n)z_2^2$, every inert prime $q \in \Peven^-(n)$ divides $z_1$ if and only if $\bigl(\frac{\delta_2(n)}{q}\bigr) = +1$ and divides $z_2$ if and only if $\bigl(\frac{\delta_2(n)}{q}\bigr) = -1$. Writing $u_1 = \widetilde{B}_1 W_1^2$ and $u_2 = \widetilde{B}_2 W_2^2$, where $W_i := \prod_{q \in \Peven^+(n),\, q \mid z_i} q^{v_q(n!)/2}$ and $\widetilde{B}_i := u_i / W_i^2$ absorb the $2$-adic, odd-valuation, and inert-even prime powers, the remaining split-even factor $A_{n,+}^* := \prod_{q \in \Peven^+(n)} q^{v_q(n!)} = (W_1 W_2)^2$ satisfies the dual quartic equations
\begin{equation}\label{eq:5.10}
\widetilde{B}_2 A_{n,+}^* = \widetilde{B}_1 W_1^4 + W_1^2, \qquad \widetilde{B}_1 A_{n,+}^* = \widetilde{B}_2 W_2^4 - W_2^2.
\end{equation}
\end{enumerate}
\end{theorem}

\begin{proof}
(i) Let $\mathcal{A}_n := \{\mathfrak{a}(y) = \prod_{p \in \mathcal{L}_0(n)} \mathfrak{p}_p^{y_p} : y \in \F_2^{\mathcal{L}_0(n)}\} \simeq \F_2^{\mathcal{L}_0(n)}$ be the elementary abelian 2-group of square-free products of ramified prime ideals of $\mathcal{O}_{K_n}$ (augmented at $p = 2$ when $2 \mid D_n$). For each $p \in \mathcal{L}_0(n)$, the row of $R_n^{\mathrm{ram}}$ at $p$ asserts that $\bigl(\frac{c_2(y)}{p}\bigr) = +1$ when $y_p = 0$ and $\bigl(\frac{-c_1(y)}{p}\bigr) = +1$ when $y_p = 1$, which is equivalent to the local Hilbert symbol condition $\bigl(\frac{c_2(y), q_n}{v}\bigr) = +1$ at every place $v$ of $\Q$. By Hasse's norm theorem and Gauss's genus theory \cite{Gauss1801}, $R_n^{\mathrm{ram}} y = \mathbf{0}$ if and only if $c_2(y) \in \mathrm{N}_{K_n/\Q}(K_n^{\times,+})$, i.e., $[\mathfrak{a}(y)]_+ \in \Cl^+(q_n)[2] \cap 2\Cl^+(q_n)$. Moreover, $\mathfrak{a}(y) \in \ker(\Psi_+)$ if and only if $\mathfrak{a}(y) = (\gamma)\mathcal{O}_{K_n}$ for some totally positive $\gamma \in K_n^{\times,+}$, in which case $u := \sigma(\gamma)/\gamma \in \mathcal{O}_{K_n}^{\times,+}$ defines a 1-cocycle class $[u] \in H^1(\Gal(K_n/\Q), \mathcal{O}_{K_n}^{\times,+}) = \mathcal{O}_{K_n}^{\times,+}/(\mathcal{O}_{K_n}^{\times,+})^2 = \{\overline{1}, \overline{\eta_n}\} \simeq \F_2$ (since $\mathrm{N}(\eta_n) = +1$ by Lemma~\ref{lem:3.1}(i)). The trivial class $[u] = \overline{1}$ gives $y = \mathbf{0}$ ($k = 2$), while the non-trivial class $[u] = \overline{\eta_n}$ gives $\gamma^2 = c_2(y)\varepsilon_n$, which is the unique continued-fraction midpoint vector $y = y^{\mathrm{mid}}(n)$ ($k = 1$). By the R\'edei--Reichardt theorem, $\dim_{\F_2}(\Cl^+(q_n)[2] \cap 2\Cl^+(q_n)) = \rank_4 \Cl^+(q_n)$, proving \eqref{eq:5.8} and \eqref{eq:5.9}.

(ii) For each inert prime $q \in \Peven^-(n)$, $\bigl(\frac{-\delta_1(n)}{q}\bigr)\bigl(\frac{\delta_2(n)}{q}\bigr) = \bigl(\frac{-q_n}{q}\bigr) = -1$. In $\delta_2(n)z_2^2 - \delta_1(n)z_1^2 = 1$, $q \mid z_1 z_2$ with $\gcd(z_1, z_2) = 1$: if $q \mid z_1$, reducing modulo $q$ gives $\bigl(\frac{\delta_2(n)}{q}\bigr) = +1$; if $q \mid z_2$, reducing modulo $q$ gives $\bigl(\frac{-\delta_1(n)}{q}\bigr) = +1 \iff \bigl(\frac{\delta_2(n)}{q}\bigr) = -1$. Absorbing these deterministically partitioned prime powers into $\widetilde{B}_1, \widetilde{B}_2$ gives $u_1 = \widetilde{B}_1 W_1^2$ and $u_2 = \widetilde{B}_2 W_2^2$ with $(W_1 W_2)^2 = A_{n,+}^* = \prod_{q \in \Peven^+(n)} q^{v_q(n!)}$ and $\widetilde{B}_2 W_2^2 - \widetilde{B}_1 W_1^2 = 1$. Multiplying by $W_1^2$ and by $W_2^2$ yields \eqref{eq:5.10}.
\end{proof}

\section{Pairwise $\SL_2(\Z)$ Duality, the Capacity Diamond, and Proof of Theorem~\ref{thm:1.3}(i)}\label{sec:6}

Having established the single-solution arithmetic structure of the $k = 2$ and $k = 1$ regimes, we now investigate the simultaneous Diophantine structure of pairs of solutions $4 \le n_1 < n_2$ to \eqref{eq:1.1}. By combining a universal $\SL_2(\Z)$ four-GCD factorization system with the cross-symmetry law of Theorem~\ref{thm:5.2}(i) and odd-harmonic Vaaler diagonal non-vanishing, we prove the consecutive and square-shift exclusion laws, the Universal Capacity Diamond, and Theorem~\ref{thm:1.3}(i).

\subsection{The Universal $\SL_2(\Z)$ Four-GCD System and Twin Norm Identities in $\Z[\sqrt{P}]$}
Let $(n_1, m_1)$ and $(n_2, m_2)$ be any two solutions to $n! + 1 = m^2$ with $4 \le n_1 < n_2$, and set $P := \frac{n_2!}{n_1!} = \prod_{j=n_1+1}^{n_2} j$. Define the six positive integers
\begin{equation}\label{eq:6.1}
A_i := \frac{m_i - 1}{2}, \qquad B_i := \frac{m_i + 1}{2} \quad (i = 1, 2), \qquad U := \frac{m_2 - m_1}{2}, \qquad V := \frac{m_2 + m_1}{2}.
\end{equation}
Notice that $\gcd(A_1, B_1) = \gcd(A_2, B_2) = 1$, $A_1 B_1 = \frac{n_1!}{4}$, $A_2 B_2 = P\frac{n_1!}{4}$, $UV = (P - 1)\frac{n_1!}{4}$, and
\begin{equation}\label{eq:6.2}
U = A_2 - A_1 = B_2 - B_1, \qquad V = A_2 + B_1 = B_2 + A_1.
\end{equation}

\begin{proposition}[{Universal $\SL_2(\Z)$ Four-GCD System and Twin Norm Identities in $\Z[\sqrt{P}]$}]\label{prop:6.1}
Let $(n_1, m_1)$ and $(n_2, m_2)$ be any two solutions to $n! + 1 = m^2$ with $4 \le n_1 < n_2$, and set $P = \frac{n_2!}{n_1!}$. Then the four greatest common divisors $X := \gcd(A_1, A_2)$, $Y := \gcd(B_1, B_2)$, $Z := \gcd(A_1, B_2)$, and $W := \gcd(B_1, A_2)$ are pairwise coprime positive integers satisfying $A_1 = XZ$, $B_1 = YW$, and $XYZW = \frac{n_1!}{4}$. Moreover:
\begin{enumerate}
\item[(i)] There exist unique positive integers $a, b, c, d$ with $ab = P - 1$ and $cd = P$ such that $A_2 = cXW$, $B_2 = dYZ$, $U = aXY$, $V = bZW$,
\begin{equation}\label{eq:6.3}
\begin{pmatrix} X \\ Z \end{pmatrix} = \begin{pmatrix} b & -d \\ c & -a \end{pmatrix} \begin{pmatrix} W \\ Y \end{pmatrix}, \qquad dY^2 - cX^2 = b, \qquad cW^2 - dZ^2 = a,
\end{equation}
and the explicit four-squares identities hold:
\begin{equation}\label{eq:6.4}
\begin{aligned}
2ac\,X^2 &= Pm_1 - m_2 - (P - 1), \qquad & 2ad\,Y^2 &= Pm_1 - m_2 + (P - 1), \\
2bd\,Z^2 &= Pm_1 + m_2 - (P - 1), \qquad & 2bc\,W^2 &= Pm_1 + m_2 + (P - 1).
\end{aligned}
\end{equation}
\item[(ii)] In the quadratic ring $\Z[\sqrt{P}]$, the elements $\alpha_1 := dY + X\sqrt{P}$ (of norm $bd$) and $\alpha_2 := cW + Z\sqrt{P}$ (of norm $ac$) satisfy the exact algebraic identities
\begin{equation}\label{eq:6.5}
\alpha_1 \alpha_2 = Pm_1 + m_2\sqrt{P}, \qquad \alpha_1 \overline{\alpha_2} = P - \sqrt{P}, \qquad \frac{\alpha_1}{\alpha_2} = \frac{P - \sqrt{P}}{ac}.
\end{equation}
\item[(iii)] At every odd prime $p \le n_1$ with $p \nmid P(P - 1)$, the four characters $\chi_X(p) = \bigl(\frac{bd}{p}\bigr)$, $\chi_Y(p) = \bigl(\frac{-bc}{p}\bigr)$, $\chi_Z(p) = \bigl(\frac{ac}{p}\bigr)$, $\chi_W(p) = \bigl(\frac{-ad}{p}\bigr)$ of conductor $\le 4P(P - 1)$ satisfy $\chi_T(p) = +1$ whenever $p \mid T$ ($T \in \{X, Y, Z, W\}$), and obey the Quadrant IV lock
\begin{equation}\label{eq:6.6}
\left(1 + \left(\frac{-P}{p}\right)\right)\left(1 + \left(\frac{-(P - 1)}{p}\right)\right)\left(1 - \left(\frac{bd}{p}\right)\right) = 0.
\end{equation}
\end{enumerate}
\end{proposition}

\begin{proof}
Since $\gcd(A_1, B_1) = 1$ and $\gcd(A_2, B_2) = 1$, the four greatest common divisors $X = \gcd(A_1, A_2)$, $Y = \gcd(B_1, B_2)$, $Z = \gcd(A_1, B_2)$, and $W = \gcd(B_1, A_2)$ are pairwise coprime. Every prime power $p^e \parallel A_1 B_1 = n_1!/4$ divides either $A_1$ or $B_1$, and since $A_1 B_1 \mid A_2 B_2$, it also divides either $A_2$ or $B_2$. Thus $A_1 = \gcd(A_1, A_2 B_2) = XZ$, $B_1 = \gcd(B_1, A_2 B_2) = YW$, and $XYZW = A_1 B_1 = n_1!/4$.

(i) By definition, $XW \mid A_2$ and $YZ \mid B_2$ with $\gcd(XW, YZ) = 1$, while from \eqref{eq:6.2}, $X = \gcd(A_1, A_2) \mid U$ and $Y = \gcd(B_1, B_2) \mid U$ with $\gcd(X, Y) = 1$, so $XY \mid U$; similarly, $ZW \mid V$. Writing $A_2 = cXW$, $B_2 = dYZ$, $U = aXY$, and $V = bZW$ gives $cd = \frac{A_2 B_2}{XYZW} = P$ and $ab = \frac{UV}{XYZW} = P - 1$. Substituting these into the four linear relations of \eqref{eq:6.2} and dividing by $X, Y, W, Z$ respectively yields
\[
Z = cW - aY, \qquad W = dZ - aX, \qquad Y = bZ - cX, \qquad X = bW - dY,
\]
which is the $\SL_2(\Z)$ matrix system in \eqref{eq:6.3} of determinant $cd - ab = P - (P - 1) = 1$. Substituting $W = \frac{X + dY}{b}$ and $Z = \frac{Y + cX}{b}$ into $YW - XZ = B_1 - A_1 = 1$ gives $dY^2 - cX^2 = b$, and substituting $Y = \frac{cW - Z}{a}$ and $X = \frac{dZ - W}{a}$ into $YW - XZ = 1$ gives $cW^2 - dZ^2 = a$.

Multiplying $dY^2 - cX^2 = b$ by $a$ gives $adY^2 - acX^2 = ab = P - 1$. Squaring both sides and using $cd = P$, $2aXY = 2U = m_2 - m_1$, and $m_2^2 = Pm_1^2 - (P - 1)$ yields
\[
a^2\bigl(dY^2 + cX^2\bigr)^2 = (P - 1)^2 + P(m_2 - m_1)^2 = P^2 m_1^2 - 2Pm_1 m_2 + m_2^2 = \bigl(Pm_1 - m_2\bigr)^2.
\]
Since $Pm_1 > m_2$, taking positive square roots gives $adY^2 + acX^2 = Pm_1 - m_2$; adding and subtracting $adY^2 - acX^2 = P - 1$ gives $2acX^2 = Pm_1 - m_2 - (P - 1)$ and $2adY^2 = Pm_1 - m_2 + (P - 1)$, and the identical calculation on $cW^2 - dZ^2 = a$ with $2bZW = m_2 + m_1$ gives $2bdZ^2 = Pm_1 + m_2 - (P - 1)$ and $2bcW^2 = Pm_1 + m_2 + (P - 1)$.

(ii)--(iii) In $\Z[\sqrt{P}]$, expanding $\alpha_1\alpha_2 = (dY + X\sqrt{P})(cW + Z\sqrt{P}) = cdYW + PXZ + (cXW + dYZ)\sqrt{P} = P(YW + XZ) + (A_2 + B_2)\sqrt{P} = Pm_1 + m_2\sqrt{P}$ and $\alpha_1\overline{\alpha_2} = (dY + X\sqrt{P})(cW - Z\sqrt{P}) = P(YW - XZ) - (dYZ - cXW)\sqrt{P} = P - \sqrt{P}$ establishes \eqref{eq:6.5}. Finally, for any odd prime $p \le n_1$ with $p \nmid P(P - 1)$, $p$ divides $XYZW = n_1!/4$. Reducing $dY^2 - cX^2 = b$ and $cW^2 - dZ^2 = a$ modulo $p$ gives $\chi_X(p) = \bigl(\frac{bd}{p}\bigr) = +1$ if $p \mid X$, $\chi_Y(p) = \bigl(\frac{-bc}{p}\bigr) = +1$ if $p \mid Y$, $\chi_Z(p) = \bigl(\frac{ac}{p}\bigr) = +1$ if $p \mid Z$, and $\chi_W(p) = \bigl(\frac{-ad}{p}\bigr) = +1$ if $p \mid W$. Since $\chi_Y(p) = \bigl(\frac{-P}{p}\bigr)\chi_X(p)$, $\chi_Z(p) = \bigl(\frac{P(P-1)}{p}\bigr)\chi_X(p)$, and $\chi_W(p) = \bigl(\frac{-(P-1)}{p}\bigr)\chi_X(p)$, all four characters coincide with $\bigl(\frac{bd}{p}\bigr)$ whenever $\bigl(\frac{-P}{p}\bigr) = \bigl(\frac{-(P-1)}{p}\bigr) = +1$, proving \eqref{eq:6.6}.
\end{proof}

\subsection{Consecutive-Index Rigidity and Unconditional Square-Shift Exclusions}
Specializing Proposition~\ref{prop:6.1} to consecutive indices $(n, n + 1)$ and to square-ratio shifts $P = t^2$ yields the following unconditional exclusion laws.

\begin{proposition}[{Consecutive-Index Rigidity and Unconditional Square-Shift Exclusions}]\label{prop:6.2}
Let $(n_1, m_1)$ and $(n_2, m_2)$ be two solutions to $n! + 1 = m^2$ with $4 \le n_1 < n_2$.
\begin{enumerate}
\item[(i)] \textbf{Uniqueness of the $n = 4$ Boundary Collapse:} For consecutive indices $(n, n + 1)$, the boundary collapse $m_{n+1} = \frac{n!}{2} - 1$ (i.e., $X = W = 1$) occurs if and only if $n! = 4(n + 2)$, whose unique positive integer solution is $n = 4$, yielding the consecutive pair $(4, 5)$ and $(5, 11)$.
\item[(ii)] \textbf{Quarter-Factorial Equipartition and Exclusion of $(k_n, k_{n+1}) = (2, 2)$:} For consecutive solutions $(n, n + 1)$ with $n \ge 6$, all four GCD components satisfy $X, Y, Z, W > 1$ and $\log X, \log Y, \log Z, \log W = \frac{1}{4}\log(n!) + O(\log n)$, and $(k_n, k_{n+1}) \ne (2, 2)$ for all sufficiently large $n$.
\item[(iii)] \textbf{Unconditional Square-Shift Exclusions:} No two solutions $n_1 < n_2$ with $P < 4(n_1! + 1)$ can have $P = \frac{n_2!}{n_1!} = t^2$ a perfect square; in particular, for every integer $t \ge 2$, $(t^2 - 1)! + 1$ and $(t^2)! + 1$ can \textbf{never} both be perfect squares. Likewise, for consecutive squares $n_1 = j^2$ and $n_2 = (j + 1)^2$, $\log X, \log Y, \log Z, \log W = \frac{1}{4}\log((j^2)!) + O(j\log j)$ unconditionally excludes $(k_{j^2}, k_{(j+1)^2}) = (2, 2)$ for all sufficiently large $j$.
\end{enumerate}
\end{proposition}

\begin{proof}
(i) For $n \ge 4$, $M_n := \frac{n!}{2}$ satisfies $M_n > m_{n+1} > m_n > 1$. The boundary collapse $m_{n+1} = M_n - 1$ (so $B_{n+1} = dYZ = M_n/2$ and $X = W = 1$) gives upon squaring:
\[
(n + 1)n! + 1 = m_{n+1}^2 = (M_n - 1)^2 = \frac{(n!)^2}{4} - n! + 1 \iff n! = 4(n + 2).
\]
Since $n! < 4(n + 2)$ for $1 \le n \le 3$, $4! = 24 = 4(6)$ at $n = 4$, and $n! > 4(n + 2)$ strictly for all $n \ge 5$ by induction, $n = 4$ is the unique positive integer solution.

(ii) For $n \ge 6$, the four-squares identity $2acX^2 = (n + 1)m_n - m_{n+1} - n \ge (n + 1 - \sqrt{n + 1})\sqrt{n!} - n > 2n(n + 1) \ge 2ac$ gives $X, Y, Z, W > 1$ and $\log X, \log Y, \log Z, \log W = \frac{1}{4}\log(n!) + O(\log n)$. If $k_n = k_{n+1} = 2$, then $m_n + 1 = 2YW = 2X_1(q_n)^2$ and $m_{n+1} + 1 = 2dYZ = 2X_1(q_{n+1})^2$, so no prime $p \in \Podd(n)$ with $p \nmid (n + 1)$ can divide $Y, W$, or $Z$. Thus every prime $p \in \Podd(n)$ with $p \nmid (n + 1)$ divides $X$ to its full factorial power $p^{v_p(n!)}$, forcing $\log X \ge \Sodd(n) - O(n\log\log n) = (\frac{1}{2} - o(1))n\log n$, which contradicts $\log X = (\frac{1}{4} + o(1))n\log n$.

(iii) Suppose two solutions $4 \le n_1 < n_2$ satisfy $P = \frac{n_2!}{n_1!} = t^2$ for some integer $t \ge 2$ with $P < 4(n_1! + 1)$. Then $m_2^2 - 1 = P(m_1^2 - 1) = t^2(m_1^2 - 1)$ rearranges to
\[
t^2 m_1^2 - m_2^2 = t^2 - 1 \iff (tm_1 - m_2)(tm_1 + m_2) = t^2 - 1.
\]
Since $t^2 m_1^2 - m_2^2 = t^2 - 1 > 0$, we have $tm_1 > m_2$, so $tm_1 - m_2 \ge 1$ is a positive integer, forcing $tm_1 + m_2 \le t^2 - 1$. However, $tm_1 + m_2 > tm_1 = t\sqrt{n_1! + 1} > t^2 - 1$ whenever $t < 2\sqrt{n_1! + 1}$. In particular, for consecutive integers $n_1 = t^2 - 1$ and $n_2 = t^2$ ($t \ge 2$), we have $P = t^2 < 4((t^2 - 1)! + 1)$, so $(t^2 - 1)! + 1$ and $(t^2)! + 1$ can \textbf{never} both be perfect squares. Similarly, for consecutive squares $n_1 = j^2$ and $n_2 = (j + 1)^2$, we have $\log P = (2j + 1)\log(j^2) + O(j) = O(j\log j) = o(\log((j^2)!))$, so \eqref{eq:6.4} again forces $\log X, \log Y, \log Z, \log W = \frac{1}{4}\log((j^2)!) + O(j\log j)$, which unconditionally rules out $(k_{j^2}, k_{(j+1)^2}) = (2, 2)$ for all large $j$ by the same $\Podd(j^2)$ weight saturation $\Sodd(j^2) = (\frac{1}{2} - o(1))\log((j^2)!) > (\frac{1}{4} + o(1))\log((j^2)!)$.
\end{proof}

\subsection{Odd-Harmonic Vaaler Diagonal Non-Vanishing, the Universal Capacity Diamond, and Proof of Theorem~\ref{thm:1.3}(i)}
We now prove the Universal Capacity Diamond and complete the proof of Theorem~\ref{thm:1.3}(i). The crucial analytic observation is that when we correlate the valuation parities $(-1)^{v_p(n_1!)}$ and $(-1)^{v_p(n_2!)}$ of two solutions $n_1 < n_2$ with gap $h := n_2 - n_1 = n_1^\alpha$ ($\alpha \in [0, 1)$), the Vaaler trigonometric polynomial $\widetilde{\Psi}_H(x)$ for $(-1)^{\lfloor x \rfloor}$ is supported exclusively on odd harmonics $2m + 1 \ne 0$, so on the diagonal $2m_1 + 1 = -(2m_2 + 1)$ the difference phase $(2m_2 + 1)\frac{h}{2p}$ never vanishes.

\begin{lemma}[{Odd-Harmonic Vaaler Diagonal Non-Vanishing for Parity-Flip Primes}]\label{lem:6.3}
Let $4 \le n_1 < n_2$ with $h := n_2 - n_1 = n_1^\alpha$ ($\alpha \in [0, 1)$), and set $P = n_2!/n_1!$. Then the logarithmic factorial weight of the parity-flip primes $\{p \le n_1 : 2 \nmid v_p(P)\} = \{p \le n_1 : v_p(n_1!) \not\equiv v_p(n_2!) \pmod 2\}$ satisfies unconditionally as $n_1 \to \infty$:
\begin{equation}\label{eq:6.7}
\sum_{\substack{3 \le p \le n_1 \\ 2 \nmid v_p(P)}} v_p(n_1!)\log p = \left(\frac{\alpha}{2} + o(1)\right)n_1\log n_1.
\end{equation}
\end{lemma}

\begin{proof}
Set $y_2 := h(\log n_1)^2 + 2$ and $y_1 := \max\bigl(3, \frac{h}{(\log n_1)^{12}}\bigr)$.
\begin{enumerate}
\item[(1)] \emph{Upper Range $p \ge y_2$:} Any prime $p \ge y_2$ with $2 \nmid v_p(P)$ satisfies $v_p(P) \ge 1$, so $p \mid P = \prod_{j=n_1+1}^{n_1+h} j$. Since $v_p(n_1!) \le \frac{n_1}{p-1} \le \frac{n_1}{y_2-1}$ for all $p \ge y_2$, we have
\[
\sum_{\substack{y_2 \le p \le n_1 \\ 2 \nmid v_p(P)}} v_p(n_1!)\log p \le \frac{n_1}{y_2 - 1} \sum_{p \mid P} \log p \le \frac{n_1\log P}{y_2 - 1} \le \frac{n_1 h\log(2n_1)}{h(\log n_1)^2} = O\!\left(\frac{n_1}{\log n_1}\right).
\]
\item[(2)] \emph{Transition Collar $y_1 < p < y_2$ and Tiny Range $p \le n_1^\delta$ ($\delta \in (0, \alpha/4)$):} By Mertens' theorem, $\sum_{y_1 < p < y_2} v_p(n_1!)\log p \le n_1 \sum_{y_1 < p < y_2} \frac{\log p}{p-1} \ll n_1\log\log n_1 = o(n_1\log n_1)$, and $\sum_{3 \le p \le n_1^\delta} v_p(n_1!)\log p = (\delta + o(1))n_1\log n_1$.
\item[(3)] \emph{Oscillatory Range $n_1^\delta < p \le y_1$:} On this range, write $\mathbf{1}_{2 \nmid v_p(P)} = \frac{1}{2} - \frac{1}{2}(-1)^{v_p(n_1!)}(-1)^{v_p(n_2!)}$. The main term $\frac{1}{2}$ contributes $\frac{1}{2}\sum_{n_1^\delta < p \le y_1} v_p(n_1!)\log p = (\frac{\alpha-\delta}{2} + o(1))n_1\log n_1$. For the correlation sum $\sum_{n_1^\delta < p \le y_1} (-1)^{v_p(n_1!) + v_p(n_2!)} v_p(n_1!)\log p$, expanding $(-1)^{\lfloor n_1/p \rfloor}$ and $(-1)^{\lfloor n_2/p \rfloor}$ via the Vaaler polynomial $\widetilde{\Psi}_H$ of Lemma~\ref{lem:2.3} ($H = \lfloor (\log n_1)^3 \rfloor$), note that in $\widetilde{\Psi}_H(x) = 2\Psi_H(x) - 4\Psi_{2H+1}(x/2)$, every even harmonic cancels identically, leaving only odd frequencies $h_1 = 2m_1 + 1$ and $h_2 = 2m_2 + 1$. Thus for every pair of odd frequencies $(h_1, h_2)$, the $p^{-1}$ phase is
\[
\frac{h_1 n_1 + h_2 n_2}{2p} = \frac{(h_1 + h_2)n_1 + h_2 h}{2p}.
\]
If $h_1 + h_2 \ne 0$, then $|(h_1 + h_2)n_1 + h_2 h| \ge 2n_1 - 2Hh \gg n_1$, giving derivative parameter $T \gg n_1/P$. If $h_1 + h_2 = 0$ (the diagonal $h_1 = -h_2 = -(2m_2 + 1)$), then because $h_2 = 2m_2 + 1 \ne 0$ is odd, the phase is $\frac{(2m_2+1)h}{2p} \ne 0$ with derivative parameter $T = \frac{|2m_2+1|h}{2P} \ge \frac{h}{2y_1} \ge \frac{1}{2}(\log n_1)^{12}$. Applying Lemma~\ref{lem:2.4} on dyadic intervals in $(n_1^\delta, y_1]$ cancels the correlation sum to $o(n_1\log n_1)$. Letting $\delta \downarrow 0$ yields \eqref{eq:6.7}.
\end{enumerate}
\end{proof}

\begin{theorem}[{Universal Capacity Diamond and Proof of Theorem~\ref{thm:1.3}(i)}]\label{thm:6.4}
For any solution $(n, m)$ to $n! + 1 = m^2$, define its odd-core polarization weight $\theta(n) := \frac{\log B^{(2)}(n)}{\Sodd(n)} \in [0, 1]$ and, for any $\varepsilon \in (0, 1/4)$, its upper-band polarization weight
\[
\theta_\varepsilon(n) := \frac{\displaystyle\sum_{n^{1-\varepsilon/2} < p \le n,\, p \mid B^{(2)}(n)} v_p(n!)\log p}{\displaystyle\sum_{n^{1-\varepsilon/2} < p \le n,\, p \in \Podd(n)} v_p(n!)\log p} \in [0, 1]
\]
(where $\theta(n) = \theta_\varepsilon(n) = 0$ if $k_n = 2$). For any two solutions $4 \le n_1 < n_2$ with $n_2 - n_1 = n_1^\alpha$ ($\alpha \in [0, 1)$), the pair $(\theta(n_1), \theta(n_2))$ lies inside the \textbf{Capacity Diamond}
\begin{equation}\label{eq:6.8}
\frac{1 - \alpha}{2} - o(1) \le \theta(n_1) + \theta(n_2) \le \frac{3 + \alpha}{2} + o(1), \qquad \bigl|\theta(n_1) - \theta(n_2)\bigr| \le \frac{1 + \alpha}{2} + o(1).
\end{equation}
Consequently, any two solutions $n_1 < n_2$ in the $k = 2$ regime ($\theta(n_1) = \theta(n_2) = 0$) are separated by $n_2 - n_1 > n_1^{1-\varepsilon}$ for every $\varepsilon > 0$, proving Theorem~\ref{thm:1.3}(i), and for every $\eta \in (0, 1/2]$, any interval $[N', N' + (N')^{1-\varepsilon}]$ with $N' \ge N_0(\varepsilon, \eta)$ contains at most $\lfloor (\eta^2\varepsilon)^{-1} \rfloor + 1$ solutions with $|\theta_\varepsilon(n) - 1/2| \ge \eta$, whence $\#\{4 \le n \le N : |\theta_\varepsilon(n) - 1/2| \ge \eta\} \le C_\varepsilon \eta^{-2} N^\varepsilon$ for all $N \ge N_0(\varepsilon, \eta)$.
\end{theorem}

\begin{proof}
Let $4 \le n_1 < n_2$ be any two solutions with $h = n_2 - n_1 = n_1^\alpha$ ($\alpha \in [0, 1)$). From \eqref{eq:6.4}, since $\log P \le h\log n_2 = O(n_1^\alpha\log n_1) = o(n_1\log n_1)$, both $ZW$ and $XY$ satisfy
\begin{equation}\label{eq:6.9}
\log(ZW) = \frac{1}{2}\log(n_1!) + O(\log P), \qquad \log(XY) = \frac{1}{2}\log(n_1!) + O(\log P).
\end{equation}
Decompose $A_i^* = B^{(1)}(n_i)Z_1^*(n_i)^2$ and $B_i^* = B^{(2)}(n_i)Z_2^*(n_i)^2$ ($i = 1, 2$) as in Theorem~\ref{thm:5.2}(i). Consider any odd prime $p \le n_1$ with $2 \mid v_p(P)$ (so $v_p(n_1!) \equiv v_p(n_2!) \pmod 2$):
\begin{itemize}
\item If $p \mid Z = \gcd(A_1, B_2)$, then either $p \in \Podd(n_1) \cap \Podd(n_2)$ (so $p^{v_p(n_1!)} \mid B^{(2)}(n_2)$) or $p \in \Peven(n_1) \cap \Peven(n_2)$ (so $p^{v_p(n_1!)} \mid Z_1^*(n_1)^2$).
\item If $p \mid W = \gcd(B_1, A_2)$, then either $p \in \Podd(n_1) \cap \Podd(n_2)$ (so $p^{v_p(n_1!)} \mid B^{(2)}(n_1)$) or $p \in \Peven(n_1) \cap \Peven(n_2)$ (so $p^{v_p(n_1!)} \mid Z_1^*(n_2)^2$).
\end{itemize}
Consequently, $ZW$ divides $2^{v_2(n_1!)} B^{(2)}(n_1)B^{(2)}(n_2)Z_1^*(n_1)^2 Z_1^*(n_2)^2 \prod_{3 \le p \le n_1,\, 2 \nmid v_p(P)} p^{v_p(n_1!)}$. Taking logarithms and applying the cross-symmetry law \eqref{eq:5.4} ($\log(B^{(2)}(n_i)Z_1^*(n_i)^2) = (\theta(n_i) + o(1))n_1\log n_1$) together with Lemma~\ref{lem:6.3} and \eqref{eq:6.9} gives
\[
\frac{1}{2}n_1\log n_1 \le \left(\theta(n_1) + \theta(n_2) + \frac{\alpha}{2} + o(1)\right)n_1\log n_1 \implies \theta(n_1) + \theta(n_2) \ge \frac{1 - \alpha}{2} - o(1).
\]
Repeating the identical divisibility argument with $(B^{(1)}, Z_2^*)$ (where $\log(B^{(1)}(n_i)Z_2^*(n_i)^2) = (1 - \theta(n_i) + o(1))n_1\log n_1$) gives $\theta(n_1) + \theta(n_2) \le \frac{3+\alpha}{2} + o(1)$, and applying both arguments to $XY = \gcd(A_1, A_2)\gcd(B_1, B_2)$ gives $|\theta(n_1) - \theta(n_2)| \le \frac{1+\alpha}{2} + o(1)$, establishing \eqref{eq:6.8}.

When both $n_1$ and $n_2$ lie in the $k = 2$ regime, $\theta(n_1) = \theta(n_2) = 0$, so \eqref{eq:6.8} forces $\frac{1-\alpha}{2} \le o(1) \implies \alpha \ge 1 - o(1)$, proving $n_2 - n_1 > n_1^{1-\varepsilon}$ and $\#\mathcal{B}_{k=2}(N) \ll_\varepsilon N^\varepsilon$ (Theorem~\ref{thm:1.3}(i)).

Finally, in the weighted Hilbert space $\ell^2(w)$ on odd primes $p \le N'$ with weights $w_p = \frac{v_p(N'!)\log p}{\Sodd(N') + \Seven(N')}$, each solution $n \in I = [N', N' + (N')^{1-\varepsilon}]$ defines a unit sign vector $\mathbf{u}_n$ with coordinates $u_n(p) = \varepsilon_p(n) \in \{\pm 1\}$ given by $\varepsilon_p(n) := +1$ if $p \mid \frac{m+1}{2}$ and $\varepsilon_p(n) := -1$ if $p \mid \frac{m-1}{2}$. For any two solutions $n_1 < n_2$ in $I$, we have $\varepsilon_p(n_1)\varepsilon_p(n_2) = +1$ for $p \mid XY$ and $\varepsilon_p(n_1)\varepsilon_p(n_2) = -1$ for $p \mid ZW$, so $|\langle \mathbf{u}_{n_1}, \mathbf{u}_{n_2} \rangle_w| = \bigl|\frac{\log(XY) - \log(ZW)}{\log(N'!)}\bigr| + O((N')^{-\varepsilon}) \le 3(N')^{-\varepsilon}$ by \eqref{eq:6.9}, while since $\sum_{p \in \Podd(N') \cap ((N')^{1-\varepsilon/2}, N']} w_p = \frac{\varepsilon}{4} + o(1)$ by Theorem~\ref{thm:2.5}(iii), any $\eta$-polarized solution ($|\theta_\varepsilon(n) - 1/2| \ge \eta$) has inner product $\ge \eta\sqrt{\varepsilon} - o(1)$ in absolute value with the normalized indicator vector $\mathbf{v}_I$ of $\Podd(N') \cap ((N')^{1-\varepsilon/2}, N']$. Applying Rankin's spherical cap Cauchy--Schwarz inequality $t^2\eta_0^2 \le t + 3t(t - 1)(N')^{-\varepsilon}$ bounds the number of $\eta$-polarized solutions in $I$ by $\lfloor (\eta^2\varepsilon)^{-1} \rfloor + 1$ for $N' \ge N_0(\varepsilon, \eta)$, and summing over dyadic intervals yields $\#\{4 \le n \le N : |\theta_\varepsilon(n) - 1/2| \ge \eta\} \le C_\varepsilon \eta^{-2} N^\varepsilon$.
\end{proof}

\section{Support-Reduced Pad\'e--Wronskian Approximations and Proof of Theorem~\ref{thm:1.3}(ii)}\label{sec:7}

In this section we prove Theorem~\ref{thm:1.3}(ii). Let $P(x) \in \Z[x]$ be a polynomial of degree $r \ge 2$ and let $s \in \Z \setminus \{0\}$. By the classical depression substitution $x \mapsto r a_r x + a_{r-1}$ and scaling of $s$ by $(r a_r)^{2r}$ (Berend and Osgood \cite[Lemma~1]{BerendOsgood1992}; Bui, Pratt, and Zaharescu \cite[Section~3]{BuiPrattZaharescu2023}), we may assume without loss of generality that
\begin{equation}\label{eq:7.1}
P(x) = a_r x^r + a_{r-\nu} x^{r-\nu} + \dots + a_0 \in \Z[x], \qquad a_r > 0, \quad s > 0,
\end{equation}
where $a_{r-1} = 0$ and $\nu := r - \deg(P(x) - a_r x^r) \in \{2, \dots, r\}$ denotes the degree gap between the leading term $a_r x^r$ and the next non-zero term of $P(x)$ (for Brocard's equation $n! + 1 = m^2$, $P(x) = x^2 - 1$ has $r = 2$ and $\nu = 2$).

\subsection{The $(m + 1)$-Variable Binomial-Product Pad\'e System and Siegel's Lemma}
Fix an integer $m \ge 2$ and a real parameter $\varepsilon_0 \in (0, m - 1)$. Let $M := \lfloor N^\theta \rfloor$ with $0 < \theta < \frac{1}{m+2}$. Let $r \le \beta_1 < \beta_2 < \dots < \beta_m \le M$ be positive integer multiples of $r$, and consider the $m + 1$ algebraic functions
\begin{equation}\label{eq:7.2}
\omega_0(x) := 1, \qquad \omega_i(x) := \prod_{j=1}^{\beta_i-1} (1 - jx)^{-1/r} = \sum_{\ell=0}^\infty b_{i,\ell}\,x^\ell \qquad (1 \le i \le m)
\end{equation}
analytic in the disk $|x| < \frac{1}{\beta_m - 1}$. By the binomial expansion $(1 - jx)^{-1/r} = \sum_{u=0}^\infty \binom{-1/r}{u}(-jx)^u$ and Cauchy's integral estimates \cite[Lemma~4.2]{BuiPrattZaharescu2023}, each Taylor coefficient $b_{i,\ell} \in \Q$ satisfies $r^{2\ell} b_{i,\ell} \in \Z$ and $|b_{i,\ell}| \le 2^{\beta_i}(2\beta_i)^\ell \le 2^M(2M)^\ell$.

\begin{lemma}[{Initial $(m + 1)$-Variable Pad\'e Polynomials}]\label{lem:7.1}
Let $D \ge 10m$ be an integer, let $\varepsilon_0 \in (0, m - 1)$, and set $O := \lfloor (m + 1 - \varepsilon_0)D \rfloor$. There exist polynomials $P_0(x), P_1(x), \dots, P_m(x) \in \Z[x]$ of degree at most $D$, not all identically zero, such that the Pad\'e remainder series
\begin{equation}\label{eq:7.3}
R(x) := \sum_{i=0}^m P_i(x)\,\omega_i(x) = \sum_{v=0}^\infty r_v\,x^v
\end{equation}
satisfies $r_v = 0$ for all $0 \le v \le O$ (so $\ord_{x=0} R(x) \ge O + 1$), and the coefficients $p_{i,k}$ of $P_i(x) = \sum_{k=0}^D p_{i,k} x^k$ ($0 \le i \le m$) and Taylor coefficients $r_v$ ($v > O$) satisfy the upper bounds
\begin{equation}\label{eq:7.4}
\begin{aligned}
|p_{i,k}| &\le 4^{m\varepsilon_0^{-1}M} \bigl(2mr^{\frac{2m}{m-\varepsilon_0}} M\bigr)^{\frac{(m-\varepsilon_0)(m+1-\varepsilon_0)}{\varepsilon_0}D}, \\
|r_v| &\le (m+1)(D+1)\, 4^{m\varepsilon_0^{-1}M} \bigl(2mr^{\frac{2m}{m-\varepsilon_0}} M\bigr)^{\frac{(m-\varepsilon_0)(m+1-\varepsilon_0)}{\varepsilon_0}D} (2M)^v.
\end{aligned}
\end{equation}
\end{lemma}

\begin{proof}
Write $P_i(x) = \sum_{k=0}^D p_{i,k} x^k$. Expanding Cauchy products in \eqref{eq:7.3} shows that the condition $r_v = 0$ for $0 \le v \le O$ is equivalent to the homogeneous linear system of $O + 1$ equations
\begin{equation}\label{eq:7.5}
r^{2O} p_{0,v} \cdot \mathbf{1}_{v \le D} + \sum_{i=1}^m \sum_{k=0}^{\min(v,D)} p_{i,k}\bigl(r^{2O} b_{i,v-k}\bigr) = 0 \qquad (0 \le v \le O),
\end{equation}
where every coefficient $r^{2O} b_{i,v-k}$ is an integer of absolute value at most $r^{2O} 2^M (2M)^O$. Consider first the subsystem of \eqref{eq:7.5} indexed by $D + 1 \le v \le O$: since $p_{0,v} \cdot \mathbf{1}_{v \le D} = 0$ for $v \ge D + 1$, this subsystem consists of
\[
M_{\mathrm{eq}} := O - D \le (m - \varepsilon_0)\,D
\]
homogeneous linear equations with integer coefficients in the
\[
N_{\mathrm{var}} := m(D + 1) > mD
\]
unknowns $\{p_{i,k} : 1 \le i \le m,\, 0 \le k \le D\}$. Since $\varepsilon_0 \in (0, m - 1)$, we have $N_{\mathrm{var}} - M_{\mathrm{eq}} > \varepsilon_0 D > 0$, so the Dirichlet box-principle exponent in Siegel's lemma \cite[Lemma~4.3]{BuiPrattZaharescu2023} is bounded above by
\[
\frac{M_{\mathrm{eq}}}{N_{\mathrm{var}} - M_{\mathrm{eq}}} \le \frac{(m - \varepsilon_0)D}{\varepsilon_0 D} = \frac{m - \varepsilon_0}{\varepsilon_0}.
\]
Applying Siegel's lemma yields a non-zero integer solution $\{\tilde{p}_{i,k} : 1 \le i \le m,\, 0 \le k \le D\}$ bounded in absolute value by $\bigl(m(D + 1)r^{2O} 2^M (2M)^O\bigr)^{\frac{m-\varepsilon_0}{\varepsilon_0}}$. Multiplying each $\tilde{p}_{i,k}$ by $r^{2O}$ to define $p_{i,k} := r^{2O}\tilde{p}_{i,k}$ ($1 \le i \le m$) ensures that $r^{2O} \mid p_{i,k}$, so the remaining $D + 1$ equations of \eqref{eq:7.5} for $0 \le v \le D$ uniquely define integers
\[
p_{0,v} := -\sum_{i=1}^m \sum_{k=0}^v \frac{p_{i,k}}{r^{2O}}\bigl(r^{2O} b_{i,v-k}\bigr) \in \Z \qquad (0 \le v \le D).
\]
Using $O \le (m + 1 - \varepsilon_0)D$ and $|b_{i,v-k}| \le 2^M(2M)^v$ yields \eqref{eq:7.4}.
\end{proof}

\subsection{Puiseux Linear Independence, the Support-Reduced Wronskian, and Gauss's Lemma Multiplicity Lock}
Define the logarithmic derivatives $A_0(x) := 0$ and
\begin{equation}\label{eq:7.6}
A_i(x) := \frac{\omega_i'(x)}{\omega_i(x)} = \sum_{j=1}^{\beta_i-1} \frac{j}{r(1 - jx)} \quad (1 \le i \le m), \qquad T(x) := r\prod_{j=1}^M (1 - jx) \in \Z[x],
\end{equation}
so that $T(x)A_i(x) \in \Z[x]$ and $\deg(T) = M$. For every integer $k \ge 0$ and $0 \le i \le m$, define the differential iterates
\begin{equation}\label{eq:7.7}
\begin{aligned}
P_i^{[k]}(x) &:= \left(T(x)\frac{d}{dx} + T(x)A_i(x)\right)^k P_i(x) \in \Z[x], \\
R^{[k]}(x) &:= \left(T(x)\frac{d}{dx}\right)^k R(x) = \sum_{i=0}^m P_i^{[k]}(x)\,\omega_i(x).
\end{aligned}
\end{equation}

\begin{lemma}[{Linear Independence of $\omega_0, \dots, \omega_m$ and Monomial Products over $\C(x)$}]\label{lem:7.2}
For any integer $d \ge 1$, any shifts $r \le \beta_1 < \beta_2 < \dots < \beta_d$ ($r \mid \beta_i$), and any degree bound $1 \le E \le r - 1$, the $\binom{d+E}{d}$ monomial-product functions
\begin{equation}\label{eq:7.8}
\Omega_{\mathbf{e}}(x) := \prod_{i=1}^d \omega_i(x)^{e_i} = \prod_{i=1}^d \prod_{j=1}^{\beta_i-1} (1 - jx)^{-e_i/r}, \qquad \mathbf{e} = (e_1, \dots, e_d) \in \Z_{\ge 0}^d, \quad |\mathbf{e}| := \sum_{i=1}^d e_i \le E,
\end{equation}
are linearly independent over $\C(x)$. In particular, taking $E = 1$ and $d = m$, the functions $\omega_0(x), \omega_1(x), \dots, \omega_m(x)$ are linearly independent over $\C(x)$.
\end{lemma}

\begin{proof}
We proceed by induction on $d \ge 1$. Suppose there exists a relation $\sum_{|\mathbf{e}| \le E} U_{\mathbf{e}}(x)\Omega_{\mathbf{e}}(x) = 0$ with $U_{\mathbf{e}}(x) \in \C[x]$. Grouping terms according to the last exponent $e_d \in \{0, 1, \dots, E\}$ gives
\begin{equation}\label{eq:7.9}
\sum_{e_d=0}^E \Phi_{e_d}(x)\,\omega_d(x)^{e_d} = 0, \qquad \Phi_{e_d}(x) := \sum_{\substack{(e_1,\dots,e_{d-1}) \\ e_1 + \dots + e_{d-1} \le E - e_d}} U_{(e_1,\dots,e_d)}(x) \prod_{i=1}^{d-1} \omega_i(x)^{e_i}.
\end{equation}
Consider the local Puiseux expansion in a punctured neighborhood of the branch point $x_* := \frac{1}{\beta_d - 1}$. Because $\beta_i - 1 < \beta_d - 1$ for all $1 \le i \le d - 1$, each function $\omega_i(x)$ ($1 \le i \le d - 1$) is holomorphic and non-vanishing at $x = x_*$. Consequently, each coefficient $\Phi_{e_d}(x)$ is holomorphic at $x = x_*$, so if $\Phi_{e_d}(x) \not\equiv 0$, its order of vanishing at $x = x_*$ is a non-negative integer $k_{e_d} \in \Z_{\ge 0}$. On the other hand, $\omega_d(x) = (1 - (\beta_d - 1)x)^{-1/r} \prod_{j=1}^{\beta_d-2} (1 - jx)^{-1/r}$ has fractional Puiseux valuation $-\frac{1}{r}$ at $x = x_*$, so the $e_d$-th term $\Phi_{e_d}(x)\omega_d(x)^{e_d}$ has Puiseux valuation
\[
v_{x_*}\bigl(\Phi_{e_d}(x)\omega_d(x)^{e_d}\bigr) = k_{e_d} - \frac{e_d}{r} \in -\frac{e_d}{r} + \Z.
\]
Since $0 \le e_d \le E \le r - 1$, the cosets $-\frac{e_d}{r} + \Z$ ($0 \le e_d \le E$) in $\Q/\Z$ are pairwise distinct. Therefore no two terms with different $e_d$ can cancel in the Puiseux expansion at $x = x_*$, forcing $\Phi_{e_d}(x) \equiv 0$ for every $0 \le e_d \le E$, which by the induction hypothesis on $d - 1$ forces $U_{\mathbf{e}}(x) \equiv 0$ for all $|\mathbf{e}| \le E$.
\end{proof}

When $\varepsilon_0 \ge 1$ in Lemma~\ref{lem:7.1}, some of the polynomials $P_i(x)$ may vanish identically, causing the full $(m + 1) \times (m + 1)$ Wronskian determinant $\det(P_i^{[k]}(x))_{0 \le k,i \le m}$ to have a column of zeros. We overcome this obstacle by passing to the non-zero support $S \subset \{0, 1, \dots, m\}$.

\begin{lemma}[{Support-Reduced Wronskian Non-Vanishing and Gauss's Lemma Multiplicity Bound}]\label{lem:7.3}
Let $(P_0, \dots, P_m)$ be the Pad\'e polynomials constructed in Lemma~\ref{lem:7.1} for $\varepsilon_0 \in (0, m - 1)$, and let
\begin{equation}\label{eq:7.10}
S := \bigl\{0 \le i \le m : P_i(x) \not\equiv 0\bigr\} = \{i_1 < i_2 < \dots < i_s\}
\end{equation}
be the non-zero support of $(P_0, \dots, P_m)$, of cardinality $s := |S|$. Assume $D > (m + 1)^2 M$ and, if $m + 1 - \varepsilon_0 \notin \Z$, $D > \frac{\binom{m+1}{2}M + m + 1}{\{m + 1 - \varepsilon_0\}}$ (where $\{x\} := x - \lfloor x \rfloor$). Then:
\begin{enumerate}
\item[(i)] The $s \times s$ support-reduced Wronskian determinant
\begin{equation}\label{eq:7.11}
\Delta_S(x) := \det\Bigl(P_{i_a}^{[k]}(x)\Bigr)_{0 \le k \le s-1,\, 1 \le a \le s} \in \Z[x]
\end{equation}
is \textbf{not identically zero} ($\Delta_S(x) \not\equiv 0$), has degree $\deg(\Delta_S) \le sD + \binom{s}{2}M$, and vanishes at $x = 0$ to order $\ord_{x=0} \Delta_S(x) \ge O - (s - 1)$. Consequently, the support size satisfies $\lceil m + 1 - \varepsilon_0 \rceil \le s \le m + 1$ (in particular, $s \ge 2$), and the sum of the multiplicities of all non-zero roots of $\Delta_S(x)$ in $\C^\times$ satisfies
\begin{equation}\label{eq:7.12}
\sum_{\alpha \in \C^\times} \ord_{x=\alpha} \Delta_S(x) \le \deg(\Delta_S) - \ord_{x=0} \Delta_S(x) \le \bigl(s - (m + 1 - \varepsilon_0)\bigr)D + \binom{s}{2}M + s.
\end{equation}
\item[(ii)] For any $\alpha = 1/n_0 \in \Q \setminus \{0\}$ with $n_0 > M$, let $a := \ord_{x=\alpha} \Delta_S(x) \ge 0$, and define the normalized operators $P_i^k(x) := \frac{T(x)^k}{k!}\bigl(\frac{d}{dx} + A_i(x)\bigr)^k P_i(x) \in \Q[x]$. Then the $(a + s) \times s$ rational matrix $\bigl(P_{i_a}^k(\alpha)\bigr)_{0 \le k \le a+s-1,\, 1 \le a \le s}$ has full column rank $s$.
\item[(iii)] For any distinct integers $N \le n_1 < n_2 < \dots < n_t < 2N$, let $a_r := \ord_{x=1/n_r} \Delta_S(x)$. Then $\prod_{r=1}^t (n_r x - 1)^{a_r}$ divides $\Delta_S(x)$ in $\Z[x]$, so $\prod_{r=1}^t n_r^{a_r}$ divides the non-zero integer leading coefficient $c_{\mathrm{lead}}(\Delta_S) \in \Z \setminus \{0\}$, yielding the Gauss's Lemma multiplicity bound
\begin{equation}\label{eq:7.13}
\sum_{r=1}^t a_r \le s\frac{(m - \varepsilon_0)(m + 1 - \varepsilon_0)}{\varepsilon_0}\frac{\log M}{\log N}D + O_m\!\left(\frac{D + M\log N}{\log N}\right).
\end{equation}
\end{enumerate}
\end{lemma}

\begin{proof}
(i) Since $(P_0, \dots, P_m) \ne (0, \dots, 0)$, we have $s = |S| \ge 1$. Because $P_{i_a}(x) \not\equiv 0$ for every $1 \le a \le s$ by definition of $S$ and $\omega_{i_1}(x), \dots, \omega_{i_s}(x)$ are linearly independent over $\C(x)$ by Lemma~\ref{lem:7.2}, Siegel's differential Wronskian lemma \cite[Chapter~2, Lemma~4]{Siegel1949} (cf.\ \cite[Lemma~5.4]{BuiPrattZaharescu2023}) applies directly: if $\Delta_S(x) \equiv 0$, then the rows of $(P_{i_a}^{[k]}(x))_{0 \le k \le s-1,\, 1 \le a \le s}$ satisfy a non-trivial relation $\sum_{k=0}^\mu B_k(x) P_{i_a}^{[k]}(x) = 0$ over $\C[x]$ of minimal order $\mu \le s - 1$, which implies that the $s$ non-zero functions $f_a(x) := P_{i_a}(x)\omega_{i_a}(x)$ ($1 \le a \le s$) all satisfy the homogeneous linear differential equation $\sum_{k=0}^\mu B_k(x)(T\frac{d}{dx})^k f = 0$ of order $\mu \le s - 1$. Since the solution space of an order-$\mu$ linear ODE has dimension $\mu \le s - 1$ over $\C$, the functions $f_1, \dots, f_s$ would be $\C$-linearly dependent, yielding a non-trivial relation $\sum_{a=1}^s c_a P_{i_a}(x)\omega_{i_a}(x) = 0$ with $c_a \in \C$ and $P_{i_a}(x) \not\equiv 0$, contradicting Lemma~\ref{lem:7.2}. Thus $\Delta_S(x) \not\equiv 0$.

Since $P_i(x) \equiv 0$ for all $i \notin S$, we have $R^{[k]}(x) = \sum_{a=1}^s P_{i_a}^{[k]}(x)\omega_{i_a}(x)$ for every $k \ge 0$. Multiplying the first column of $\Delta_S(x)$ by $\omega_{i_1}(x)$ and adding $\omega_{i_a}(x)$ times the $a$-th column for $2 \le a \le s$ replaces the first column by $(R^{[k]}(x))_{0 \le k \le s-1}$. Expanding by cofactors $\Delta_{1,k}(x) \in \Z[x]$ along the first column gives
\begin{equation}\label{eq:7.14}
\Delta_S(x)\,\omega_{i_1}(x) = \sum_{k=0}^{s-1} \Delta_{1,k}(x)\,R^{[k]}(x).
\end{equation}
Since $\ord_{x=0} R(x) \ge O + 1$ and each application of $T(x)\frac{d}{dx}$ decreases the order of vanishing at $x = 0$ by at most 1, we have $\ord_{x=0} R^{[k]}(x) \ge O + 1 - k \ge O - (s - 1)$ for all $0 \le k \le s - 1$. Since $\omega_{i_1}(0) = 1 \ne 0$, \eqref{eq:7.14} gives $\ord_{x=0} \Delta_S(x) \ge O - (s - 1)$. Moreover, $\deg(P_{i_a}^{[k]}) \le D + kM$, so $\deg(\Delta_S) \le sD + \binom{s}{2}M$. Comparing $\deg(\Delta_S) \ge \ord_{x=0} \Delta_S(x)$ gives $sD + \binom{s}{2}M \ge (m + 1 - \varepsilon_0)D - s$, which under our lower bound on $D$ forces $s \ge \lceil m + 1 - \varepsilon_0 \rceil \ge 2$ and proves \eqref{eq:7.12}.

(ii) Since $T(\alpha) = r\prod_{j=1}^M (1 - j/n_0) \ne 0$ for $n_0 > M$, the operators $(T\frac{d}{dx})^k$ and $\frac{T(x)^k}{k!}\frac{d^k}{dx^k}$ span the same $\Q$-vector space of differential operators at $x = \alpha$; differentiating \eqref{eq:7.14} up to order $a = \ord_{x=\alpha} \Delta_S(x)$ at $x = \alpha$ via Leibniz's rule shows that $(P_{i_a}^k(\alpha))_{0 \le k \le a+s-1,\, 1 \le a \le s}$ has rank $s$ \cite[Lemmas~5.6 and 6.3]{BuiPrattZaharescu2023}.

(iii) Since $n_r \ge N \ge 2$, each linear polynomial $n_r x - 1 \in \Z[x]$ is primitive and irreducible in the unique factorization domain $\Z[x]$, and distinct integers $n_1 < \dots < n_t$ give pairwise coprime polynomials $n_r x - 1$ in $\Q[x]$ and hence non-associate prime elements in $\Z[x]$. Since $(n_r x - 1)^{a_r}$ divides $\Delta_S(x)$ in $\Q[x]$ and $n_r x - 1$ is primitive, Gauss's Lemma implies that $\prod_{r=1}^t (n_r x - 1)^{a_r}$ divides $\Delta_S(x)$ in $\Z[x]$:
\[
\Delta_S(x) = Q(x) \prod_{r=1}^t (n_r x - 1)^{a_r}, \qquad Q(x) \in \Z[x] \setminus \{0\}.
\]
Comparing the leading coefficients on both sides gives $\prod_{r=1}^t n_r^{a_r} \mid c_{\mathrm{lead}}(\Delta_S) \in \Z \setminus \{0\}$. Since $n_r \ge N$, we have $N^{\sum_{r=1}^t a_r} \le \prod_{r=1}^t n_r^{a_r} \le |c_{\mathrm{lead}}(\Delta_S)|$. Bounding $|c_{\mathrm{lead}}(\Delta_S)|$ by Hadamard's inequality and \eqref{eq:7.4} gives $|c_{\mathrm{lead}}(\Delta_S)| \le (2mr^{\frac{2m}{m-\varepsilon_0}} M)^{s\frac{(m-\varepsilon_0)(m+1-\varepsilon_0)}{\varepsilon_0}D + O_m(D/\log M + M)}$, and taking logarithms and dividing by $\log N$ yields \eqref{eq:7.13}.
\end{proof}

Before completing the proof of Theorem~\ref{thm:1.3}(ii), we record the explicit Leibniz derivative and factorial-root approximation estimates at $x = 1/n_0$.

\begin{lemma}[{Leibniz Derivative Estimates and Factorial-Root Approximation at $x = 1/n_0$}]\label{lem:7.4}
Under the hypotheses of Lemma~\ref{lem:7.1} and Lemma~\ref{lem:7.3}, let $N \le n_0 < 2N$ with $n_0 > 2M$, and set $P_i^k(x) := \frac{T(x)^k}{k!}\bigl(\frac{d}{dx} + A_i(x)\bigr)^k P_i(x) \in \Q[x]$ and $R^k(x) := \frac{T(x)^k}{k!}\frac{d^k}{dx^k} R(x) = \sum_{i=0}^m P_i^k(x)\omega_i(x)$.
\begin{enumerate}
\item[(i)] For every $0 \le k \le a + s - 1$, the rational numbers $P_i^k(1/n_0)$ ($0 \le i \le m$) have a common denominator dividing $Z := r^{a+s} n_0^{D + (a+s)M}$, and satisfy the upper bounds
\begin{equation}\label{eq:7.15}
\bigl|P_i^k(1/n_0)\bigr| \le u\,U^D, \qquad \bigl|R^k(1/n_0)\bigr| \le w\,W^{-D},
\end{equation}
where $U = (2mr^{\frac{2m}{m-\varepsilon_0}} M)^{\frac{(m-\varepsilon_0)(m+1-\varepsilon_0)}{\varepsilon_0}}$, $W = N^{m+1-\varepsilon_0}(8mr^{\frac{2m}{m-\varepsilon_0}} M)^{-\frac{m(m+1-\varepsilon_0)}{\varepsilon_0}}$, and the sub-exponential factors are $u = (4rN)^{(a+s)M + a + s} 10^{m\varepsilon_0^{-1}M}$ and $w = (m+1)(D+1)\allowbreak(2rN)^{2(a+s)M + 2(a+s)}\allowbreak \cdot 4^{m\varepsilon_0^{-1}M}$.
\item[(ii)] Suppose $sn_0! = P(x_0)$ and $s(n_0 - \beta_i)! = P(x_i)$ ($1 \le i \le d$), where $P(x) = a_r x^r + a_{r-\nu} x^{r-\nu} + \dots + a_0$ has degree $r \ge 2$ and degree gap $\nu \in \{2, \dots, r\}$. Then for every multi-index $\mathbf{e} = (e_1, \dots, e_d) \in \Z_{\ge 0}^d$ with $|\mathbf{e}| \le E \le \min(r - 1, \nu - 1)$, setting $q_{\mathbf{e}} := x_0^{E - |\mathbf{e}|} \prod_{i=1}^d \bigl(n_0^{\beta_i/r} x_i\bigr)^{e_i} \in \N$, we have
\begin{equation}\label{eq:7.16}
\left|\Omega_{\mathbf{e}}\!\left(\frac{1}{n_0}\right) - \frac{q_{\mathbf{e}}}{x_0^E}\right| \ll_{P,E} M^E x_0^{-\nu} \le x_0^{-\nu + o(1)}.
\end{equation}
\end{enumerate}
\end{lemma}

\begin{proof}
(i) Since $T(x) = r\prod_{j=1}^M (1 - jx)$ and $T(x)A_i(x) \in \Z[x]$ have degree $\le M$, induction on $k$ via Leibniz's rule shows that $r^k k!\,P_i^k(x) = r^k P_i^{[k]}(x) + \sum_{j=0}^{k-1} c_{k,j}(x) P_i^{[j]}(x)$ is a polynomial in $\Z[x]$ of degree $\le D + kM$. Evaluating at $x = 1/n_0$ shows that $r^{a+s} n_0^{D+(a+s)M} P_i^k(1/n_0) \in \Z$ for all $0 \le k \le a + s - 1$. Cauchy's integral formula on $|z - 1/n_0| = \frac{1}{2n_0}$ combined with \eqref{eq:7.4} bounds $|P_i^k(1/n_0)| \le uU^D$, while differentiating $R(x) = \sum_{v=O+1}^\infty r_v x^v$ term-by-term at $x = 1/n_0 \le 1/N$ gives
\[
\bigl|R^k(1/n_0)\bigr| \le \frac{|T(1/n_0)|^k}{k!} \sum_{v=O+1}^\infty v(v - 1)\cdots(v - k + 1)|r_v|\,n_0^{-(v-k)} \le w\,W^{-D},
\]
since $O + 1 \ge (m + 1 - \varepsilon_0)D$ and $\sum_{v > O} \binom{v}{k}(2M/n_0)^v \le 2\binom{O+1}{k}(4M/N)^{O+1}$ for $N > 8M$.

(ii) Dividing $sn_0! = a_r x_0^r(1 + O(x_0^{-\nu}))$ by $s(n_0 - \beta_i)! = a_r x_i^r(1 + O(x_i^{-\nu}))$ gives
\[
\frac{n_0!}{(n_0 - \beta_i)!} = n_0^{\beta_i} \prod_{j=1}^{\beta_i-1} \left(1 - \frac{j}{n_0}\right) = \left(\frac{x_i}{x_0}\right)^r \bigl(1 + O(x_0^{-\nu})\bigr).
\]
Raising both sides to the power $-1/r$ and multiplying by $n_0^{\beta_i/r}$ (which is an integer since $r \mid \beta_i$) yields $\omega_i(1/n_0) = \prod_{j=1}^{\beta_i-1} (1 - j/n_0)^{-1/r} = \frac{n_0^{\beta_i/r} x_i}{x_0}\bigl(1 + O(x_0^{-\nu})\bigr)$. Taking products over $1 \le i \le d$ with exponents $e_i$ ($\sum e_i \le E$) gives $\Omega_{\mathbf{e}}(1/n_0) = \frac{q_{\mathbf{e}}}{x_0^E}(1 + O(Ex_0^{-\nu})) = \frac{q_{\mathbf{e}}}{x_0^E} + O(M^E x_0^{-\nu})$, since $\omega_i(1/n_0) = 1 + O(\beta_i^2/n_0) \ll M$.
\end{proof}

\subsection{Proof of Theorem~\ref{thm:1.3}(ii)}

\begin{proof}[Proof of Theorem~\ref{thm:1.3}(ii)]
\textbf{Step 1 (The $(m + 1)$-Variable Support-Reduced Pad\'e--Wronskian Bound):} Fix $m \ge 2$, $\varepsilon_0 \in (0, m - 1)$, and $0 < \theta < \min\bigl(\frac{\varepsilon_0(m - 1 - \varepsilon_0)}{(m + 1 - \varepsilon_0)(2m - \varepsilon_0)}, \frac{1}{m+2}\bigr)$, and set $M := \lfloor N^\theta \rfloor$. Let $N \le n_1 < n_2 < \dots < n_R < 2N$ be the solutions in $[N, 2N)$ to $s \cdot n! = P(x)$. Grouping these solutions into disjoint $(mr + 1)$-tuples $(n_q, n_{q+1}, \dots, n_{q+mr})$, at most $(mr + 1)\frac{N}{M}$ solutions belong to tuples of diameter $n_{q+mr} - n_q > M$. In every tuple of diameter $\le M$, the Pigeonhole Principle modulo $r$ produces $m + 1$ solutions $n - \beta_m < \dots < n - \beta_1 < n$ in the same residue class modulo $r$, so $r \le \beta_1 < \dots < \beta_m \le M$ and $r \mid \beta_i$. Covering $[N, 2N)$ by $\le \log N + 1$ short intervals $I_{N'} = [N', N' + \frac{N}{\log N})$ and summing over the $\le M^m$ choices of shifts $(\beta_1, \dots, \beta_m)$ gives
\begin{equation}\label{eq:7.17}
R \ll_{P,m} \frac{N}{M} + M^m(\log N) \max_{\substack{r \le \beta_1 < \dots < \beta_m \le M \\ r \mid \beta_i}} \max_{N \le N' < 2N} t(\beta_1, \dots, \beta_m; N'),
\end{equation}
where $t = t(\beta_1, \dots, \beta_m; N')$ is the number of integers $n \in I_{N'}$ satisfying $s(n - \beta_i)! = P(x_i)$ for $0 \le i \le m$ ($\beta_0 = 0$).

To establish $R \ll_{P,s,m} \frac{N}{M} = N^{1-\theta}$, it suffices to prove that $t \le \frac{N}{M^{m+1}\log N}$. Suppose for contradiction that $t > \frac{N}{M^{m+1}\log N}$. By Lemma~\ref{lem:7.3}(i), there exists at least one $n_0 \in I_{N'}$ among these $t$ solutions such that $a := \ord_{x=1/n_0} \Delta_S(x) \le \frac{\deg(\Delta_S)}{t} \ll_m \frac{DM^{m+1}\log N}{N}$. Set
\[
U := \bigl(2mr^{\frac{2m}{m-\varepsilon_0}} M\bigr)^{\frac{(m-\varepsilon_0)(m+1-\varepsilon_0)}{\varepsilon_0}}, \qquad W := \frac{N^{m+1-\varepsilon_0}}{\bigl(8mr^{\frac{2m}{m-\varepsilon_0}} M\bigr)^{\frac{m(m+1-\varepsilon_0)}{\varepsilon_0}}}, \qquad C := 2N,
\]
and choose $D := 1 + \bigl\lfloor \frac{\log(bcw((N + N/\log N)!)^{1/r})}{\log(W/C)} \bigr\rfloor \asymp_{P,s,m} N$ as in \cite[Proposition~3.4]{BuiPrattZaharescu2023}, where $c = (2rN)^{(a+s)M + a + s}$, $u = (4rN)^{(a+s)M + a + s}\allowbreak \cdot 10^{m\varepsilon_0^{-1}M}$, and $w = (m+1)(D+1)(2rN)^{2(a+s)M + 2(a+s)}\allowbreak \cdot 4^{m\varepsilon_0^{-1}M}$. Since $\theta < \frac{1}{m+2}$, we have $aM \ll_m M^{m+2}\log N = o(N)$, so $c, u, w = (N!)^{o(1)}$. By Lemma~\ref{lem:7.3}(ii), there exist $0 \le k_1 < k_2 < \dots < k_s \le a + s - 1$ such that the $s \times s$ matrix $p_{i_a,j} := P_{i_a}^{k_j}(1/n_0)$ ($1 \le a, j \le s$) is invertible, with common denominator $Z := r^{a+s} n_0^{D+(a+s)M} \le cC^D$, $|p_{i_a,j}| \le uU^D$, and
\[
\left|\sum_{a=1}^s p_{i_a,j}\,\omega_{i_a}(1/n_0)\right| = \bigl|R^{k_j}(1/n_0)\bigr| \le wW^{-D} \qquad (1 \le j \le s).
\]
Let $x, x_1, \dots, x_m \in \N$ satisfy $sn_0! = P(x)$ and $s(n_0 - \beta_i)! = P(x_i)$, and set $p_0 := x > 0$ and $p_i := n_0^{\beta_i/r} x_i > 0$ ($1 \le i \le m$). Since $P(x)$ is depressed ($\nu \ge 2$), $\rho := \max_{1 \le i \le m} \bigl|\omega_i(1/n_0) - \frac{p_i}{x}\bigr| \ll_P x^{-2+o(1)} = (N!)^{\frac{1}{r}(-2+o(1))}$. Because $p_i > 0$ for every $0 \le i \le m$, the $s$-dimensional subvector $(p_{i_1}, \dots, p_{i_s})$ is non-zero, and since $\det(p_{i_a,j}) \ne 0$, there exists $j \in \{1, \dots, s\}$ such that $\sum_{a=1}^s p_{i_a,j} p_{i_a} \ne 0$. Clearing the denominator $Z \le cC^D$ and applying the triangle inequality to $p_{i,j} p_i = -xp_{i,j}(\omega_i(1/n_0) - \frac{p_i}{x}) + xp_{i,j}\omega_i(1/n_0)$ gives
\[
\frac{1}{cC^D} \le \frac{1}{Z} \le \left|\sum_{a=1}^s p_{i_a,j} p_{i_a}\right| \le sxuU^D\rho + xwW^{-D}.
\]
By our choice of $D$, $xwW^{-D} \le \frac{1}{2cC^D}$, so $\rho \gg_{P,s,m} x^{-1}(CU)^{-D} \gg (N!)^{\frac{1}{r}(-1-\chi-o(1))}$, where
\[
\chi = \frac{\log(CU)}{\log(W/C)} = \frac{1 + \theta\frac{(m-\varepsilon_0)(m+1-\varepsilon_0)}{\varepsilon_0}}{m - \varepsilon_0 - \theta\frac{m(m+1-\varepsilon_0)}{\varepsilon_0}} + O\!\left(\frac{1}{\log N}\right).
\]
This contradicts $\rho \ll (N!)^{\frac{1}{r}(-2+o(1))}$ whenever $\chi < 1 - \delta$, which rearranges to
\begin{equation}\label{eq:7.18}
\theta < \theta_m(\varepsilon_0) := \frac{\varepsilon_0(m - 1 - \varepsilon_0)}{(m + 1 - \varepsilon_0)(2m - \varepsilon_0)}.
\end{equation}
Evaluating at $m = 5$ and $\varepsilon_0 = 5 - \sqrt{5} \in (2, 3)$ gives $m - 1 - \varepsilon_0 = \sqrt{5} - 1$, $m + 1 - \varepsilon_0 = \sqrt{5} + 1$, and $2m - \varepsilon_0 = 5 + \sqrt{5} = \sqrt{5}(\sqrt{5} + 1)$, so
\[
\frac{\varepsilon_0(4 - \varepsilon_0)}{(6 - \varepsilon_0)(10 - \varepsilon_0)} = \frac{(5 - \sqrt{5})(\sqrt{5} - 1)}{(\sqrt{5} + 1)\cdot\sqrt{5}(\sqrt{5} + 1)} = \left(\frac{\sqrt{5} - 1}{\sqrt{5} + 1}\right)^2 = \frac{7 - 3\sqrt{5}}{2} \approx 0.145898 > \frac{1}{7} \approx 0.142857,
\]
so $\theta_5(5 - \sqrt{5}) = \frac{7 - 3\sqrt{5}}{2} > \frac{1}{7}$. Thus every $\theta < \frac{1}{m+2} = \frac{1}{7}$ is admissible at $m = 5$ (see Table~\ref{tab:1}), and summing over dyadic intervals $[N/2^j, N/2^{j-1})$ establishes \eqref{eq:1.10}.

\begin{table}[!htbp]
\centering
\caption{Progression of unconditional power-saving exponents $\theta$ and upper bounds $O(N^{1-\theta})$ for $\mathcal{N}_{P,s}(N)$ and $\mathcal{N}_{\mathrm{Brocard}}(N)$ under the Support-Reduced Pad\'e--Wronskian method.}\label{tab:1}
\resizebox{\linewidth}{!}{%
\begin{tabular}{cccccc}
\toprule
Shifts $m$ & Wronskian Size & Optimal $\varepsilon_0 \in (0, m - 1)$ & $\theta_m(\varepsilon_0)$ & Pigeonhole Cap $\frac{1}{m+2}$ & Unconditional Counting Bound \\
\midrule
2 \cite{BuiPrattZaharescu2023} & $3 \times 3$ & $2 - \sqrt{2}$ & $17 - 12\sqrt{2} \approx 0.02944$ & $1/4$ & $O(N^{33/34})$ \\
3 (Support-Reduced) & $s \times s$ ($s \le 4$) & $3 - \sqrt{3}$ & $7 - 4\sqrt{3} \approx 0.07180$ & $1/5$ & $O(N^{0.92821})$ \\
4 (Support-Reduced) & $s \times s$ ($s \le 5$) & $2$ & $1/9 \approx 0.11111$ & $1/6$ & $O_\varepsilon(N^{8/9+\varepsilon})$ \\
\textbf{5 (Support-Reduced)} & $\mathbf{s \times s}$ ($\mathbf{4 \le s \le 6}$) & $\mathbf{5 - \sqrt{5}}$ & $\mathbf{\frac{7 - 3\sqrt{5}}{2} \approx 0.14590 > \frac{1}{7}}$ & $\mathbf{1/7}$ & $\mathbf{O_\varepsilon(N^{6/7+\varepsilon})}$ \\
\bottomrule
\end{tabular}}
\end{table}

\textbf{Step 2 (Decoupled Monomial-Product Pad\'e--Wronskian Exponents for $\nu \ge 3$):} Suppose $r \ge 3$ and $\nu \ge 3$. Fix $d \ge 1$ and $1 \le E \le \min(r - 1, \nu - 1)$ with $m := \binom{d+E}{d} - 1 \ge 2$. By Lemma~\ref{lem:7.2}, the $m + 1 = \binom{d+E}{d}$ monomial-product functions $\Omega_{\mathbf{e}}(x) = \prod_{i=1}^d \omega_i(x)^{e_i}$ ($|\mathbf{e}| \le E$) are linearly independent over $\C(x)$. Because all $m + 1$ basis functions are generated from only $d$ shifts $\beta_1 < \dots < \beta_d \le M$, the Pigeonhole Principle over $[1, M]^d$ replaces $M^m$ in \eqref{eq:7.17} by $M^d$, so $t > \frac{N}{M^{d+1}\log N}$ yields root multiplicity $a \ll \frac{DM^{d+1}\log N}{N}$ and $(a + s)M = o(D)$ holds for all $\theta < \frac{1}{d+2}$. At $x = 1/n_0$, Lemma~\ref{lem:7.4}(ii) gives $\Omega_{\mathbf{e}}(1/n_0) = \frac{q_{\mathbf{e}}}{x_0^E} + O(x_0^{-\nu+o(1)})$ with common denominator $x_0^E$. Clearing $x_0^E$ in the non-zero support Wronskian row $\sum_{a=1}^s p_{i_a,j} q_{\mathbf{e}_a} \ne 0$ and choosing $D$ so that $x_0^E wW^{-D} \le \frac{1}{2cC^D}$ yields $\rho := \max_{|\mathbf{e}| \le E} |\Omega_{\mathbf{e}}(1/n_0) - q_{\mathbf{e}}/x_0^E| \gg x_0^{-E}(CU)^{-D} \gg x_0^{-E(1+\chi+o(1))}$, which contradicts $\rho \ll x_0^{-\nu+o(1)}$ whenever $E(1 + \chi) < \nu \iff \chi < \frac{\nu - E}{E}$, i.e.,
\begin{equation}\label{eq:7.19}
\begin{aligned}
\theta < \theta_{d,E,\nu}(\varepsilon_0) &:= \frac{\varepsilon_0\bigl((\nu - E)(m - \varepsilon_0) - E\bigr)}{(m + 1 - \varepsilon_0)\bigl(\nu m - E\varepsilon_0\bigr)}, \\
\theta_{d,E,\nu}^* &:= \min\!\left(\max_{\varepsilon_0 \in (0, m-1)} \theta_{d,E,\nu}(\varepsilon_0),\, \frac{1}{d + 2}\right).
\end{aligned}
\end{equation}
Evaluating \eqref{eq:7.19} for $\nu \ge 3$ (Table~\ref{tab:2}) completes the proof of Theorem~\ref{thm:1.3}(ii).
\end{proof}

\begin{table}[!htbp]
\centering
\caption{Unconditional counting bounds $\mathcal{N}_{P,s}(N) \ll_{P,s,\varepsilon} N^{1-\theta_{d,E,\nu}^*+\varepsilon}$ obtained via the Decoupled Monomial-Product Pad\'e--Wronskian method for polynomial-factorial equations $s \cdot n! = P(x)$ with degree gap $\nu \ge 2$.}\label{tab:2}
\resizebox{\linewidth}{!}{%
\begin{tabular}{ccccccc}
\toprule
Degree Gap $\nu$ & Shifts $d$ & Product Degree $E$ & Basis Functions $m + 1 = \binom{d+E}{d}$ & Optimal $\varepsilon_0$ & $\theta_{d,E,\nu}^*$ & Unconditional Counting Bound \\
\midrule
$\nu = 2$ (e.g.\ $n! + 1 = m^2$) & 5 & 1 & 6 ($m = 5$) & $5 - \sqrt{5} \approx 2.7639$ & $1/7 \approx 0.14286$ & $O_\varepsilon(N^{6/7+\varepsilon})$ \\
$\nu = 3$ (e.g.\ $s \cdot n! + A = x^3$) & 4 & 1 & 5 ($m = 4$) & $\frac{40 - 2\sqrt{85}}{9} \approx 2.3957$ & $1/6 \approx 0.16667$ & $O_\varepsilon(N^{5/6+\varepsilon})$ \\
$\nu = 4$ (e.g.\ $s \cdot n! + A = x^4$) & 3 & 1 & 4 ($m = 3$) & $\frac{18 - 2\sqrt{21}}{5} \approx 1.7670$ & $1/5 = 0.20000$ & $O_\varepsilon(N^{4/5+\varepsilon})$ \\
$\nu = 5$ (e.g.\ $s \cdot n! + A = x^5$) & 2 & 2 & 6 ($m = 5$) & $\frac{90 - 7\sqrt{30}}{17} \approx 3.0388$ & $0.21062$ & $O(N^{0.78939})$ \\
$\nu \ge 6$ (with $d = 2, E = 2$) & 2 & 2 & 6 ($m = 5$) & $\frac{60 - 3\sqrt{70}}{11} \approx 3.1728$ & $1/4 = 0.25000$ & $O_\varepsilon(N^{3/4+\varepsilon})$ \\
$\nu \to \infty$ (single shift $d = 1$) & 1 & $E \ge 6$ & $E + 1$ ($m = E$) & $4$ & $1/3 \approx 0.33333$ & $O_\varepsilon(N^{2/3+\varepsilon})$ \\
\bottomrule
\end{tabular}}
\end{table}

\section{Frey--Hellegouarch Elliptic Curves and Effective Class Group Bounds}\label{sec:8}

In this final section we construct the Frey--Hellegouarch elliptic curves attached to \eqref{eq:1.1} and prove the effective Tatuzawa class number and narrow $2^J$-rank bounds for $K_n = \Q(\sqrt{q_n})$.

\subsection{Frey--Hellegouarch Elliptic Curves, Base-Change Unit Factorizations, and Level-2 CM Reduction}

\begin{theorem}[{First-Power Factorial Frey Curve $\mathcal{E}_n$, Quadratic Base-Change Unit Factorization, and Level-2 CM Reduction}]\label{thm:8.1}
Let $(n, m)$ be a solution to $n! + 1 = m^2$ with $n \ge 8$, so $k \in \{1, 2\}$ by Theorem~\ref{thm:1.1}.
\begin{enumerate}
\item[(i)] The elliptic curve $\mathcal{E}_n : Y^2 = X^3 + 2mX^2 + X$ over $\Q$ is modular and minimal and semistable at all odd primes $p \le n$, with Weierstrass discriminant $\Delta(\mathcal{E}_n) = 64\,n!$ (minimal at $p = 2$ up to a $2^{12}$ Weierstrass coordinate scaling $X = 4x + 1$, $Y = 8y + 4x$ when $m \equiv -1 \pmod{32}$), conductor $N(\mathcal{E}_n) = 2^\alpha \prod_{3 \le p \le n} p$ ($\alpha \in \{0, 1, 2, 3\}$), and Frobenius traces $a_p(\mathcal{E}_n) = \bigl(\frac{-m}{p}\bigr)$ for all $3 \le p \le n$. Over $K_n = \Q(\sqrt{q_n})$, $\mathcal{E}_n$ acquires full $K_n$-rational 2-torsion and factors into units in $\mathcal{O}_{K_n}^\times$:
\begin{equation}\label{eq:8.1}
\mathcal{E}_n / K_n : Y^2 = X\bigl(X + \varepsilon_n^k\bigr)\bigl(X + \varepsilon_n^{-k}\bigr).
\end{equation}
\item[(ii)] In the $k = 2$ regime, the Level-2 Frey curve $F_n^{(2)} : Y^2 = X\bigl(X - 2(X_1(q_n) - 1)\bigr)(X + 4)$ has Weierstrass discriminant $\Delta(F_n^{(2)}) = 2^{12} q_n^2 Y_1(q_n)^4$ (minimal at all odd primes $p$, and at $p = 2$ up to a $2^{12}$ coordinate scaling), slashed conductor $N(F_n^{(2)}) = 2^{\alpha'} \frac{\prod_{p \le n} p}{\radodd(X_1(q_n))}$ ($\alpha' \in \{0, 1, 2, 3\}$), surjective mod-$\ell$ Galois representation $\bar{\rho}_{F_n^{(2)},\ell} : \Gal(\overline{\Q}/\Q) \to \GL_2(\F_\ell)$ for every prime $\ell \ge 5$, and at every odd prime $p \mid X_1(q_n)$ reduces modulo $p$ to the fixed CM curve $E_{\mathrm{CM}} : Y^2 = V^3 - 4V$ of conductor 32 with $j(F_n^{(2)}) \equiv 1728 \pmod{X_1(q_n)^2}$. Moreover, the coprime triple $(A_n, B_n, C_n) = \bigl(\frac{X_1(q_n)-1}{c_0}, \frac{2}{c_0}, \frac{X_1(q_n)+1}{c_0}\bigr)$ has $abc$ quality $q(A_n, B_n, C_n) \ge \bigl(\frac{1}{4\log 2} - o(1)\bigr)\log n \to \infty$.
\end{enumerate}
\end{theorem}

\begin{proof}
(i) The Weierstrass equation $\mathcal{E}_n : Y^2 = X^3 + 2mX^2 + X$ has invariants $b_2 = 8m$, $b_4 = 2$, $b_6 = 0$, $b_8 = -1$, so
\[
c_4(\mathcal{E}_n) = b_2^2 - 24b_4 = 16(4m^2 - 3) = 16(4n! + 1), \qquad \Delta(\mathcal{E}_n) = -b_2^2 b_8 - 8b_4^3 = 64(m^2 - 1) = 64\,n!.
\]
For every odd prime $p \le n$, we have $p \mid n!$ and $c_4(\mathcal{E}_n) \equiv 16 \not\equiv 0 \pmod p$, so $\mathcal{E}_n$ is minimal at $p$ with multiplicative reduction of exact first-power factorial valuation $v_p(\Delta(\mathcal{E}_n)) = v_p(n!)$. Reducing modulo $p$, where $m \equiv \varepsilon_p \in \{\pm 1\} \pmod p$, gives $Y^2 \equiv X(X + \varepsilon_p)^2 \pmod p$ with node at $X \equiv -\varepsilon_p \pmod p$ and tangent cone $Y^2 \equiv -\varepsilon_p(X + \varepsilon_p)^2 \pmod p$, so $a_p(\mathcal{E}_n) = \bigl(\frac{-\varepsilon_p}{p}\bigr) = \bigl(\frac{-m}{p}\bigr)$. Moreover, $\mathcal{E}_n$ is $\Q$-rationally 2-isogenous via $x = Y^2/X^2$, $y = Y(1 - X^2)/X^2$ to $\widehat{\mathcal{E}}_n : y^2 = x(x - 2(m - 1))(x - 2(m + 1))$ of discriminant $2^{12}(n!)^2$, whose quadratic twist by $-1$ is $E_n^{(2)} : Y^2 = X(X - 2(m - 1))(X + 4)$, which satisfies $a_p(E_n^{(2)}) = +1$ for all primes $p \le n$ with $p \equiv 1 \pmod 4$. Over $K_n = \Q(\sqrt{q_n})$, since $2m = \varepsilon_n^k + \varepsilon_n^{-k}$ and $\varepsilon_n^k \cdot \varepsilon_n^{-k} = 1$, the cubic factors as $X^3 + (\varepsilon_n^k + \varepsilon_n^{-k})X^2 + X = X(X + \varepsilon_n^k)(X + \varepsilon_n^{-k})$, proving \eqref{eq:8.1}.

(ii) In the $k = 2$ regime, shifting $V = X + 2$ in $F_n^{(2)} : Y^2 = X\bigl(X - 2(X_1(q_n) - 1)\bigr)(X + 4)$ gives the model
\begin{equation}\label{eq:8.2}
F_n^{(2)} : Y^2 = (V^2 - 4)\bigl(V - 2X_1(q_n)\bigr) = V^3 - 2X_1(q_n)V^2 - 4V + 8X_1(q_n),
\end{equation}
with root differences $4$, $2(X_1(q_n) - 1)$, and $2(X_1(q_n) + 1)$, discriminant $\Delta(F_n^{(2)}) = 2^{12}(X_1(q_n)^2 - 1)^2 = 2^{12} q_n^2 Y_1(q_n)^4$, and $c_4(F_n^{(2)}) = 64(X_1(q_n)^2 + 3)$. At every odd prime $p \mid X_1(q_n)$, we have $p \nmid q_n Y_1(q_n)$, so $F_n^{(2)}$ has good reduction at $p$ (slashing $\radodd(X_1(q_n))$ from the conductor $N(F_n^{(2)})$), and reducing \eqref{eq:8.2} modulo $p \mid X_1(q_n)$ yields $Y^2 \equiv V^3 - 4V \pmod p$, which is the fixed CM elliptic curve $E_{\mathrm{CM}}$ of conductor 32 and $j$-invariant 1728; indeed,
\begin{equation}\label{eq:8.3}
j\bigl(F_n^{(2)}\bigr) - 1728 = \frac{64\,X_1(q_n)^2\bigl(X_1(q_n)^2 - 9\bigr)^2}{\bigl(X_1(q_n)^2 - 1\bigr)^2} \equiv 0 \pmod{X_1(q_n)^2}.
\end{equation}
By the modularity theorem \cite{Wiles1995,BreuilConradDiamondTaylor2001}, $\mathcal{E}_n$, $E_n^{(2)}$, and $F_n^{(2)}$ are modular. For any prime $p \in I_1 = (n/2, n]$, $v_p(j(F_n^{(2)})) = -2 < 0$, so $j(F_n^{(2)}) \notin \Z$. Since $F_n^{(2)}(\Q)[2] \simeq (\Z/2\Z)^2$, Mazur's and Kenku's rational isogeny theorems \cite{Mazur1978,Kenku1982} imply that $\bar{\rho}_{F_n^{(2)},\ell}$ is irreducible for every prime $\ell \ge 5$, and Tate's $p$-adic uniformization at $p \in I_1$ shows that inertia $I_p$ acts via a non-trivial transvection of order $\ell$, forcing $\Imag(\bar{\rho}_{F_n^{(2)},\ell}) = \GL_2(\F_\ell)$ by Serre's open image theorem \cite{Serre1972}. Furthermore, for $(A_n, B_n, C_n) = \bigl(\frac{X_1(q_n)-1}{c_0}, \frac{2}{c_0}, \frac{X_1(q_n)+1}{c_0}\bigr)$, Theorem~\ref{thm:4.2}(i) gives $\log\rad(A_n B_n C_n) = \log\rad(2q_n Y_1(q_n)) \le (\log 2 + o(1))n$ while $\log C_n = (\frac{1}{4} + o(1))n\log n$, yielding $q(A_n, B_n, C_n) \ge (\frac{1}{4\log 2} - o(1))\log n \to \infty$, while writing $n! = u_n^3 v_n$ ($v_n$ cube-free) and setting $x_n = u_n v_n$, $y_n = m v_n$ gives $0 < y_n^2 - x_n^3 = v_n^2 \le x_n^{\frac{12+o(1)}{\log n}}$.
\end{proof}

\subsection{Effective Class Group Bounds}

\begin{theorem}[{Effective Tatuzawa Class Number Explosion and Narrow $2^J$-Rank Dichotomy}]\label{thm:8.2}
Let $\varepsilon \in (0, 1/8)$. With at most one exceptional integer $n_* \in \N$, every solution $n \ge 8$ to $n! + 1 = m^2$ satisfies the effective non-genus class number lower bound
\begin{equation}\label{eq:8.4}
h_{\mathrm{non\text{-}gen}}(q_n) := \frac{h(q_n)}{2^{\omega(D_n)-2}} > \frac{0.655\,\varepsilon\,k\,q_n^{1/2-\varepsilon}}{n\log n \cdot 2^{\pi(n)}} = q_n^{1/2-\varepsilon-o(1)},
\end{equation}
and in the narrow class group $\Cl^+(q_n) \simeq \mathcal{A}_2(q_n) \oplus \Cl_{\mathrm{odd}}^+(q_n)$, either $|\Cl_{\mathrm{odd}}^+(q_n)| \ge q_n^{1/4-\varepsilon}$, or $|\mathcal{A}_2(q_n)| \ge 2^{(\frac{\log 2}{4} - \varepsilon)n}$ and $\exp(\mathcal{A}_2(q_n)) \ge n^{\frac{1}{4\log 2}-\varepsilon}$ (with $\sum_{j \ge 1} \rank_{2^j}(\mathcal{A}_2(q_n)) = \log_2|\mathcal{A}_2(q_n)| \ge (\frac{\log 2}{4} - \varepsilon)n$).
\end{theorem}

\begin{proof}
By Dirichlet's analytic class number formula \cite[Chapter~6]{Davenport2000}, $h(q_n)R(q_n) = \frac{1}{2}\sqrt{D_n} L(1, \chi_{D_n})$. By Tatuzawa's effective refinement of Siegel's theorem \cite{Siegel1935,Tatuzawa1951}, for any $\varepsilon \in (0, 1/8)$, $L(1, \chi_D) > 0.655\,\varepsilon\,D^{-\varepsilon}$ holds for all fundamental discriminants $D > 0$ with at most one exception $D_*$. By Theorem~\ref{thm:1.1}, distinct solutions $n_1 < n_2$ have distinct square-free cores $q_{n_1} \ne q_{n_2}$, so $D_n = D_*$ can occur for at most one solution $n_* \in \N$. Substituting $\sqrt{D_n} \ge \sqrt{q_n}$, $L(1, \chi_{D_n}) > 0.655\,\varepsilon\,(4q_n)^{-\varepsilon}$, and $R(q_n) \le \log \varepsilon_n \le \frac{1}{2k}n\log n$ yields \eqref{eq:8.4}. In $\Cl^+(q_n) \simeq \mathcal{A}_2(q_n) \oplus \Cl_{\mathrm{odd}}^+(q_n)$, Gauss's genus theory gives $\rank_2(\mathcal{A}_2(q_n)) = \omega(D_n) - 1 = (\log 2 + o(1))\frac{n}{\log n}$; thus if $|\Cl_{\mathrm{odd}}^+(q_n)| < q_n^{1/4-\varepsilon}$, then $\log_2|\mathcal{A}_2(q_n)| \ge (\frac{1}{4} - \frac{\varepsilon}{2})\log_2 q_n = (\frac{\log 2}{4} - \frac{\varepsilon}{2} - o(1))n$, which forces $|\mathcal{A}_2(q_n)| \ge 2^{(\frac{\log 2}{4} - \varepsilon)n}$ and $\log_2\exp(\mathcal{A}_2(q_n)) \ge \frac{\log_2|\mathcal{A}_2(q_n)|}{\rank_2(\mathcal{A}_2(q_n))} \ge (\frac{1}{4\log 2} - \varepsilon)\log n$.
\end{proof}

\section*{Acknowledgments}
This manuscript was developed over four months of intensive human--AI mathematical research collaboration between the author and Google Gemini. The AI system assisted in exploring analytical strategies, checking symbolic identities, and drafting and formatting the \LaTeX{} manuscript. All mathematical definitions, lemmas, propositions, theorems, proofs, and bibliographical citations have been independently verified by the author, who takes full responsibility for the correctness and content of this work.

\begin{thebibliography}{99}

\bibitem{Berend1997}
D.~Berend, \emph{On the parity of exponents in the factorization of $n!$}, J. Number Theory \textbf{64} (1997), no.~1, 13--19.

\bibitem{BerendHarmse2006}
D.~Berend and J.~E. Harmse, \emph{On polynomial-factorial Diophantine equations}, Trans. Amer. Math. Soc. \textbf{358} (2006), no.~4, 1741--1779.

\bibitem{BerendOsgood1992}
D.~Berend and C.~F. Osgood, \emph{On the equation $P(x) = n!$ and a question of Erd\H{o}s}, J. Number Theory \textbf{42} (1992), no.~2, 189--193.

\bibitem{BerndtGalway2000}
B.~C. Berndt and W.~F. Galway, \emph{On the Brocard--Ramanujan Diophantine equation $n! + 1 = m^2$}, Ramanujan J. \textbf{4} (2000), no.~1, 41--42.

\bibitem{BreuilConradDiamondTaylor2001}
C.~Breuil, B.~Conrad, F.~Diamond, and R.~Taylor, \emph{On the modularity of elliptic curves over $\Q$: wild $3$-adic exercises}, J. Amer. Math. Soc. \textbf{14} (2001), no.~4, 843--939.

\bibitem{Breusch1932}
R.~Breusch, \emph{Zur Verallgemeinerung des Bertrandschen Postulates, dass zwischen $x$ und $2x$ stets Primzahlen liegen}, Math. Z. \textbf{34} (1932), no.~1, 505--526.

\bibitem{Brocard1876}
H.~Brocard, \emph{Question 166}, Nouv. Corresp. Math. \textbf{2} (1876), 287.

\bibitem{Brocard1885}
H.~Brocard, \emph{Question 1532}, Nouv. Ann. Math. (3) \textbf{4} (1885), 391.

\bibitem{Broughan2002}
K.~A. Broughan, \emph{Asymptotic order of the square-free part of $n!$}, Integers \textbf{2} (2002), A10, 1--6.

\bibitem{BuiPrattZaharescu2023}
H.~M. Bui, K.~Pratt, and A.~Zaharescu, \emph{Power savings for counting solutions to polynomial-factorial equations}, Adv. Math. \textbf{422} (2023), 109021.

\bibitem{Coates1977}
J.~Coates, \emph{$p$-adic $L$-functions and Iwasawa's theory}, in \emph{Algebraic Number Fields: $L$-Functions and Galois Properties} (Proc. Sympos., Univ. Durham, 1975), Academic Press, London, 1977, pp.~269--353.

\bibitem{Dabrowski1996}
A.~D\k{a}browski, \emph{On the Diophantine equation $x! + A = y^2$}, Nieuw Arch. Wiskd. (4) \textbf{14} (1996), no.~3, 321--324.

\bibitem{DabrowskiUlas2013}
A.~D\k{a}browski and M.~Ulas, \emph{Variations on the Brocard--Ramanujan equation}, J. Number Theory \textbf{133} (2013), 1168--1185.

\bibitem{Davenport2000}
H.~Davenport, \emph{Multiplicative Number Theory}, 3rd ed., Graduate Texts in Mathematics, vol.~74, Springer-Verlag, New York, 2000.

\bibitem{EpsteinGlickman2020}
A.~Epstein and K.~Glickman, \emph{C++ Brocard SDK: Verification of Brocard's problem up to $10^{15}$}, GitHub Repository (2020), cited in \cite{BuiPrattZaharescu2023}.

\bibitem{ErdosOblath1937}
P.~Erd\H{o}s and R.~Obl\'ath, \emph{\"Uber diophantische Gleichungen der Form $n! = x^p \pm y^p$ und $n! \pm m! = x^p$}, Acta Litt. Sci. Szeged \textbf{8} (1937), 241--255.

\bibitem{ErdosSelfridge1975}
P.~Erd\H{o}s and J.~L. Selfridge, \emph{The product of consecutive integers is never a power}, Illinois J. Math. \textbf{19} (1975), 292--301.

\bibitem{FerreroWashington1979}
B.~Ferrero and L.~C. Washington, \emph{The Iwasawa invariant $\mu_p$ vanishes for abelian number fields}, Ann. of Math. (2) \textbf{109} (1979), no.~2, 377--395.

\bibitem{FukudaKomatsu1986}
T.~Fukuda and K.~Komatsu, \emph{On $\Z_p$-extensions of real quadratic fields}, J. Math. Soc. Japan \textbf{38} (1986), no.~1, 95--102.

\bibitem{Gauss1801}
C.~F. Gauss, \emph{Disquisitiones Arithmeticae}, Gerhard Fleischer, Leipzig, 1801.

\bibitem{Gras2016}
G.~Gras, \emph{Les $\theta$-r\'egulateurs locaux d'un nombre alg\'ebrique: conjectures $p$-adiques}, Canad. J. Math. \textbf{68} (2016), no.~3, 571--624.

\bibitem{Greenberg1976}
R.~Greenberg, \emph{On the Iwasawa invariants of totally real number fields}, Amer. J. Math. \textbf{98} (1976), no.~1, 263--284.

\bibitem{Greenberg1978}
R.~Greenberg, \emph{On the structure of certain Galois groups}, Invent. Math. \textbf{47} (1978), no.~1, 85--99.

\bibitem{IwaniecKowalski2004}
H.~Iwaniec and E.~Kowalski, \emph{Analytic Number Theory}, AMS Colloquium Publications, vol.~53, American Mathematical Society, Providence, RI, 2004.

\bibitem{Jaulent2024}
J.-F. Jaulent, \emph{Sur la trivialit\'e de certains modules d'Iwasawa}, Funct. Approx. Comment. Math. \textbf{70} (2024), no.~1, 29--39.

\bibitem{Kenku1982}
M.~A. Kenku, \emph{On the number of $\Q$-isomorphism classes of elliptic curves in each $\Q$-isogeny class}, J. Number Theory \textbf{15} (1982), no.~2, 199--202.

\bibitem{KihelLuca2008}
O.~Kihel and F.~Luca, \emph{Variants of the Brocard--Ramanujan equation}, J. Th\'eor. Nombres Bordeaux \textbf{20} (2008), no.~2, 353--363.

\bibitem{Krasniqi2026}
V.~Krasniqi, \emph{The Factorial Separation Conjecture}, Preprint \texttt{arXiv:2610.25661v1} (2026).

\bibitem{Legendre1798}
A.~M. Legendre, \emph{Essai sur la th\'eorie des nombres}, Duprat, Paris, 1798.

\bibitem{Luca2002}
F.~Luca, \emph{The Diophantine equation $P(x) = n!$ and a result of M.~Overholt}, Glas. Mat. Ser. III \textbf{37(57)} (2002), no.~2, 269--273.

\bibitem{Matson2017}
R.~D. Matson, \emph{Brocard's problem 4th solution search utilizing quadratic residues}, Unsolved Problems in Number Theory, Logic and Cryptography (2017), 1--10.

\bibitem{Mazur1978}
B.~Mazur, \emph{Rational isogenies of prime degree}, Invent. Math. \textbf{44} (1978), no.~2, 129--162.

\bibitem{MazurWiles1984}
B.~Mazur and A.~Wiles, \emph{Class fields of abelian extensions of $\Q$}, Invent. Math. \textbf{76} (1984), no.~2, 179--330.

\bibitem{Nagura1952}
J.~Nagura, \emph{On the interval containing at least one prime number}, Proc. Japan Acad. \textbf{28} (1952), no.~4, 177--181.

\bibitem{Overholt1993}
M.~Overholt, \emph{The Diophantine equation $n! + 1 = m^2$}, Bull. London Math. Soc. \textbf{25} (1993), no.~2, 104.

\bibitem{Ozaki1997}
M.~Ozaki, \emph{The class group of $\Z_p$-extensions over totally real number fields}, Tohoku Math. J. (2) \textbf{49} (1997), no.~3, 431--435.

\bibitem{OzakiTaya1998}
M.~Ozaki and H.~Taya, \emph{A note on the Iwasawa $\lambda$-invariants of real abelian number fields}, Interdiscip. Inform. Sci. \textbf{4} (1998), no.~2, 109--116.

\bibitem{Perron1954}
O.~Perron, \emph{Die Lehre von den Kettenbr\"uchen}, Band~I, 3rd ed., B.~G. Teubner, Stuttgart, 1954.

\bibitem{PollackShapiro1973}
R.~M. Pollack and H.~N. Shapiro, \emph{The next to last case of a factorial Diophantine equation}, Comm. Pure Appl. Math. \textbf{26} (1973), 313--325.

\bibitem{QiStokes2024}
P.~Qi and M.~Stokes, \emph{On the non $p$-rationality and Iwasawa invariants of certain real quadratic fields}, Preprint \texttt{arXiv:2408.03836v1} (2024).

\bibitem{Ramanujan1913}
S.~Ramanujan, \emph{Question 469}, J. Indian Math. Soc. \textbf{5} (1913), 59.

\bibitem{RosserSchoenfeld1962}
J.~B. Rosser and L.~Schoenfeld, \emph{Approximate formulas for some functions of prime numbers}, Illinois J. Math. \textbf{6} (1962), no.~1, 64--94.

\bibitem{Sander2001}
J.~W. Sander, \emph{On the parity of exponents in the prime factorization of factorials}, J. Number Theory \textbf{90} (2001), no.~2, 316--328.

\bibitem{Schoenfeld1976}
L.~Schoenfeld, \emph{Sharper bounds for the Chebyshev functions $\theta(x)$ and $\psi(x)$. II}, Math. Comp. \textbf{30} (1976), no.~134, 337--360.

\bibitem{Serre1972}
J.-P. Serre, \emph{Propri\'et\'es galoisiennes des points d'ordre fini des courbes elliptiques}, Invent. Math. \textbf{15} (1972), no.~4, 259--331.

\bibitem{Siegel1935}
C.~L. Siegel, \emph{\"Uber die Classenzahl quadratischer Zahlk\"orper}, Acta Arith. \textbf{1} (1935), no.~1, 83--86.

\bibitem{Siegel1949}
C.~L. Siegel, \emph{Transcendental Numbers}, Annals of Mathematics Studies, no.~16, Princeton University Press, Princeton, NJ, 1949.

\bibitem{Tatuzawa1951}
T.~Tatuzawa, \emph{On a theorem of Siegel}, Jpn. J. Math. \textbf{21} (1951), 163--178.

\bibitem{Taya1996}
H.~Taya, \emph{On cyclotomic $\Z_p$-extensions of real quadratic fields}, Acta Arith. \textbf{74} (1996), no.~2, 107--119.

\bibitem{Vaaler1985}
J.~D. Vaaler, \emph{Some extremal functions in Fourier analysis}, Bull. Amer. Math. Soc. (N.S.) \textbf{12} (1985), no.~2, 183--216.

\bibitem{Washington1997}
L.~C. Washington, \emph{Introduction to Cyclotomic Fields}, 2nd ed., Graduate Texts in Mathematics, vol.~83, Springer-Verlag, New York, 1997.

\bibitem{Wiles1995}
A.~Wiles, \emph{Modular elliptic curves and Fermat's Last Theorem}, Ann. of Math. (2) \textbf{141} (1995), no.~3, 443--551.

\bibitem{Yanagisawa2026}
T.~Yanagisawa, \emph{On finiteness properties of the unramified Iwasawa module of a $\Z_p$-extension with restricted ramification}, Preprint \texttt{arXiv:2606.22324v2} (2026).

\end{thebibliography}

\end{document}
